% Ressemblances et différences au sein de l’espèce humaine — SVTEEHB 3ème — Cours signé OnBuch+
% Programme officiel MINESEC. Compiler avec : tectonic N01-ressemblances-differences-sein-espece-humaine.tex
\newcommand{\DOCMATIERE}{SVTEEHB}
\newcommand{\DOCNIVEAU}{3ème}
\newcommand{\DOCMODULE}{SVTEEHB – classe de 3ème}
\newcommand{\DOCLECON}{N01}
\newcommand{\DOCTITRE}{Ressemblances et différences au sein de l’espèce humaine}
% OnBuch+ — préambule « cours signés » ENRICHI (compilé avec tectonic / XeLaTeX)
% Système visuel « Pop » d'OnBuch+ : fond crème (jamais blanc), bordures encre
% épaisses, ombres « dures » décalées, titres Archivo Black en capitales.
% Version MATHÉMATIQUES : graphes (pgfplots/tikz), boîtes pédagogiques
% (définition, propriété, méthode, attention, le savais-tu, exercice résolu,
% exercice, TP, bilan), page de garde et signature OnBuch+ ancrée en bas.
%
% Macros attendues dans main.tex AVANT \input{preamble} :
%   \DOCMATIERE  \DOCNIVEAU  \DOCMODULE  \DOCLECON (numéro)  \DOCTITRE
\documentclass[11pt]{article}

\usepackage{fontspec}
\usepackage[french]{babel}
\usepackage[a4paper,margin=2cm,top=3.4cm,bottom=2.8cm,headheight=1.4cm,footskip=1.3cm]{geometry}
\usepackage{amsmath,amssymb,amsfonts}
\usepackage{mathtools} % \xmapsto, \coloneqq... souvent utilisés par les modèles
\usepackage{cancel} % simplifications barrées (bilans de forces)
% \abs{z} (module, valeur absolue) et \norm{u} (norme), utilisés par les modèles
\providecommand{\abs}{}\let\abs\relax\DeclarePairedDelimiter{\abs}{\lvert}{\rvert}
\providecommand{\norm}{}\let\norm\relax\DeclarePairedDelimiter{\norm}{\lVert}{\rVert}
\usepackage[table]{xcolor} % [table] : fournit aussi \rowcolors (alternance de lignes)
\usepackage{pagecolor}
\usepackage{tikz}
\usetikzlibrary{arrows.meta,positioning,calc,shapes.geometric,shapes.misc,shapes.arrows,decorations.pathmorphing,decorations.pathreplacing,decorations.markings,patterns,shadows,mindmap,trees,fit,backgrounds,matrix,intersections,angles,quotes}
\usepackage{pgfplots}
\usepackage[version=4]{mhchem} % équations nucléaires \ce{^{235}_{92}U}
\usepackage[european,siunitx]{circuitikz} % schémas électriques
\pgfplotsset{compat=1.18}
\usepgfplotslibrary{fillbetween,groupplots} % aires entre courbes (calcul intégral)
\usepackage{enumitem}
\usepackage{fancyhdr}
% [most] charge toutes les bibliothèques tcolorbox (listings, minted, raster,
% documentation...) et gonfle inutilement la mémoire TeX sur un document long
% (~300 boîtes) : on ne charge que ce qui est réellement utilisé.
\usepackage[skins,breakable]{tcolorbox}
\usepackage{graphicx}
\usepackage{colortbl}
\usepackage{booktabs}
\usepackage{tabularx}
\usepackage{multirow}
\usepackage{diagbox} % case d'en-tête diagonale des tableaux à double entrée
% Colonnes extensibles centrée / gauche / droite (C, L, R), souvent utilisées par les modèles
\newcolumntype{C}{>{\centering\arraybackslash}X}
\newcolumntype{L}{>{\raggedright\arraybackslash}X}
\newcolumntype{R}{>{\raggedleft\arraybackslash}X}
\usepackage{array}
\usepackage{multicol}
\usepackage{siunitx}
% parse-numbers=false : accepte aussi \SI{2,5 \times 10^{-3}}{...}
\sisetup{locale=FR,output-decimal-marker={,},group-minimum-digits=4,parse-numbers=false}
\DeclareSIUnit\ohm{\ensuremath{\symup{\Omega}}} % U+2126 absent de la police
\DeclareSIUnit\division{div} % réglages d'oscilloscope (ms/div, V/div)
\DeclareSIUnit\tr{tr} % tours (vitesse de rotation, ex. \SI{1500}{\tr\per\minute})
\DeclareSIUnit\molar{M} % concentration molaire (ex. \SI{3}{\molar}, \SI{1}{\micro\molar})
\DeclareSIUnit\an{an} % année (le modèle confond parfois avec \year, un compteur TeX) — ex. \SI{5}{\kilo\meter\per\an}
\DeclareSIUnit\FCFA{FCFA} % franc CFA (ex. \SI{500}{\FCFA\per\kilo\gram})
\DeclareSIUnit\degreeBrix{\textdegree Brix} % teneur en sucre d'un sirop/jus (réfractométrie)
\DeclareSIUnit\month{mois} % le modèle utilise parfois \month, un compteur TeX, comme unité de durée
\DeclareSIUnit\microliter{\micro\liter} % le modèle écrit parfois \microliter en un seul token
\DeclareSIUnit\million{millions} % dénombrement (ex. \SI{15}{\million\per\mL} spermatozoïdes)
\usepackage{qrcode}
% \degree : commande fréquemment utilisée par les modèles pour un angle ou une
% température en dehors de \SI{}{} (non fournie par les paquets chargés).
\providecommand{\degree}{^{\circ}}
% \qed (fin de démonstration), utilisé par les modèles sans amsthm
\providecommand{\qed}{\hfill$\square$}
% \barre : double barre des valeurs interdites dans un tableau de variations (tkz-tab)
\providecommand{\barre}{\ensuremath{\|}}
% environnement proof (amsthm non chargé) utilisé par les modèles
\ifdefined\proof\else\newenvironment{proof}[1][Démonstration]{\par\noindent\textit{#1.}\ }{\qed\par}\fi
\usepackage{lastpage}
\usepackage{hyperref}
\hypersetup{hidelinks}

% ============================================================
% Palette « Pop » OnBuch+ — mêmes hex que la classe P (plus_ui.dart)
% ============================================================
\definecolor{poporange}{HTML}{F59321}
\definecolor{popdark}{HTML}{E07A0C}
\definecolor{popcream}{HTML}{FFF6E8}
\definecolor{popink}{HTML}{1C1714}
\definecolor{poppurple}{HTML}{7A5AE0}
\definecolor{popgreen}{HTML}{2E9E6A}
\definecolor{poppink}{HTML}{E0457B}
\definecolor{popblue}{HTML}{2F7DE1}
\definecolor{popgold}{HTML}{C98A12}
\definecolor{popmuted}{HTML}{8A7F76}
\definecolor{popline}{HTML}{EADFD2}
\definecolor{popgray}{HTML}{8A7F76}
\colorlet{popgrey}{popgray}
% Alias tolérants (le modèle nomme parfois les couleurs différemment)
\colorlet{popwhite}{white}
\colorlet{popblack}{popink}
\colorlet{popred}{poppink}
\colorlet{popyellow}{popgold}
% Teintes claires (fonds de tableaux/encadrés) — le modèle nomme parfois
% les couleurs avec un suffixe L (ex. popblueL) sans les avoir définies.
\colorlet{poporangeL}{poporange!20}
\colorlet{popdarkL}{popdark!20}
\colorlet{poppurpleL}{poppurple!20}
\colorlet{popgreenL}{popgreen!20}
\colorlet{poppinkL}{poppink!20}
\colorlet{popblueL}{popblue!20}
\colorlet{popgoldL}{popgold!20}
\colorlet{popredL}{popred!20}
\colorlet{popgrayL}{popgray!20}
% Teintes claires pour fonds de tableaux / zones
\colorlet{popblueL}{popblue!10!white}
\colorlet{poporangeL}{poporange!14!white}
\colorlet{popgreenL}{popgreen!12!white}
\colorlet{poppurpleL}{poppurple!12!white}
\colorlet{poppinkL}{poppink!12!white}
\colorlet{popgoldL}{popgold!14!white}

\pagecolor{popcream}

% ============================================================
% Typographie
% ============================================================
\defaultfontfeatures{Ligatures=TeX}
\setmainfont{PlusJakartaSans-Medium.ttf}[Path=../../../fonts/,BoldFont=PlusJakartaSans-Bold.ttf]
% Maths en sans-serif (Fira Math) avec chiffres et lettres droites de Plus
% Jakarta, pour que les formules se fondent dans le texte.
\usepackage[math-style=ISO,bold-style=ISO]{unicode-math}
\setmathfont{FiraMath-Regular.otf}[Path=../../../fonts/]
\setmathfont{PlusJakartaSans-Medium.ttf}[Path=../../../fonts/,range={up/num,up/{Latin,latin},\mathup}]
\usepackage{icomma} % virgule décimale sans espace en mode math
% Caractères mathématiques Unicode absents des polices de texte (Plus Jakarta,
% Archivo Black) : rendus par leur équivalent mathématique, en texte comme en
% formule — sinon ils s'affichent en carré vide (ex. « Arithmétique dans ℤ »).
\usepackage{newunicodechar}
\newunicodechar{ℝ}{\ensuremath{\mathbb{R}}}
\newunicodechar{ℤ}{\ensuremath{\mathbb{Z}}}
\newunicodechar{ℂ}{\ensuremath{\mathbb{C}}}
\newunicodechar{ℕ}{\ensuremath{\mathbb{N}}}
\newunicodechar{ℚ}{\ensuremath{\mathbb{Q}}}
\newunicodechar{𝔻}{\ensuremath{\mathbb{D}}}
\newunicodechar{α}{\ensuremath{\alpha}}
\newunicodechar{β}{\ensuremath{\beta}}
\newunicodechar{γ}{\ensuremath{\gamma}}
\newunicodechar{δ}{\ensuremath{\delta}}
\newunicodechar{ε}{\ensuremath{\varepsilon}}
\newunicodechar{θ}{\ensuremath{\theta}}
\newunicodechar{λ}{\ensuremath{\lambda}}
\newunicodechar{μ}{\ensuremath{\mu}}
\newunicodechar{σ}{\ensuremath{\sigma}}
\newunicodechar{τ}{\ensuremath{\tau}}
\newunicodechar{φ}{\ensuremath{\varphi}}
\newunicodechar{ψ}{\ensuremath{\psi}}
\newunicodechar{ω}{\ensuremath{\omega}}
\newunicodechar{Σ}{\ensuremath{\Sigma}}
% majuscules produites par \MakeUppercase dans les titres (α-aminés -> Α)
\newunicodechar{Α}{\ensuremath{\alpha}}
\newunicodechar{Β}{\ensuremath{\beta}}
\newunicodechar{Γ}{\ensuremath{\Gamma}}
\newunicodechar{Δ}{\ensuremath{\Delta}}
\newunicodechar{Ε}{\ensuremath{\varepsilon}}
\newunicodechar{Θ}{\ensuremath{\Theta}}
\newunicodechar{Λ}{\ensuremath{\Lambda}}
\newunicodechar{Μ}{\ensuremath{\mu}}
\newunicodechar{Τ}{\ensuremath{\tau}}
\newunicodechar{Φ}{\ensuremath{\Phi}}
\newunicodechar{Ψ}{\ensuremath{\Psi}}
\newunicodechar{Ω}{\ensuremath{\Omega}}
\newunicodechar{↦}{\ensuremath{\mapsto}}
\newunicodechar{∈}{\ensuremath{\in}}
\newunicodechar{∉}{\ensuremath{\notin}}
\newunicodechar{∧}{\ensuremath{\wedge}}
\newunicodechar{⇔}{\ensuremath{\Leftrightarrow}}
\newunicodechar{⟺}{\ensuremath{\Longleftrightarrow}}
\newunicodechar{⇒}{\ensuremath{\Rightarrow}}
\newunicodechar{⟹}{\ensuremath{\Longrightarrow}}
\newunicodechar{∘}{\ensuremath{\circ}}
\newunicodechar{∪}{\ensuremath{\cup}}
\newunicodechar{∩}{\ensuremath{\cap}}
\newunicodechar{≡}{\ensuremath{\equiv}}
\newunicodechar{⊂}{\ensuremath{\subset}}
\newunicodechar{⊥}{\ensuremath{\perp}}
\newunicodechar{∃}{\ensuremath{\exists}}
\newunicodechar{∀}{\ensuremath{\forall}}
\newunicodechar{∛}{\ensuremath{\sqrt[3]{\,}}}
\newunicodechar{∜}{\ensuremath{\sqrt[4]{\,}}}
\newunicodechar{⃗}{\ensuremath{{}^{\to}}}
\newunicodechar{ⁿ}{\ifmmode^{n}\else\textsuperscript{\ensuremath{n}}\fi}
\newunicodechar{ˣ}{\ifmmode^{x}\else\textsuperscript{\ensuremath{x}}\fi}
\newunicodechar{ᵇ}{\ifmmode^{b}\else\textsuperscript{\ensuremath{b}}\fi}
\newunicodechar{ʳ}{\ifmmode^{r}\else\textsuperscript{\ensuremath{r}}\fi}
\newunicodechar{ᵘ}{\ifmmode^{u}\else\textsuperscript{\ensuremath{u}}\fi}
\newunicodechar{ʸ}{\ifmmode^{y}\else\textsuperscript{\ensuremath{y}}\fi}
\newunicodechar{ᵃ}{\ifmmode^{a}\else\textsuperscript{\ensuremath{a}}\fi}
\newunicodechar{ᵉ}{\ifmmode^{e}\else\textsuperscript{\ensuremath{e}}\fi}
\newunicodechar{ᵏ}{\ifmmode^{k}\else\textsuperscript{\ensuremath{k}}\fi}
\newunicodechar{ᵖ}{\ifmmode^{p}\else\textsuperscript{\ensuremath{p}}\fi}
\newunicodechar{ᴺ}{\ifmmode^{N}\else\textsuperscript{\ensuremath{N}}\fi}
\newunicodechar{ᶜ}{\ifmmode^{c}\else\textsuperscript{\ensuremath{c}}\fi}
\newunicodechar{ʰ}{\ifmmode^{h}\else\textsuperscript{\ensuremath{h}}\fi}
\newunicodechar{ᵠ}{\ifmmode^{q}\else\textsuperscript{\ensuremath{q}}\fi}
\newunicodechar{ₙ}{\ifmmode_{n}\else\textsubscript{\ensuremath{n}}\fi}
\newunicodechar{ₐ}{\ifmmode_{a}\else\textsubscript{\ensuremath{a}}\fi}
\newunicodechar{ᵢ}{\ifmmode_{i}\else\textsubscript{\ensuremath{i}}\fi}
\newunicodechar{ᵣ}{\ifmmode_{r}\else\textsubscript{\ensuremath{r}}\fi}
\newunicodechar{ₖ}{\ifmmode_{k}\else\textsubscript{\ensuremath{k}}\fi}
\newunicodechar{ᵦ}{\ifmmode_{\beta}\else\textsubscript{\ensuremath{\beta}}\fi}
\newunicodechar{ₓ}{\ifmmode_{x}\else\textsubscript{\ensuremath{x}}\fi}
\newfontfamily\popdisplay{ArchivoBlack-Regular.ttf}[Path=../../../fonts/]
\newfontfamily\poplogo{SpaceGrotesk-Bold.ttf}[Path=../../../fonts/]
\newfontfamily\popsemi{PlusJakartaSans-SemiBold.ttf}[Path=../../../fonts/]
\newfontfamily\popxbold{PlusJakartaSans-ExtraBold.ttf}[Path=../../../fonts/]
\newfontfamily\popxboldls{PlusJakartaSans-ExtraBold.ttf}[Path=../../../fonts/,LetterSpace=3.0]

\setlist[enumerate]{leftmargin=1.7em,itemsep=5pt,topsep=4pt}
\setlist[itemize]{leftmargin=1.7em,itemsep=4pt,topsep=4pt,label={\color{poporange}\textbullet}}
\setlength{\parindent}{0pt}
\setlength{\parskip}{5pt}
\renewcommand{\arraystretch}{1.25}

% pgfplots : style Pop par défaut
\pgfplotsset{
  every axis/.append style={axis line style={popink,line width=1pt},tick style={popink},
    grid style={popline,line width=0.5pt},label style={font=\small},tick label style={font=\footnotesize}},
  popcurve/.style={poporange,line width=1.8pt,no markers,smooth},
}
% Même style disponible dans un \draw TikZ simple (hors pgfplots)
\tikzset{popcurve/.style={poporange,line width=1.8pt}}
\tikzset{popgrid/.style={popline,very thin}}
\colorlet{popcurve}{poporange} % le nom du style est parfois utilisé comme couleur

% ============================================================
% Pill / squiggle
% ============================================================
\newcommand{\poppill}[3]{%
  \tikz[baseline=-0.55ex]\node[draw=popink,line width=1.2pt,rounded corners=2.6pt,%
    fill=#2,inner xsep=6.5pt,inner ysep=3pt]{%
    \popxboldls\fontsize{7}{7}\selectfont\color{#3}\MakeUppercase{#1}};%
}
\newcommand{\popsquiggle}[1]{%
  \tikz[baseline=-0.4ex]{
    \draw[#1,line width=2.6pt,line cap=round]
      plot [smooth] coordinates {(0,0.09) (0.34,0.01) (0.68,0.21) (1.02,0.01) (1.36,0.10)};
  }%
}
% Mot-clé surligné (marqueur orange sous le texte)
\newcommand{\cle}[1]{{\popxbold\color{popdark}#1}}

% ============================================================
% En-tête / pied
% ============================================================
\pagestyle{fancy}
\fancyhf{}
\renewcommand{\headrulewidth}{0pt}
\renewcommand{\footrulewidth}{0pt}
\lhead{\raisebox{-0.15ex}{\poplogo\fontsize{13}{13}\selectfont\color{popink}OnBuch\color{poporange}+}}
\rhead{\poppill{\DOCMATIERE\ \textbullet\ \DOCNIVEAU\ \textbullet\ Leçon \DOCLECON}{white}{popink}}
\lfoot{\popxbold\fontsize{7.6}{7.6}\selectfont\color{popmuted}\MakeUppercase{Cours signé OnBuch+}\ \textbullet\ {\popsemi\DOCTITRE}}
\rfoot{\tikz[baseline=-0.6ex]\node[rounded corners=8.5pt,fill=popink,minimum height=17pt,minimum width=17pt,inner xsep=5pt,inner sep=0]{\popxbold\fontsize{8}{8}\selectfont\color{popcream}\thepage\,/\,\pageref*{LastPage}};}

% ============================================================
% Titres
% ============================================================
\newcommand{\chaptitre}[2]{%
  \begin{tcolorbox}[enhanced,colback=poporange,colframe=popink,boxrule=2.6pt,arc=11pt,
      drop shadow=popink,left=15pt,right=15pt,top=12pt,bottom=11pt,before skip=3pt,after skip=18pt]
    \poppill{#2}{popink}{popcream}\par\vspace{9pt}
    {\raggedright\hyphenpenalty=10000\exhyphenpenalty=10000\popdisplay\fontsize{22}{24}\selectfont\color{popcream}\MakeUppercase{#1}\par}
  \end{tcolorbox}
}
% Section numérotée : pastille orange avec le numéro + titre Archivo
\newcounter{popsec}
\newcommand{\coursec}[1]{%
  \stepcounter{popsec}\setcounter{popsub}{0}%
  \par\vspace{16pt}\noindent
  \begin{minipage}{\linewidth}
  \tikz[baseline=-0.7ex]\node[draw=popink,line width=1.6pt,fill=poporange,rounded corners=4pt,minimum size=22pt,inner sep=0]{\popdisplay\fontsize{12}{12}\selectfont\color{popcream}\thepopsec};%
  \hspace{8pt}{\popdisplay\fontsize{15}{17}\selectfont\color{popink}\MakeUppercase{#1}}\par
  \vspace{4pt}\noindent{\color{poporange}\rule{36pt}{3.6pt}}
  \end{minipage}\par\nopagebreak\vspace{8pt}
}
\newcounter{popsub}
\newcommand{\courssub}[1]{%
  \stepcounter{popsub}%
  \par\vspace{9pt}\noindent{\popxbold\fontsize{11.5}{13}\selectfont\color{popink}{\color{poporange}\thepopsec.\thepopsub}\hspace{6pt}#1}\par\nopagebreak\vspace{4pt}
}

% ============================================================
% Boîtes pédagogiques — même recette : fond blanc, bordure épaisse colorée,
% ombre dure encre, badge plein. Titre optionnel [#1].
% ============================================================
\tcbset{popbase/.style={breakable,enhanced,colback=white,boxrule=2.3pt,arc=9pt,
  shadow={3pt}{-3pt}{0pt}{popink},left=14pt,right=14pt,top=11pt,bottom=11pt,before skip=11pt,after skip=15pt,
  fontupper=\popsemi\fontsize{10.6}{14.5}\selectfont\color{popink},
  fontlower=\popsemi\fontsize{10.6}{14.5}\selectfont\color{popink}}}
% Coche dessinée (le glyphe ✓ n'existe pas dans les polices du document)
\newcommand{\popcheck}[1]{\tikz[baseline=-0.5ex]\draw[#1,line width=1.6pt,line cap=round,line join=round](-0.13,0.0)--(-0.04,-0.09)--(0.13,0.1);}
\providecommand{\coche}{\popcheck{popgreen}}
\providecommand{\resultat}[1]{\par\smallskip\noindent\fcolorbox{popgreen}{popgreenL}{\parbox{\dimexpr\linewidth-2\fboxsep-2\fboxrule\relax}{\textbf{Résultat.} #1}}\par\smallskip}
\newcommand{\popboxhead}[3]{\poppill{#1}{#2}{white}\ifx\relax#3\relax\else\hspace{6pt}{\popxbold\fontsize{10.6}{12}\selectfont\color{popink}#3}\fi\par\vspace{7pt}}

\newtcolorbox{aretenir}[1][]{popbase,colframe=popgreen,before upper={\popboxhead{À retenir}{popgreen}{#1}}}
\newtcolorbox{exemplebox}[1][]{popbase,colframe=poporange,before upper={\popboxhead{Exemple}{poporange}{#1}}}
\newtcolorbox{definition}[1][]{popbase,colframe=popblue,before upper={\popboxhead{Définition}{popblue}{#1}}}
\newtcolorbox{propriete}[1][]{popbase,colframe=poppurple,before upper={\popboxhead{Propriété}{poppurple}{#1}}}
\newtcolorbox{methode}[1][]{popbase,colframe=popgold,before upper={\popboxhead{Méthode}{popgold}{#1}}}
\newtcolorbox{attention}[1][]{popbase,colframe=poppink,before upper={\popboxhead{Attention}{poppink}{#1}}}
\newtcolorbox{experience}[1][]{popbase,colframe=popblue,colback=popblueL,before upper={\popboxhead{Expérience}{popblue}{#1}}}
% « Le savais-tu ? » : carte violette pleine (contexte, culture, Cameroun)
\newtcolorbox{savaistu}[1][]{popbase,colback=poppurple,colframe=popink,
  fontupper=\popsemi\fontsize{10.4}{14.2}\selectfont\color{white},
  before upper={\poppill{Le savais-tu\,?}{white}{poppurple}\ifx\relax#1\relax\else\hspace{6pt}{\popxbold\color{white}#1}\fi\par\vspace{7pt}}}
% Exercice résolu : énoncé en haut, \tcblower puis la solution sur fond crème
\newcounter{popexo}\newcounter{popexr}
\newtcolorbox{exoresolu}[1][]{popbase,colframe=poporange,colbacklower=popcream,
  segmentation style={popink,line width=1.2pt,dashed},
  before upper={\stepcounter{popexr}\popboxhead{Exercice résolu \thepopexr}{poporange}{#1}},
  before lower={\poppill{Solution}{popink}{popcream}\par\vspace{6pt}}}
% Exercice (non résolu en place) — la correction est dans la partie Corrigés
\newtcolorbox{exercice}[1][]{popbase,colframe=popink,
  before upper={\stepcounter{popexo}\popboxhead{Exercice \thepopexo}{popink}{#1}}}
% Corrigé
\newtcolorbox{corrige}[1][]{popbase,colframe=popgreen,colback=popgreenL,
  before upper={\popboxhead{Corrigé}{popgreen}{#1}}}
% Objectifs / prérequis / bilan
\newtcolorbox{objectifs}{popbase,colframe=popink,colback=poporangeL,
  before upper={\popboxhead{Objectifs de la leçon}{popink}{}}}
\newtcolorbox{prerequis}{popbase,colframe=popink,colback=popblueL,
  before upper={\popboxhead{Prérequis}{popblue}{}}}
\newtcolorbox{bilan}{popbase,colframe=popink,colback=white,boxrule=2.8pt,
  before upper={\popboxhead{Fiche bilan}{poporange}{L'essentiel à connaître pour l'examen}}}
% Figure encadrée avec légende
\newtcolorbox{popfigure}[1][]{enhanced,breakable=false,colback=white,colframe=popline,boxrule=1.4pt,arc=8pt,
  left=10pt,right=10pt,top=10pt,bottom=8pt,before skip=10pt,after skip=12pt,halign=center,
  lower separated=false,halign lower=center,
  fontlower=\popsemi\fontsize{9}{11.5}\selectfont\color{popmuted},
  #1}
\newcommand{\legende}[1]{\par\vspace{6pt}{\popsemi\fontsize{9}{11.5}\selectfont\color{popmuted}#1}}

% ============================================================
% Signature OnBuch+
% ============================================================
\newcommand{\poplogoinline}[1]{{\poplogo\fontsize{#1}{#1}\selectfont\color{popink}OnBuch\color{poporange}+}}
\newcommand{\signatureonbuch}{%
  \begin{tcolorbox}[enhanced,colback=popink,colframe=popink,boxrule=2.2pt,arc=10pt,
      drop shadow=poporange,left=14pt,right=14pt,top=10pt,bottom=10pt,before skip=0pt,after skip=0pt]
    \tikz[baseline=-0.7ex]\node[circle,fill=popgreen,draw=popcream,line width=1.3pt,minimum size=22pt,inner sep=0]{\popcheck{white}};%
    \hspace{9pt}\begin{minipage}[c]{0.62\linewidth}
      {\popxbold\fontsize{12}{13}\selectfont\color{popcream}Signé\ \textbullet\ {\poplogo OnBuch\color{poporange}+}}\par\vspace{2pt}
      {\raggedright\popsemi\fontsize{8}{10.5}\selectfont\color{popline}\MakeUppercase{Équipe pédagogique OnBuch+}\\\MakeUppercase{Conforme au programme officiel MINESEC}\par}
    \end{minipage}\hfill
    {\poplogo\fontsize{22}{22}\selectfont\color{popcream}OnBuch\color{poporange}+}
  \end{tcolorbox}%
}

% ============================================================
% Page de garde
% #1 titre · #2 matière · #3 niveau · #4 module · #5 n° leçon · #6 durée ·
% #7 description / objectifs (texte ou itemize)
% ============================================================
\newcommand{\pagegarde}[7]{%
  \thispagestyle{empty}
  \newgeometry{a4paper,margin=1.7cm,top=1.6cm,bottom=1.4cm}%
  \begin{tikzpicture}[remember picture,overlay]
    \fill[poporange!18] ([xshift=-3.2cm,yshift=-3.6cm]current page.north east) circle (4.6cm);
    \fill[poppurple!14] ([xshift=2.4cm,yshift=7.6cm]current page.south west) circle (3.8cm);
  \end{tikzpicture}%
  {\poplogo\fontsize{30}{30}\selectfont\color{popink}OnBuch\color{poporange}+}\hfill
  \raisebox{0.6ex}{\poppill{Cours signé \textbullet\ Premium}{popink}{popcream}}\par
  \vspace{1pt}\popsquiggle{poppurple}\par
  \vspace{8pt}
  \begin{tcolorbox}[enhanced,colback=poporange,colframe=popink,boxrule=2.8pt,arc=13pt,
      drop shadow=popink,left=18pt,right=18pt,top=12pt,bottom=13pt,after skip=10pt]
    \poppill{#2 \textbullet\ #3}{popink}{popcream}\hspace{5pt}\poppill{#4}{white}{popink}\par\vspace{8pt}
    {\popdisplay\fontsize{14}{14}\selectfont\color{popink}LEÇON #5}\par\vspace{4pt}
    {\raggedright\hyphenpenalty=10000\exhyphenpenalty=10000\popdisplay\fontsize{25}{27}\selectfont\color{popcream}\MakeUppercase{#1}\par}
    \vspace{8pt}
    {\popxbold\fontsize{9.5}{9.5}\selectfont\color{popink}\MakeUppercase{Durée conseillée : #6}}
  \end{tcolorbox}
  \begin{tcolorbox}[enhanced,colback=white,colframe=popink,boxrule=2.2pt,arc=10pt,
      drop shadow=popink,left=16pt,right=16pt,top=10pt,bottom=10pt,after skip=8pt]
    \poppill{Dans cette leçon}{poppurple}{white}\par\vspace{5pt}
    \popsemi\fontsize{9.3}{12.6}\selectfont\color{popink}#7
  \end{tcolorbox}
  \noindent\begin{minipage}[t]{0.487\linewidth}
    \begin{tcolorbox}[enhanced,colback=popgreen,colframe=popink,boxrule=2.2pt,arc=10pt,
        drop shadow=popink,left=11pt,right=11pt,top=10pt,bottom=10pt,height=3.1cm,valign=center]
      \tcbox[colback=white,colframe=popink,boxrule=1.3pt,arc=5pt,boxsep=0pt,nobeforeafter,
        left=3pt,right=3pt,top=3pt,bottom=3pt,valign=center]{\qrcode[height=1.9cm]{https://play.google.com/store/apps/details?id=com.luvvix.onbuch&hl=fr}}%
      \hspace{8pt}\begin{minipage}[c]{0.5\linewidth}
        {\popdisplay\fontsize{11}{12.5}\selectfont\color{white}\MakeUppercase{Télécharge OnBuch}}\par\vspace{4pt}
        {\raggedright\popsemi\fontsize{8.4}{11}\selectfont\color{white}Gratuit \textbullet\ Google Play\\Tuteur IA, annales, concours, résultats\par}
      \end{minipage}
    \end{tcolorbox}
  \end{minipage}\hfill
  \begin{minipage}[t]{0.487\linewidth}
    \begin{tcolorbox}[enhanced,colback=poppurple,colframe=popink,boxrule=2.2pt,arc=10pt,
        drop shadow=popink,left=11pt,right=11pt,top=10pt,bottom=10pt,height=3.1cm,valign=center]
      \tikz[baseline=-0.7ex]\node[circle,fill=white,draw=popink,line width=1.3pt,minimum size=1.6cm,inner sep=0]{\popdisplay\fontsize{24}{24}\selectfont\color{poppurple}+};%
      \hspace{8pt}\begin{minipage}[c]{0.6\linewidth}
        {\popdisplay\fontsize{11}{12.5}\selectfont\color{white}\MakeUppercase{Passe à OnBuch+}}\par\vspace{4pt}
        {\raggedright\popsemi\fontsize{8.4}{11}\selectfont\color{white}Tous les cours signés, Tuteur IA illimité, fiches PDF — onglet OnBuch+ de l'app\par}
      \end{minipage}
    \end{tcolorbox}
  \end{minipage}
  \vspace{8pt}
  \popcontenu
  \vfill
  \signatureonbuch
  \restoregeometry
  \newpage
}

% Rangée « Ce cours contient » (page de garde)
\newcommand{\poptile}[3]{% #1 couleur #2 gros texte #3 légende
  \begin{tcolorbox}[enhanced,colback=#1,colframe=popink,boxrule=2pt,arc=9pt,shadow={2.5pt}{-2.5pt}{0pt}{popink},
      left=6pt,right=6pt,top=9pt,bottom=9pt,halign=center,valign=center,height=1.75cm,width=\linewidth,nobeforeafter]
    {\popdisplay\fontsize{12.5}{13}\selectfont\color{white}\MakeUppercase{#2}}\par\vspace{3pt}
    {\centering\popsemi\fontsize{7.8}{9.5}\selectfont\color{white}#3\par}
  \end{tcolorbox}}
\newcommand{\popcontenu}{%
  \par\noindent{\popxboldls\fontsize{7.5}{7.5}\selectfont\color{popmuted}CE COURS SIGNÉ CONTIENT}\par\vspace{7pt}
  \noindent\begin{minipage}[t]{0.235\linewidth}\poptile{popblue}{Cours}{complet, illustré, pas à pas}\end{minipage}\hfill
  \begin{minipage}[t]{0.235\linewidth}\poptile{poporange}{Méthodes}{et exercices résolus}\end{minipage}\hfill
  \begin{minipage}[t]{0.235\linewidth}\poptile{poppink}{Exercices}{type examen + corrigés}\end{minipage}\hfill
  \begin{minipage}[t]{0.235\linewidth}\poptile{popgreen}{Fiche bilan}{l'essentiel à retenir}\end{minipage}\par}

% Sommaire
\newtcolorbox{sommaire}{enhanced,colback=white,colframe=popink,boxrule=2.2pt,arc=10pt,
  drop shadow=popink,left=18pt,right=18pt,top=13pt,bottom=13pt,before skip=6pt,after skip=20pt,
  before upper={\poppill{Sommaire}{poppurple}{white}\par\vspace{9pt}\popsemi\fontsize{10.6}{14.5}\selectfont\color{popink}}}

% Signature de fin de cours — poussée tout en bas de la dernière page
\newcommand{\findecours}{%
  \par\vfill\nopagebreak
  \begin{center}\popsquiggle{poporange}\end{center}\vspace{4pt}
  \signatureonbuch
}
\AtBeginDocument{\def\llbracket{\mathopen{[\mkern-3.5mu[}}\def\rrbracket{\mathclose{]\mkern-3.5mu]}}} % intervalles d'entiers (glyphe ⟦ absent de la police)
% Glyphes absents de Fira Math : redéfinis à partir de glyphes disponibles
\AtBeginDocument{%
  \def\setminus{\mathbin{\backslash}}\let\smallsetminus\setminus
  \def\vdots{\mathord{\vbox{\baselineskip3.2pt\lineskiplimit0pt\kern4pt\hbox{.}\hbox{.}\hbox{.}}}}%
  \def\ddots{\mathinner{\mkern1mu\raise6.5pt\hbox{.}\mkern2mu\raise3.6pt\hbox{.}\mkern2mu\raise0.7pt\hbox{.}\mkern1mu}}%
  \def\triangle{\mathord{\scalebox{0.9}{$\symup{\Delta}$}}}%
  \def\nabla{\mathord{\raisebox{\depth}{\scalebox{1}[-1]{$\symup{\Delta}$}}}}%
  \def\checkmark{\popcheck{popdark}}%
  \def\star{\mathbin{\ast}}\def\bigstar{\mathord{\ast}}%
  \def\diamond{\mathbin{\scalebox{0.6}{$\Diamond$}}}%
  \def\bigcup{\mathop{\mathchoice{\raisebox{-0.3em}{\scalebox{1.7}{$\cup$}}}{\scalebox{1.25}{$\cup$}}{\cup}{\cup}}\displaylimits}%
  \def\bigcap{\mathop{\mathchoice{\raisebox{-0.3em}{\scalebox{1.7}{$\cap$}}}{\scalebox{1.25}{$\cap$}}{\cap}{\cap}}\displaylimits}%
  \def\lesseqqgtr{\mathrel{\lessgtr}}\def\bigoplus{\mathop{\oplus}\displaylimits}%
}
\providecommand{\ssi}{si et seulement si} % abréviation fréquente
\AtBeginDocument{\providecommand{\wideparen}{\overparen}} % arc de cercle

%% FIGFIX-BEGIN v5 — correctifs globaux des figures (docs/onbuch-generation/tools/figfix)
\usepackage{adjustbox}\usepackage{etoolbox}\tcbuselibrary{raster}
% toute figure TikZ/pgfplots plus large que la ligne est réduite à la largeur disponible
\BeforeBeginEnvironment{tikzpicture}{\begin{adjustbox}{max width=\linewidth}}
\AfterEndEnvironment{tikzpicture}{\end{adjustbox}}
% légendes pgfplots lisibles par-dessus les courbes
\pgfplotsset{every axis/.append style={legend style={fill=white,fill opacity=0.92,text opacity=1,draw=black!15,inner sep=3pt}}}
% étiquettes d'arêtes (syntaxe edge["..."]) avec fond blanc
\tikzset{every edge quotes/.append style={fill=white,inner sep=1.2pt}}
%% FIGFIX-END
\begin{document}

\pagegarde{Ressemblances et différences au sein de l’espèce humaine}{SVTEEHB}{3ème}{SVTEEHB – classe de 3ème}{N01}{3 séances (fiche de progression harmonisée)}{Cette leçon établit, à partir d'observations et de résultats expérimentaux, les caractères qui définissent l'unité de l'espèce humaine, puis distingue les différences héréditaires des différences acquises sous l'effet de l'environnement. Elle fournit les outils pour analyser des situations concrètes de la vie courante (ressemblances familiales, jumeaux, variations individuelles).
\vspace{4pt}
\begin{multicols}{2}\small
\begin{itemize}
\item À la fin de cette leçon, je sais définir une espèce et le caractère d'un individu.
\item À la fin de cette leçon, je sais citer et illustrer les caractères communs à tous les êtres humains (unité de l'espèce).
\item À la fin de cette leçon, je sais identifier, à partir de documents, les caractères qui différencient les individus.
\item À la fin de cette leçon, je sais distinguer un caractère héréditaire d'un caractère modifié par l'environnement, en justifiant ma réponse.
\item À la fin de cette leçon, je sais interpréter des résultats d'observations (jumeaux, familles, enquêtes statistiques) pour conclure sur l'origine d'un caractère.
\item À la fin de cette leçon, je sais appliquer la notion d'unité de l'espèce humaine à des situations concrètes de la vie courante (mariages, transfusions, rejet des discriminations).
\end{itemize}\end{multicols}}

\begin{sommaire}
\begin{multicols}{2}
\begin{enumerate}
\item Diversité des individus et notion de caractère
\item Les ressemblances : les caractères de l'espèce humaine
\item Les différences entre individus : inventaire et mesure
\item Les caractères héréditaires
\item Les caractères modifiés par l'environnement
\item Bilan : unité et diversité, applications à la vie courante
\item Activité d'intégration
\item Méthodes et exercices résolus
\item Exercices
\item Corrigés des exercices
\item Fiche bilan
\end{enumerate}
\end{multicols}
\end{sommaire}

\begin{objectifs}
\begin{itemize}
\item À la fin de cette leçon, je sais définir une espèce et le caractère d'un individu.
\item À la fin de cette leçon, je sais citer et illustrer les caractères communs à tous les êtres humains (unité de l'espèce).
\item À la fin de cette leçon, je sais identifier, à partir de documents, les caractères qui différencient les individus.
\item À la fin de cette leçon, je sais distinguer un caractère héréditaire d'un caractère modifié par l'environnement, en justifiant ma réponse.
\item À la fin de cette leçon, je sais interpréter des résultats d'observations (jumeaux, familles, enquêtes statistiques) pour conclure sur l'origine d'un caractère.
\item À la fin de cette leçon, je sais appliquer la notion d'unité de l'espèce humaine à des situations concrètes de la vie courante (mariages, transfusions, rejet des discriminations).
\item À la fin de cette leçon, je sais rédiger et justifier une démarche, une réponse ou une conclusion à partir de documents fournis.
\end{itemize}
\end{objectifs}

\begin{prerequis}
\begin{itemize}
\item Notion d'espèce et de reproduction (classe de 4ème : reproduction des êtres vivants).
\item Notion de fécondation : union d'un spermatozoïde et d'un ovule (4ème).
\item Vocabulaire de base : individu, population, caractère observable.
\item Lecture et interprétation de tableaux et de graphiques simples.
\item Notion de milieu de vie et d'influence des conditions de vie (hygiène, alimentation).
\end{itemize}
\end{prerequis}

\newpage

\coursec{Diversité des individus et notion de caractère}

Au marché Mokolo à Yaoundé, Amina observe la foule qui passe devant l'étal de sa mère. Tous les clients ont deux yeux, un nez, deux bras, deux jambes, et tous marchent debout. Pourtant, aucun ne ressemble exactement à un autre : l'un est grand et mince, l'autre petit et corpulent ; ici une peau très foncée, là une peau plus claire ; au dispensaire du quartier, on apprend que les groupes sanguins varient aussi d'une personne à l'autre. Sa petite sœur, née avec une peau très claire alors que ses parents sont de teint foncé, intrigue tout le voisinage. Comment expliquer que tous les humains se ressemblent tout en étant tous différents ? Ces différences viennent-elles de nos parents ou de notre mode de vie ? C'est à ces questions que répond cette leçon.

\textbf{Problématique :} comment décrire et classer scientifiquement les ressemblances et les différences observées entre les individus de l'espèce humaine ?

\courssub{Observation : comparer des individus d'un même groupe}

\begin{experience}[Observation comparée de photographies]
\textbf{Matériel :} photographies de 5 élèves de la classe (ou de 5 habitants d'un même village), un mètre ruban, une balance, un tableau de relevés.

\textbf{Protocole :}
\begin{enumerate}
\item Observer attentivement chaque photographie et noter les traits visibles : taille, teint de la peau, forme des oreilles, forme du nez.
\item Mesurer la taille et la masse de chaque élève.
\item Consulter les carnets de santé pour connaître le groupe sanguin de chacun.
\item Reporter tous les résultats dans un tableau comparatif.
\end{enumerate}

\textbf{Résultats obtenus (exemple) :}
\begin{center}
\begin{tabularx}{\linewidth}{|l|X|X|X|X|}
\hline
\rowcolor{popblueL} \textbf{Individu} & \textbf{Taille (\si{\centi\meter})} & \textbf{Couleur de la peau} & \textbf{Groupe sanguin} & \textbf{Forme des oreilles} \\
\hline
Amina & 158 & foncé & O & lobes libres \\
\hline
Brice & 165 & foncé & A & lobes adhérents \\
\hline
Chantal & 152 & clair & B & lobes libres \\
\hline
Dieudonné & 170 & foncé & AB & lobes libres \\
\hline
Estelle & 160 & clair & O & lobes adhérents \\
\hline
\end{tabularx}
\end{center}

\textbf{Interprétation :} certains traits sont identiques chez tous (deux yeux, deux bras, station debout) : ce sont des traits \cle{communs}. D'autres varient d'un individu à l'autre (taille, teint, groupe sanguin, forme des oreilles) : ce sont des traits \cle{variables}.

\textbf{Conclusion :} les individus de l'espèce humaine présentent à la fois des ressemblances et des différences. Pour les décrire rigoureusement, il faut une notion précise : la notion de \cle{caractère}.
\end{experience}

\courssub{La notion de caractère}

Pour comparer les élèves de la classe, nous avons regardé leur taille, leur teint, leur groupe sanguin... Chacun de ces traits est ce qu'on appelle un \cle{caractère}. C'est l'outil de base du scientifique pour décrire un individu.

\begin{definition}[Caractère]
Un \cle{caractère} est tout trait \textbf{observable} ou \textbf{mesurable} qui décrit un individu. On distingue trois grandes catégories de caractères :
\begin{itemize}
\item les caractères \cle{morphologiques} : ils concernent la forme et l'aspect du corps (taille, couleur de la peau, forme du nez, des lèvres, des oreilles, empreintes digitales) ;
\item les caractères \cle{physiologiques} : ils concernent le fonctionnement du corps (tension artérielle, sécrétion d'enzymes digestives, rythme cardiaque) ;
\item les caractères \cle{biochimiques} : ils concernent la composition chimique du corps (groupe sanguin A, B, AB ou O, présence de certaines protéines dans le sang).
\end{itemize}
\end{definition}

\begin{popfigure}
\begin{center}
\begin{tikzpicture}[scale=1]
% Silhouette humaine simplifiée
\draw[fill=popcream,draw=popdark,thick] (0,3.4) circle (0.55); % tête
\draw[fill=popcream,draw=popdark,thick,rounded corners=6pt] (-0.7,0.6) rectangle (0.7,2.7); % tronc
\draw[fill=popcream,draw=popdark,thick,rounded corners=4pt] (-1.5,1.0) rectangle (-0.75,2.5); % bras gauche
\draw[fill=popcream,draw=popdark,thick,rounded corners=4pt] (0.75,1.0) rectangle (1.5,2.5); % bras droit
\draw[fill=popcream,draw=popdark,thick,rounded corners=4pt] (-0.65,-1.6) rectangle (-0.1,0.65); % jambe gauche
\draw[fill=popcream,draw=popdark,thick,rounded corners=4pt] (0.1,-1.6) rectangle (0.65,0.65); % jambe droite
% Annotations
\draw[->,popblue,thick] (0.55,3.6) -- (2.6,4.3) node[right,align=left,font=\small]{Forme du nez,\\ couleur des yeux};
\draw[->,poporange,thick] (-0.6,3.3) -- (-2.6,3.9) node[left,align=right,font=\small]{Couleur de la peau\\ (sur tout le corps)};
\draw[->,popgreen,thick] (0.7,2.2) -- (2.6,2.6) node[right,align=left,font=\small]{Groupe sanguin\\ (A, B, AB, O)};
\draw[->,poppurple,thick] (-1.5,1.2) -- (-2.6,1.4) node[left,align=right,font=\small]{Empreintes digitales\\ (au bout des doigts)};
\draw[->,popgold,thick] (0.4,-1.5) -- (2.6,-1.5) node[right,align=left,font=\small]{Taille et masse\\ (caractères mesurables)};
% Mesure de la taille
\draw[<->,popmuted,thick] (-2.1,-1.6) -- (-2.1,3.95) node[midway,left,font=\small,rotate=90]{taille};
\end{tikzpicture}
\end{center}
\legende{Silhouette humaine annotée : quelques caractères observables ou mesurables chez un individu. Les caractères morphologiques (forme du nez, couleur de la peau, empreintes digitales) sont directement visibles ; les caractères physiologiques et biochimiques (groupe sanguin) nécessitent des analyses.}
\end{popfigure}

\courssub{Caractères qualitatifs et caractères quantitatifs}

Tous les caractères ne se décrivent pas de la même façon. La couleur des yeux se décrit par des mots (noirs, marron, verts...) ; la taille se décrit par un nombre accompagné d'une unité (158 cm). On distingue donc deux types de caractères.

\begin{definition}[Caractères qualitatifs et quantitatifs]
\begin{itemize}
\item Un caractère \cle{qualitatif} se décrit par des \textbf{qualités}, sans mesure chiffrée : couleur des yeux, groupe sanguin, forme des oreilles.
\item Un caractère \cle{quantitatif} se décrit par une \textbf{mesure chiffrée} avec une unité : la taille (en \si{\centi\meter}), la masse (en \si{\kilo\gram}), la tension artérielle (en \si{\milli\meter\of{Hg}}).
\end{itemize}
\end{definition}

\begin{exemplebox}[Quelques caractères de l'espèce humaine]
\begin{center}
\begin{tabularx}{\linewidth}{|l|X|X|}
\hline
\rowcolor{popblueL} \textbf{Caractère} & \textbf{Type} & \textbf{Exemples de valeurs} \\
\hline
Couleur des yeux & qualitatif & noirs, marron, verts \\
\hline
Groupe sanguin & qualitatif & A, B, AB, O \\
\hline
Forme des oreilles & qualitatif & lobes libres ou adhérents \\
\hline
Taille & quantitatif & \SI{152}{cm}, \SI{165}{cm}, \SI{170}{cm} \\
\hline
Masse & quantitatif & \SI{45}{kg}, \SI{60}{kg} \\
\hline
Tension artérielle & quantitatif & \SI{120}{mmHg} / \SI{80}{mmHg} \\
\hline
\end{tabularx}
\end{center}
\end{exemplebox}

\begin{methode}[Comment comparer scientifiquement des individus]
\begin{enumerate}
\item Dresser d'abord la \textbf{liste des caractères} à étudier (taille, teint, groupe sanguin...).
\item Pour chaque caractère, relever sa valeur chez chaque individu (observation ou mesure).
\item Classer chaque caractère en deux catégories : \textbf{commun} (identique chez tous) ou \textbf{variable} (différent d'un individu à l'autre).
\item Conclure : les caractères communs définissent l'espèce ; les caractères variables traduisent la diversité des individus.
\end{enumerate}
\end{methode}

\begin{attention}[Ne pas confondre caractère et modalité]
Erreur fréquente au BEPC : confondre le \cle{caractère} et sa \cle{modalité} (la valeur prise par le caractère).
\begin{itemize}
\item « Couleur de la peau » est un \textbf{caractère} ; « peau claire » et « peau foncée » sont des \textbf{modalités} de ce caractère.
\item « Groupe sanguin » est un caractère ; « A », « B », « AB », « O » sont ses modalités.
\end{itemize}
Dans une copie, écris toujours : « le caractère \emph{taille} a pour modalités \SI{152}{cm}, \SI{165}{cm}... » et non « le caractère est \SI{152}{cm} ».
\end{attention}

\courssub{Chaque individu est unique : la notion d'individualité}

Reprenons notre tableau de relevés : aucun des 5 élèves ne possède exactement la même combinaison de taille, de teint, de groupe sanguin et de forme des oreilles. Si l'on ajoutait d'autres caractères (empreintes digitales, forme du nez, couleur des yeux...), il deviendrait pratiquement impossible de trouver deux individus identiques.

\begin{aretenir}[Individualité]
Chaque individu possède un \textbf{ensemble unique de caractères} : c'est la notion d'\cle{individualité}. Même les vrais jumeaux, qui se ressemblent beaucoup, présentent des différences, par exemple dans leurs empreintes digitales.
\end{aretenir}

\begin{exemplebox}[Les empreintes digitales, une signature unique]
Les dessins de la peau au bout des doigts (empreintes digitales) se forment avant la naissance et ne changent plus toute la vie. Ils permettent de distinguer deux individus avec une fiabilité quasi absolue : on n'a jamais trouvé deux personnes avec les mêmes empreintes, \textbf{pas même les vrais jumeaux}. C'est pourquoi la police et les services d'état civil les utilisent pour identifier les personnes.
\end{exemplebox}

\begin{savaistu}[8 milliards d'individus, aucun identique]
Il existe environ 8 milliards d'êtres humains sur Terre. Pourtant, aucun n'est strictement identique à un autre. Au Cameroun, les cartes nationales d'identité et les passeports utilisent les empreintes digitales et la photographie du visage pour reconnaître chaque citoyen de façon sûre : la diversité des caractères humains est ainsi mise au service de la vie quotidienne.
\end{savaistu}

\begin{exoresolu}[Identifier caractères et modalités]
\textbf{Énoncé.} Voici la description de deux élèves : « Paul mesure \SI{162}{cm}, il a la peau foncée, les yeux noirs et le groupe sanguin O. Marie mesure \SI{155}{cm}, elle a la peau claire, les yeux marron et le groupe sanguin A. »
\begin{enumerate}
\item Dresse la liste des caractères étudiés.
\item Pour chaque caractère, indique s'il est commun ou variable chez ces deux élèves, et précise ses modalités.
\item Classe ces caractères en qualitatifs et quantitatifs.
\end{enumerate}
\tcblower
\textbf{Solution détaillée.}
\begin{enumerate}
\item Les caractères étudiés sont : la \textbf{taille}, la \textbf{couleur de la peau}, la \textbf{couleur des yeux} et le \textbf{groupe sanguin}.
\item Les quatre caractères sont \textbf{variables} car leurs valeurs diffèrent entre Paul et Marie :
\begin{itemize}
\item taille : \SI{162}{cm} chez Paul, \SI{155}{cm} chez Marie ;
\item couleur de la peau : foncée chez Paul, claire chez Marie ;
\item couleur des yeux : noirs chez Paul, marron chez Marie ;
\item groupe sanguin : O chez Paul, A chez Marie.
\end{itemize}
\item Caractères \textbf{quantitatifs} (mesurables) : la taille. Caractères \textbf{qualitatifs} (décrits par des qualités) : la couleur de la peau, la couleur des yeux, le groupe sanguin.
\end{enumerate}
\textbf{Conclusion :} ces deux élèves appartiennent bien à la même espèce (ils partagent tous les caractères humains), mais leurs modalités différent pour chaque caractère étudié : ils sont uniques.
\end{exoresolu}

\begin{exercice}[À toi de jouer]
Observe trois camarades de ta classe. 1) Choisis quatre caractères (deux morphologiques, un physiologique ou biochimique, un quantitatif). 2) Construis un tableau de relevés. 3) Indique pour chaque caractère s'il est commun ou variable, et justifie ta réponse.
\end{exercice}

\begin{corrige}[Exercice 1]
Exemple de réponse attendue. Caractères choisis : taille (morphologique et quantitatif), forme des oreilles (morphologique), groupe sanguin (biochimique), couleur des yeux (morphologique). Le tableau présente en lignes les trois camarades et en colonnes les quatre caractères. La taille varie (par exemple \SI{150}{cm}, \SI{158}{cm}, \SI{163}{cm}) : caractère variable. La forme des oreilles varie (lobes libres ou adhérents) : caractère variable. Le groupe sanguin varie (A, O, B...) : caractère variable. En revanche, des caractères comme « deux yeux », « deux bras », « station debout » seraient communs à tous. La justification repose sur la comparaison des modalités : un caractère est commun si sa modalité est identique chez tous, variable sinon.
\end{corrige}

\begin{aretenir}[Bilan de la section]
\begin{itemize}
\item Un \cle{caractère} est un trait observable ou mesurable d'un individu ; il peut être \textbf{morphologique}, \textbf{physiologique} ou \textbf{biochimique}, \textbf{qualitatif} ou \textbf{quantitatif}.
\item Chaque caractère possède des \cle{modalités} (ses valeurs possibles).
\item Certains caractères sont \textbf{communs} à tous les humains, d'autres sont \textbf{variables}.
\item L'ensemble des caractères d'un individu est \textbf{unique} : c'est l'\cle{individualité}.
\end{itemize}
Dans la suite de la leçon, nous verrons que les caractères communs définissent l'espèce humaine, puis nous rechercherons l'origine des différences : hérédité ou environnement ?
\end{aretenir}

\coursec{Les ressemblances : les caractères de l'espèce humaine}

Dans la section précédente, nous avons découvert que chaque individu humain est unique : la diversité des caractères (taille, couleur de peau, groupe sanguin, etc.) est une évidence quotidienne. Pourtant, si l'on observe un nouveau-né à Yaoundé, un paysan dans le Nord, un élève à Douala ou un chercheur à Paris, on reconnaît immédiatement qu'ils appartiennent tous à la même \emph{espèce}. Comment définir rigoureusement cette unité profonde qui transcende les différences visibles ? C'est l'objet de cette section : identifier les \textbf{caractères de l'espèce humaine}, ces traits invariants partagés par les 8 milliards d'humains que compte la Terre.

\courssub{Définition biologique de l'espèce et critère de fécondité}

Avant de lister les caractères communs, rappelons ce qu'est une espèce en biologie. La notion est fondamentale car elle fournit le critère ultime pour décider si deux individus appartiennent ou non au même groupe biologique.

\begin{definition}[Espèce biologique (concept de Mayr)]
Une \cle{espèce} est un ensemble d'individus \textbf{semblables} par leurs caractères morphologiques, anatomiques, physiologiques et biochimiques, \textbf{capables de se reproduire entre eux} (fécondité naturelle) et \textbf{donnant des descendants eux-mêmes féconds}.
\end{definition}

\begin{attention}[Piège fréquent : confusion espèce / race / variété]
En biologie moderne, le terme \emph{race} n'a pas de valeur taxinomique pour l'espèce humaine. Il existe des \cle{races} chez les animaux domestiques (sélection artificielle par l'homme : races de chiens, de poules, de bovins comme le \emph{Boran} ou le \emph{Goudali} au Cameroun), mais \textbf{une seule espèce humaine} : \emph{Homo sapiens}. Les différences de couleur de peau, de morphologie faciale ou de texture de cheveux sont des \textbf{variations intraspécifiques}, non des frontières biologiques entre espèces ou sous-espèces.
\end{attention}

La preuve la plus directe et la plus indiscutable de l'unité de l'espèce humaine est la \textbf{fécondité des croisements} entre populations géographiquement et culturellement très éloignées.

\begin{experience}[Observation : les mariages interethniques et mixtes au Cameroun et dans le monde]
\textbf{Mise en situation.} Au Cameroun, pays de plus de 250 groupes ethniques, les mariages entre personnes de groupes ethniques différents (ex. Bamiléké et Bassa, Peul et Beti, Douala et Mafa) sont quotidiens. De même, les couples mixtes (Camerounais/Européen, Camerounais/Asiatique, etc.) sont fréquents dans nos villes.

\textbf{Observation (fait constant).} Tous ces couples ont des enfants qui grandissent normalement et qui, devenus adultes, sont \textbf{féconds} à leur tour : ils peuvent avoir des enfants avec n'importe quel autre humain, quelle que soit son origine.

\textbf{Interprétation.} Selon la définition de l'espèce, la fécondité naturelle et la fécondité des descendants (sur plusieurs générations) prouvent qu'il n'existe \textbf{aucune barrière reproductive} entre les populations humaines.

\textbf{Conclusion.} Tous les humains actuels appartiennent à une \textbf{seule et unique espèce biologique} : \emph{Homo sapiens}.
\end{experience}

\begin{propriete}[Unité de l'espèce humaine -- Critère décisif]
Tous les êtres humains actuels, quelles que soient leur origine géographique, leur couleur de peau, leur culture ou leur histoire, appartiennent à l'espèce unique \textbf{\emph{Homo sapiens}}. Le critère décisif est la \textbf{fécondité des croisements} entre toutes les populations humaines, sans exception.
\end{propriete}

\begin{savaistu}[Le Cameroun, carrefour de la diversité humaine]
Le Cameroun est souvent appelé « l'Afrique en miniature » pour sa diversité géographique, mais aussi humaine. On y trouve des populations de taille très grande (Peuls du Nord, \SI{1,85}{m} en moyenne chez les hommes) et de taille petite (Pygmées Baka du Sud-Est, \SI{1,50}{m}). Pourtant, les mariages entre Peuls et Baka, bien que rares pour des raisons culturelles, donnent des enfants parfaitement viables et féconds. Cette réalité de terrain illustre concrètement l'unité biologique de notre espèce.
\end{savaistu}

\courssub{Les caractères morphologiques communs : le plan d'organisation externe}

Malgré la diversité des tailles et des apparences, tout être humain partage un \textbf{plan d'organisation corporel} identique. C'est ce qui permet à un anatomiste de reconnaître instantanément un squelette humain parmi ceux d'autres primates.

\begin{popfigure}
\begin{tikzpicture}[scale=0.7, every node/.style={font=\small}]
% Human skeleton (simplified)
\def\humanx{0}
\def\humany{0}
% Vertebral column
\draw[popdark, line width=2pt] (\humanx, \humany) -- (\humanx, \humany+14) node[midway, left=2mm] {Colonne vertébrale};
% Skull
\draw[popdark, line width=1.5pt] (\humanx-1.2, \humany+14) ellipse (1.2 and 1.5) node[above=2mm] {Crâne volumineux};
% Pelvis (human: broad, short)
\draw[popblue, line width=2pt] (\humanx-2.5, \humany+7) -- (\humanx+2.5, \humany+7) node[midway, above=1mm, popblue] {Bassin large et court};
\draw[popblue, line width=1.5pt] (\humanx-2.5, \humany+7) -- (\humanx-2.5, \humany+5);
\draw[popblue, line width=1.5pt] (\humanx+2.5, \humany+7) -- (\humanx+2.5, \humany+5);
% Femurs (vertical)
\draw[popdark, line width=2pt] (\humanx-1.2, \humany+7) -- (\humanx-1.2, \humany+0.5);
\draw[popdark, line width=2pt] (\humanx+1.2, \humany+7) -- (\humanx+1.2, \humany+0.5);
% Tibias
\draw[popdark, line width=1.5pt] (\humanx-1.2, \humany+0.5) -- (\humanx-1.2, \humany-1.5);
\draw[popdark, line width=1.5pt] (\humanx+1.2, \humany+0.5) -- (\humanx+1.2, \humany-1.5);
% Feet (arched)
\draw[popdark, line width=1.5pt] (\humanx-1.5, \humany-1.5) -- (\humanx-0.5, \humany-1.5) -- (\humanx-1.2, \humany-1) -- cycle;
\draw[popdark, line width=1.5pt] (\humanx+0.5, \humany-1.5) -- (\humanx+1.5, \humany-1.5) -- (\humanx+1.2, \humany-1) -- cycle;
% Arms
\draw[popdark, line width=1.5pt] (\humanx-2.5, \humany+12) -- (\humanx-5, \humany+5);
\draw[popdark, line width=1.5pt] (\humanx+2.5, \humany+12) -- (\humanx+5, \humany+5);
% Hands (opposable thumb suggested)
\draw[popdark, fill=popcream] (\humanx-5, \humany+5) circle (0.3);
\draw[popdark, fill=popcream] (\humanx+5, \humany+5) circle (0.3);

% Great ape skeleton (simplified, to the right)
\def\apex{12}
\def\apey{0}
% Vertebral column (curved)
\draw[popdark, line width=2pt] (\apex, \apey+2) .. controls (\apex+1, \apey+6) and (\apex-0.5, \apey+10) .. (\apex, \apey+13) node[midway, right=2mm] {Colonne vertébrale en arc};
% Skull (prognathic)
\draw[popdark, line width=1.5pt] (\apex-1, \apey+13) ellipse (1.3 and 1.2) node[above=2mm] {Crâne prognathe};
\draw[popdark, line width=1.5pt] (\apex-1.5, \apey+12.5) -- (\apex-2.5, \apey+11.5); % Face projection
% Pelvis (ape: narrow, long)
\draw[poporange, line width=2pt] (\apex-1.5, \apey+6) -- (\apex+1.5, \apey+6) node[midway, above=1mm, poporange] {Bassin étroit et long};
\draw[poporange, line width=1.5pt] (\apex-1.5, \apey+6) -- (\apex-1.5, \apey+4);
\draw[poporange, line width=1.5pt] (\apex+1.5, \apey+6) -- (\apex+1.5, \apey+4);
% Femurs (angled)
\draw[popdark, line width=2pt] (\apex-1, \apey+6) -- (\apex-1.5, \apey+0);
\draw[popdark, line width=2pt] (\apex+1, \apey+6) -- (\apex+1.5, \apey+0);
% Arms (long, reach knees)
\draw[popdark, line width=1.5pt] (\apex-2, \apey+11) -- (\apex-4.5, \apey+3);
\draw[popdark, line width=1.5pt] (\apex+2, \apey+11) -- (\apex+4.5, \apey+3);
% Hands (long fingers, short thumb)
\draw[popdark, fill=popcream] (\apex-4.5, \apey+3) circle (0.3);
\draw[popdark, fill=popcream] (\apex+4.5, \apey+3) circle (0.3);

% Labels
\node[align=center, font=\bfseries\large, popblue] at (\humanx, \humany+16) {Humain (\emph{Homo sapiens})};
\node[align=center, font=\bfseries\large, poporange] at (\apex, \apey+16) {Grand singe (ex. Chimpanzé)};
% Arrows for key differences
\draw[->, popblue, thick] (\humanx-3.5, \humany+7) -- (\humanx-4.5, \humany+7) node[left] {Bassin adapté à la bipédie};
\draw[->, poporange, thick] (\apex+2.5, \apey+6) -- (\apex+3.5, \apey+6) node[right] {Bassin pour quadrupédie};
\draw[->, popblue, thick] (\humanx-1.2, \humany+3) -- (\humanx-2.5, \humany+2) node[left] {Fémurs verticaux};
\draw[->, poporange, thick] (\apex-1, \apey+3) -- (\apex-2.5, \apey+2) node[left] {Fémurs obliques};
\draw[->, popblue, thick] (\humanx-5, \humany+5.5) -- (\humanx-6, \humany+5.5) node[left] {Pouce long et opposable};
\draw[->, poporange, thick] (\apex+4.5, \apey+3.5) -- (\apex+5.5, \apey+3.5) node[right] {Pouce court, peu opposable};
\end{tikzpicture}
\legende{Comparaison schématique du squelette humain (gauche) et d'un grand singe (droite). Les adaptations à la bipédie obligatoire chez l'humain : bassin large et court (popblue), fémurs verticaux rapprochant les genoux du centre de gravité, pied arqué avec gros orteil aligné, pouce long et fortement opposable. Chez le singe : bassin étroit et long (poporange), fémurs obliques, mains et pieds préhensiles pour l'arboricolie.}
\end{popfigure}

\begin{aretenir}[Caractères morphologiques de l'espèce humaine]
\begin{itemize}
    \item \textbf{Station debout (bipédie) permanente :} le corps est vertical, le poids supporté par les deux membres inférieurs.
    \item \textbf{Bassin large, court et en entonnoir :} supporte les viscères et permet l'insertion des muscles fessiers (stabilisation du tronc sur un pied).
    \item \textbf{Colonne vertébrale à double courbure (cyphose dorsale, lordose lombaire) :} amortit les chocs de la marche et centre la gravité.
    \item \textbf{Membres inférieurs allongés, fémurs verticaux :} genoux rapprochés sous le centre de gravité.
    \item \textbf{Pied arqué (voûte plantaire) avec gros orteil aligné et non opposable :} levier de propulsion et amortisseur.
    \item \textbf{Deux membres supérieurs libérés de la locomotion :} préhension fine, manipulation d'outils.
    \item \textbf{Main avec pouce long, fort et pleinement opposable :} pince de précision (pouce/index) et pince de force.
    \item \textbf{Face plane (orthognathie), front fuyant, arcades sourcilières réduites.}
    \item \textbf{Pilosité corporelle réduite (cheveux longs sur le crâne, duvet fin ailleurs).}
\end{itemize}
\end{aretenir}

\courssub{Les caractères anatomiques et physiologiques communs : unité du plan interne}

L'unité ne s'arrête pas à l'enveloppe externe. L'organisation interne — squelette, viscères, système nerveux — suit un \textbf{plan d'organisation unique}.

\begin{itemize}
    \item \textbf{Squelette :} même nombre d'os (\SI{206}{chez l'adulte}), même nomenclature (fémur, humérus, vertèbres cervicales/thoraciques/lombaires, etc.), même articulation.
    \item \textbf{Appareils d'organes :} même disposition topographique (cœur à gauche dans le thorax, foie à droite dans l'abdomen, estomac en J, reins rétropéritonéaux, etc.), même vascularisation et innervation principales.
    \item \textbf{Système nerveux central :} un \textbf{cerveau volumineux} (moyenne \SI{1350}{cm^3} chez l'adulte), fortement plissé (circonvolutions), avec un cortex cérébral très développé (siège des fonctions cognitives supérieures : langage abstrait, raisonnement, planification, culture).
    \item \textbf{Langage articulé :} appareil phonatoire (larynx bas, cavité buccale modifiée, langue mobile, contrôle respiratoire fin) permettant une infinité de sons combinés en mots et phrases syntaxiques. Aucune autre espèce n'a un langage symbolique, syntaxique et transmissible culturellement d'une telle complexité.
\end{itemize}

\begin{exemplebox}[Exemple concret : la transplantation d'organes]
La possibilité de greffer un rein, un foie, un cœur ou une cornée d'un donneur décédé à un receveur vivant, \textbf{quelle que soit leur origine ethnique ou géographique}, repose entièrement sur l'unité anatomique, histologique et immunologique de base de l'espèce. Seule la compatibilité du \textbf{système HLA} (antigènes leucocytaires humains, très polymorphes) conditionne le succès, non l'appartenance à un « groupe racial ». Un Camerounais peut recevoir le rein d'un Européen compatible HLA, et inversement.
\end{exemplebox}

\courssub{Les caractères biochimiques et génétiques communs : l'unité au niveau moléculaire}

C'est au niveau microscopique et moléculaire que l'unité est la plus frappante et la plus rigoureuse.

\begin{aretenir}[Caractères biochimiques et génétiques invariants]
\begin{itemize}
    \item \textbf{Type cellulaire unique :} toutes nos cellules sont des \textbf{cellules eucaryotes} (noyau vrai, organites : mitochondries, réticulum, Golgi, etc.). Pas de paroi cellulosique, pas de chloroplaste.
    \item \textbf{Mêmes grandes molécules organiques :} protéines (mêmes 20 acides aminés protéinogènes), glucides, lipides, acides nucléiques (ADN, ARN). Mêmes voies métaboliques fondamentales (glycolyse, cycle de Krebs, phosphorylation oxydative, synthèse des protéines).
    \item \textbf{Même code génétique :} universel (presque), triplet de bases $\rightarrow$ acide aminé.
    \item \textbf{Même caryotype :} \textbf{46 chromosomes} organisés en \textbf{23 paires} (22 paires d'autosomes + 1 paire de gonosomes XX ou XY) dans le noyau de chaque cellule somatique des individus sains.
\end{itemize}
\end{aretenir}

\begin{exemplebox}[Exemple chiffré : le caryotype humain universel]
\begin{center}
\begin{tabularx}{\linewidth}{|l|X|X|}\hline
\rowcolor{popblueL}
\textbf{Paramètre} & \textbf{Valeur} & \textbf{Commentaire} \\ \hline
Nombre de chromosomes (cellule somatique) & $2n = 46$ & Constant chez tout humain sain \\ \hline
Nombre de paires & 23 paires & 22 autosomes + 1 paire gonosomes \\ \hline
Gonosomes & XX (femelle) ou XY (mâle) & Détermination du sexe génétique \\ \hline
Taille du génome haploïde & $\approx 3,2 \times 10^9$ paires de bases & $\approx 3,2$ Gb \\ \hline
Nombre de gènes codants & $\approx 20\,000$ à $21\,000$ & Estimation actuelle (projet ENCODE) \\ \hline
\end{tabularx}
\end{center}
Ce caryotype est \textbf{identique} chez un enfant Baka du Sud-Cameroun, un élève de Douala, un paysan du Nord ou un new-yorkais. Seules de rares anomalies chromosomiques (trisomie 21, syndrome de Turner, etc.) en dévient, et ce sont des pathologies, non des variations normales de l'espèce.
\end{exemplebox}

\begin{savaistu}[99,9 \% d'ADN partagé : la diversité tient dans 0,1 \%]
Deux êtres humains pris au hasard dans la population mondiale partagent en moyenne \textbf{99,9 \%} de leur séquence d'ADN. Toute la diversité phénotypique observable (taille, couleur de peau, forme du nez, groupes sanguins, susceptibilité à certaines maladies) est codée dans les \textbf{0,1 \% restants} (environ 3 millions de paires de bases variables, essentiellement des polymorphismes nucléotidiques simples -- SNP). C'est une preuve moléculaire éclatante de notre unité récente et commune.
\end{savaistu}

\courssub{Preuve fondamentale : la fécondité inter-populations illimitée}

Nous l'avons évoqué en introduction : le critère \emph{biologique} de l'espèce (Mayr) est la fécondité des croisements. Pour l'espèce humaine, ce critère est \textbf{absolument vérifié à l'échelle planétaire}.

\begin{experience}[Démarche d'investigation : tester l'appartenance à la même espèce]
\textbf{Question :} Les populations humaines isolées depuis des millénaires (ex. Amérindiens d'Amazonie, Papous de Nouvelle-Guinée, Inuits de l'Arctique, Pygmées d'Afrique centrale, Européens, Asiatiques) appartiennent-elles à la même espèce ?

\textbf{Hypothèse :} Si elles appartiennent à la même espèce, les croisements entre individus de ces populations doivent donner des descendants viables et \textbf{féconds}.

\textbf{Document / Observation historique et actuelle :}
\begin{itemize}
    \item Depuis les grandes explorations (XVe--XVIe s.) et les migrations modernes, les mélanges de populations sont massifs et continus.
    \item Aucune barrière de stérilité (prézygotique ou postzygotique) n'a jamais été observée entre quelque populations humaines que ce soit.
    \item Les enfants métis (terme descriptif, non péjoratif) sont la norme dans l'histoire humaine : populations d'Amérique latine (métissage amérindien/européen/africain), du Cap-Vert, de La Réunion, de Madagascar, du Maghreb, d'Afrique du Sud, etc.
    \item Au Cameroun : enfants de couples Bamiléké/Bassa, Peul/Mafa, Camerounais/Français, Camerounais/Chinois, etc. Tous grandissent normalement et ont leur propre descendance.
\end{itemize}

\textbf{Interprétation :} L'absence totale d'isolement reproductif entre \textbf{toutes} les populations humaines actuelles est un fait biologique massif.

\textbf{Conclusion :} \emph{Homo sapiens} est une \textbf{seule espèce biologique}, non subdivisée en sous-espèces ou races biologiques.
\end{experience}

\begin{popfigure}
\begin{tikzpicture}[scale=0.8, every node/.style={font=\small}]
% Phylogenetic tree: single branch for Homo sapiens
\node[align=center, font=\bfseries\large, popdark] at (0, 5.5) {Arbre phylogénétique simplifié des Homininés récents};

% Common ancestor
\node[draw, ellipse, fill=popgoldL, thick, align=center] (anc) at (0, 4.5){Ancêtre commun\\ \emph{Homo} sp. (ex. \emph{H. heidelbergensis})\\ $\sim$ 600 000 -- 800 000 ans};

% Branches for extinct species
\draw[thick, popmuted] (anc.south west) -- ++(-2.5, -1.5) node[left, popmuted, align=center]{\emph{H. neanderthalensis}\\(éteint $\sim$ 40 ka)};
\draw[thick, popmuted] (anc.south) -- ++(0, -1.5) node[below, popmuted, align=center]{\emph{H. denisova}\\(éteint $\sim$ 50 ka)};
\draw[thick, popmuted] (anc.south east) -- ++(2.5, -1.5) node[right, popmuted, align=center]{Autres lignées\\ africaines éteintes};

% Homo sapiens branch - single trunk
\draw[very thick, popblue] (anc.south) -- (0, 1.5) node[midway, right=3mm, popblue, font=\bfseries] {\emph{Homo sapiens}};

% Modern human populations as twigs at the very top
\node[draw, rectangle, rounded corners, fill=popblue!10, thick, popblue] (hs) at (0, 1) {\textbf{\emph{Homo sapiens} (actuel)}};

% Twigs representing populations (no branching = no subspecies)
\foreach \name/\x in {Afrique subsaharienne/-3.5, Afrique du Nord/-1.5, Europe/0.5, Asie de l'Est/2.5, Océanie/4.5, Amériques/6.5} {
    \draw[popblue, thick] (hs.north) -- ++(\x-0, 1.2) node[above, align=center, popdark] {\name};
}

% Note on interbreeding
\node[draw, rectangle, rounded corners, fill=popgreenL, align=center, font=\small] at (0, -1.5) {
    \textbf{Fécondité totale entre toutes les populations actuelles}\\
    $\rightarrow$ Une seule espèce, aucune sous-espèce biologique.
};
\end{tikzpicture}
\legende{Arbre phylogénétique montrant l'émergence unique de \emph{Homo sapiens} en Afrique ($\sim$300 000 ans) et sa dispersion mondiale. Toutes les populations actuelles sont des « rameaux terminaux » d'un même tronc unique, sans divergence spéciative. Les croisements avec Néandertaliens et Dénisoviens ont laissé des traces (1--4 \% d'ADN), mais ces lignées sont éteintes.}
\end{popfigure}

\courssub{Bilan : unité biologique, diversité culturelle, conséquence éthique}

\begin{aretenir}[Synthèse : les caractères de l'espèce humaine]
L'espèce humaine \emph{Homo sapiens} est définie par un \textbf{ensemble de caractères invariants} partagés par tous ses membres :
\begin{enumerate}
    \item \textbf{Morphologiques :} bipédie obligatoire, bassin large/court, main préhensile à pouce opposable, face plane, pilosité réduite.
    \item \textbf{Anatomiques/Physiologiques :} même plan d'organisation des organes, même squelette, même cerveau volumineux et plissé, langage articulé.
    \item \textbf{Biochimiques/Génétiques :} cellules eucaryotes, même métabolisme, même code génétique, \textbf{caryotype 46,XX/XY universel}, 99,9 \% d'ADN identique.
    \item \textbf{Reproductif (critère décisif) :} fécondité naturelle et fécondité des descendants entre \textbf{toutes} les populations sans exception.
\end{enumerate}
\end{aretenir}

\begin{propriete}[Conséquence majeure : le racisme n'a aucune base biologique]
Puisque \textbf{tous les humains appartiennent à une seule et même espèce}, partageant le même patrimoine génétique fondamental et une fécondité croisée totale, \textbf{la notion de « races humaines » au sens biologique (sous-espèces) est scientifiquement invalide}. Les classifications raciales du XIXe siècle (basées sur la couleur de peau, la forme du crâne, etc.) sont des constructions sociales et historiques, non des réalités biologiques. La biologie moderne fonde ainsi rigoureusement le \textbf{rejet du racisme et de toute discrimination fondée sur l'origine ou l'apparence physique}.
\end{propriete}

\begin{exoresolu}[Analyse d'un document : unité et diversité]
\textbf{Énoncé.} Un journal publie : « Les scientifiques découvrent que les Africains et les Européens n'ont pas les mêmes gènes, preuve qu'ils sont des races différentes. » En vous appuyant sur les connaissances du cours, critiquez cette affirmation.

\textbf{Solution détaillée.}
\begin{enumerate}
    \item \textbf{Erreur sur le terme « gènes » :} Tous les humains possèdent les \textbf{mêmes gènes} (mêmes loci, mêmes fonctions : gène de l'hémoglobine, gène de l'insuline, gènes du développement cérébral, etc.). Ce qui diffère, ce sont les \textbf{allèles} (versions alternatives d'un même gène) et leur fréquence dans les populations.
    \item \textbf{Erreur sur la proportion :} Les différences alléliques ne concernent qu'une infime fraction du génome (\SI{0,1}{\%}). 99,9 \% de la séquence d'ADN est \textbf{identique}.
    \item \textbf{Erreur logique (critère d'espèce) :} L'existence de fréquences alléliques différentes (ex. allèle de la drépanocytose plus fréquent en zone paludéenne) ne crée \textbf{aucune barrière reproductive}. Les mariages Afrique/Europe donnent des enfants féconds. Donc : \textbf{même espèce}.
    \item \textbf{Conclusion :} L'affirmation du journal confond \textbf{variabilité génétique intraspécifique} (normale, utile pour l'adaptation) et \textbf{spéciation}. Elle est scientifiquement fausse et véhicule un préjugé raciste.
\end{enumerate}
\end{exoresolu}

\begin{methode}[Comment argumenter l'unité de l'espèce humaine (type BEPC)]
Pour répondre à une question du type « Montrez que tous les humains appartiennent à la même espèce », structurez votre réponse en \textbf{trois arguments indépendants} :
\begin{enumerate}
    \item \textbf{Argument morpho-anatomique :} Citez 3 à 4 caractères morphologiques invariants (bipédie, bassin, main, cerveau, face plane).
    \item \textbf{Argument biochimique/génétique :} Citez l'universalité du caryotype (46 chromosomes), du code génétique, du métabolisme, et le pourcentage d'ADN partagé (99,9 \%).
    \item \textbf{Argument reproductif (le plus fort) :} Fécondité des croisements entre toutes les populations + fécondité des métis sur plusieurs générations = absence d'isolement reproductif.
\end{enumerate}
Concluez par la définition de l'espèce et la conséquence éthique (unité = égalité en dignité et en droits).
\end{methode}

\begin{exercice}[Entraînement : QCM et questions courtes]
\begin{enumerate}
    \item Parmi les propositions suivantes, laquelle n'est \textbf{PAS} un caractère de l'espèce humaine ?
    \begin{itemize}
        \item[A.] Station debout bipède permanente.
        \item[B.] Présence d'une queue vestigiale (coccyx) chez l'adulte.
        \item[C.] Pilosité corporelle épaisse et uniforme sur tout le corps.
        \item[D.] Caryotype à 46 chromosomes (23 paires).
    \end{itemize}
    \item Un couple formé d'un homme d'origine européenne et d'une femme d'origine africaine a un enfant. Cet enfant, devenu adulte, épouse une personne d'origine asiatique et a des enfants féconds. Quel critère d'espèce cette situation illustre-t-elle ?
    \item Complétez : « La diversité phénotypique observable entre humains (taille, couleur de peau...) est codée dans environ \_\_\_\_ \% du génome. »
    \item Vrai ou Faux ? « L'existence de groupes sanguins différents (A, B, AB, O) prouve qu'il existe plusieurs espèces humaines. » Justifiez.
\end{enumerate}
\end{exercice}

\begin{corrige}[Exercice 1]
\begin{enumerate}
    \item \textbf{C.} La pilosité humaine est \textbf{réduite} (cheveux longs, duvet fin), pas épaisse et uniforme (caractère de nombreux mammifères, pas de l'humain). A, B, D sont des caractères d'espèce.
    \item Le critère de \textbf{fécondité des croisements inter-populations} (et la fécondité des descendants sur plusieurs générations = absence d'isolement postzygotique).
    \item \textbf{0,1 \%} (environ 3 millions de paires de bases sur 3 milliards).
    \item \textbf{Faux.} Les groupes sanguins sont des \textbf{polymorphismes génétiques} (allèles multiples d'un même gène \emph{ABO}) au sein d'une \textbf{même espèce}. Ils ne créent aucune stérilité entre porteurs de groupes différents (transfusion possible avec compatibilité, reproduction toujours féconde).
\end{enumerate}
\end{corrige}

\coursec{Les différences entre individus : inventaire et mesure}

Dans la section précédente, nous avons vu que tous les êtres humains partagent des \cle{caractères de l'espèce humaine} : station debout, marche bipède, main préhensile, langage articulé, cerveau développé. Ces ressemblances font de nous une seule et même espèce. Pourtant, il suffit de regarder autour de soi, dans la cour du collège ou au marché, pour constater une évidence : \textbf{aucun individu ne ressemble exactement à un autre}. Même les vrais jumeaux présentent de petites différences. Cette section a pour but de dresser l'\cle{inventaire} de ces différences, de les \cle{mesurer} et de les classer, afin de dégager la grande question qui guidera la suite de la leçon : d'où viennent-elles ?

\courssub{Une enquête statistique en classe : observer pour inventorier}

Pour étudier scientifiquement les différences entre individus, la démarche la plus simple consiste à mener une \cle{enquête} sur un groupe bien défini : notre propre classe. On choisit quelques caractères faciles à observer ou à mesurer, et on consigne les résultats dans un tableau.

\begin{experience}[Enquête sur la diversité des élèves d'une classe de 3ème]
\textbf{Matériel :} une toise (ou un mètre ruban fixé au mur), une balance, un tableau de relevés, et éventuellement les carnets de santé ou les cartes de groupe sanguin.\par
\textbf{Protocole :}
\begin{enumerate}
\item Mesurer la \textbf{taille} de chaque élève, debout, pieds nus, dos contre la toise. Noter la valeur au centimètre près.
\item Peser chaque élève et noter sa \textbf{masse} en kilogrammes.
\item Relever pour chacun : le \textbf{groupe sanguin} (A, B, AB ou O), la \textbf{couleur des yeux} (noirs, marron, clairs), la \textbf{forme du lobe de l'oreille} (collé ou détaché), la \textbf{couleur de la peau}.
\item Regrouper les résultats dans un tableau et compter les effectifs de chaque modalité.
\end{enumerate}
\textbf{Observations :} les tailles et les masses prennent de nombreuses valeurs différentes, toutes proches les unes des autres ; en revanche, le groupe sanguin ou le lobe de l'oreille ne présentent que quelques possibilités nettement distinctes, sans forme intermédiaire.\par
\textbf{Interprétation :} tous les caractères ne varient pas de la même manière. Certains varient « par glissement » continu, d'autres « par catégories » bien séparées.
\end{experience}

Cette enquête illustre une règle fondamentale de la démarche scientifique : avant d'expliquer un phénomène, il faut d'abord le \textbf{décrire} et le \textbf{mesurer} avec soin.

\courssub{Exploiter les mesures : l'histogramme des tailles}

Un tableau de 30 ou 40 mesures brutes est difficile à lire. Pour y voir clair, on regroupe les valeurs par \cle{classes} (par exemple les tailles de 1,40 m à 1,45 m, puis de 1,45 m à 1,50 m, etc.) et on compte combien d'élèves appartiennent à chaque classe. On construit alors un \cle{histogramme} : en abscisse, les classes de tailles ; en ordonnée, les effectifs.

\begin{popfigure}
\begin{center}
\begin{tikzpicture}
\begin{axis}[
    width=0.9\linewidth, height=6.5cm,
    ybar, bar width=14pt,
    xlabel={Taille (m)}, ylabel={Nombre d'élèves},
    symbolic x coords={1.40--1.45, 1.45--1.50, 1.50--1.55, 1.55--1.60, 1.60--1.65, 1.65--1.70, 1.70--1.75, 1.75--1.80},
    xtick=data, x tick label style={rotate=35, anchor=east, font=\small},
    ymin=0, ymax=10,
    nodes near coords, every node near coord/.append style={font=\small},
    axis lines=left, enlarge x limits=0.08,
    ytick={0,2,4,6,8,10},
]
\addplot[fill=popblue!60, draw=popblue] coordinates {
    (1.40--1.45,1) (1.45--1.50,2) (1.50--1.55,4) (1.55--1.60,7)
    (1.60--1.65,8) (1.65--1.70,5) (1.70--1.75,2) (1.75--1.80,1)
};
\end{axis}
\end{tikzpicture}
\end{center}
\legende{Histogramme des tailles d'un échantillon de 30 élèves de 3ème. Les tailles se répartissent en cloche autour d'une valeur moyenne voisine de 1,60 m : les tailles moyennes sont les plus fréquentes, les tailles très petites ou très grandes sont rares.}
\end{popfigure}

L'histogramme dessine une \textbf{courbe en cloche} : la plupart des élèves ont une taille proche de la valeur moyenne, et plus on s'éloigne de cette moyenne (vers les très petites ou les très grandes tailles), plus les effectifs diminuent. C'est la signature d'un caractère à \cle{variation continue}.

\begin{methode}[Exploiter un histogramme]
\begin{enumerate}
\item Repérer la \textbf{valeur (ou la classe) la plus fréquente} : c'est le sommet de la cloche, ici la classe de 1,60 m à 1,65 m avec 8 élèves.
\item Déterminer l'\textbf{étendue} des valeurs : différence entre la plus grande et la plus petite valeur, ici de 1,40 m à 1,80 m, soit une étendue de 0,40 m.
\item Observer la \textbf{forme} de la répartition : en cloche (concentration autour de la moyenne) ou non.
\item \textbf{Conclure} sur la variabilité du caractère dans le groupe : ici, la taille est un caractère très variable, qui prend de nombreuses valeurs intermédiaires autour d'une moyenne d'environ 1,60 m.
\end{enumerate}
\end{methode}

\begin{exemplebox}[Un exemple chiffré]
Dans une classe de 40 élèves de 3ème, les tailles peuvent s'échelonner de \SI{1,45}{m} à \SI{1,80}{m}, avec une moyenne autour de \SI{1,60}{m}. Aucun élève ne mesure exactement 1,60 m au centimètre près, mais la majorité se situe entre 1,55 m et 1,70 m. La moyenne est donc une valeur de référence, pas une valeur obligatoire : chaque individu s'en écarte un peu.
\end{exemplebox}

\courssub{Deux grands types de caractères : variation continue et variation discontinue}

L'enquête en classe fait apparaître deux comportements très différents selon les caractères étudiés. D'où la définition suivante.

\begin{definition}[Caractères à variation continue et à variation discontinue]
Un caractère à \cle{variation continue} peut prendre \textbf{toutes les valeurs intermédiaires} entre deux extrêmes : on ne peut pas ranger les individus dans des catégories nettes. Exemples : la \textbf{taille}, la \textbf{masse}, le \textbf{teint de la peau} (du plus clair au plus foncé, avec toutes les nuances).\par
Un caractère à \cle{variation discontinue} ne présente que \textbf{quelques modalités nettes}, sans forme intermédiaire : chaque individu appartient à une catégorie précise. Exemples : le \textbf{groupe sanguin} (A, B, AB ou O), la forme du \textbf{lobe de l'oreille} (collé ou détaché).
\end{definition}

\begin{center}
\begin{tabularx}{\linewidth}{|l|X|X|}
\hline
\rowcolor{popblueL} \textbf{Critère} & \textbf{Variation continue} & \textbf{Variation discontinue} \\
\hline
Nombre de valeurs & Infinité de valeurs intermédiaires & Quelques modalités nettes (2 à 4) \\
\hline
Exemples & Taille, masse, teint de la peau & Groupe sanguin ABO, lobe de l'oreille \\
\hline
Représentation & Histogramme en cloche & Diagramme en bâtons (catégories séparées) \\
\hline
\end{tabularx}
\end{center}

\courssub{Les empreintes digitales : une signature unique pour chaque individu}

Observe le bout de tes doigts : la peau y présente de fines crêtes dessinant des motifs (boucles, arcs, tourbillons). Ce sont les \cle{empreintes digitales}. Or, même entre vrais jumeaux, ces dessins ne sont jamais identiques : \textbf{chaque individu possède des empreintes digitales qui lui sont propres}. C'est l'exemple le plus frappant de la diversité individuelle au sein de l'espèce humaine.

Cette unicité a une application très concrète : l'\textbf{identification} des personnes. Au Cameroun, les empreintes digitales sont relevées lors de l'établissement de la \textbf{carte nationale d'identité} et du \textbf{passeport} biométrique. Elles permettent de vérifier qu'une personne est bien celle qu'elle prétend être, et sont aussi utilisées par la police scientifique pour identifier les auteurs d'infractions.

\courssub{Les groupes sanguins ABO : une différence aux conséquences vitales}

Le \cle{groupe sanguin} est un caractère à variation discontinue : tout individu appartient à l'un des quatre groupes \textbf{A}, \textbf{B}, \textbf{AB} ou \textbf{O}, sans intermédiaire possible. La répartition de ces groupes n'est pas la même partout : elle varie d'une population à l'autre.

\begin{popfigure}
\begin{center}
\begin{tabularx}{0.85\linewidth}{|l|X|X|}
\hline
\rowcolor{popblueL} \textbf{Groupe sanguin} & \textbf{Effectif (sur 500 personnes)} & \textbf{Fréquence} \\
\hline
O & 255 & 51 \% \\
\hline
A & 120 & 24 \% \\
\hline
B & 100 & 20 \% \\
\hline
AB & 25 & 5 \% \\
\hline
\textbf{Total} & \textbf{500} & \textbf{100 \%} \\
\hline
\end{tabularx}
\end{center}
\legende{Répartition des groupes sanguins ABO dans un échantillon de 500 personnes d'une population camerounaise (valeurs indicatives). Le groupe O y est le plus fréquent, le groupe AB le plus rare.}
\end{popfigure}

Ce caractère a une importance pratique capitale : lors d'une \textbf{transfusion sanguine} (par exemple après un accident ou pendant une opération chirurgicale à l'hôpital), le sang donné doit être \textbf{compatible} avec celui du malade, sous peine d'accident grave. C'est pourquoi les centres de santé déterminent toujours le groupe sanguin avant toute transfusion, et pourquoi les donneurs du groupe O, dont le sang peut être transfusé à tous les autres groupes, sont si précieux.

\begin{savaistu}[Le groupe O, champion d'Afrique]
Le groupe sanguin \textbf{O} est le plus fréquent dans la plupart des populations africaines, y compris au Cameroun. Mais sa fréquence exacte varie d'une région à l'autre et d'une population à l'autre : il n'existe pas une répartition unique valable pour tout le continent. Cette diversité des fréquences est elle-même un signe de la variabilité humaine.
\end{savaistu}

\courssub{La question-problème : d'où viennent ces différences ?}

Notre inventaire est terminé : les individus de l'espèce humaine diffèrent par de nombreux caractères, à variation continue (taille, masse, teint) ou discontinue (groupe sanguin, lobe de l'oreille), et certains caractères comme les empreintes digitales sont uniques à chacun. Reste la grande question, celle qui motive toute la suite de la leçon :

\begin{aretenir}[Question-problème]
D'où viennent les différences observées entre les individus ? Sont-elles \textbf{transmises par les parents} (caractères héréditaires) ou \textbf{acquises au cours de la vie} sous l'effet de l'environnement (alimentation, climat, mode de vie) ? Les sections suivantes apporteront la réponse, caractère par caractère.
\end{aretenir}

\begin{attention}[Erreur fréquente à éviter]
On ne peut \textbf{pas} conclure sur l'origine d'un caractère (héréditaire ou acquis) à partir d'une \textbf{seule observation}. Dire « cet élève est grand parce que son père est grand » n'est qu'une hypothèse. Pour conclure rigoureusement, il faut des \textbf{comparaisons} (enfants et parents, jumeaux, populations vivant dans des milieux différents) et des \textbf{données chiffrées} nombreuses. Une observation isolée ne prouve rien.
\end{attention}

\begin{exoresolu}[Exploiter les résultats d'une enquête de classe]
\textbf{Énoncé.} Dans une classe de 3ème de 30 élèves, on a relevé : 15 élèves de groupe O, 8 de groupe A, 5 de groupe B et 2 de groupe AB ; par ailleurs, les tailles vont de 1,42 m à 1,78 m. 1) Calculer la fréquence de chaque groupe sanguin. 2) Indiquer, en justifiant, si le groupe sanguin et la taille sont des caractères à variation continue ou discontinue. 3) Un élève affirme : « Mon groupe sanguin vient forcément de ma mère. » Cette affirmation est-elle acceptable scientifiquement ?
\tcblower
\textbf{Solution détaillée.}\par
1) Les fréquences se calculent en divisant chaque effectif par l'effectif total (30) :
\begin{align*}
f(O) &= \frac{15}{30} = 0,50 \text{ soit } 50\,\% &
f(A) &= \frac{8}{30} \approx 0,267 \text{ soit } 26,7\,\% \\
f(B) &= \frac{5}{30} \approx 0,167 \text{ soit } 16,7\,\% &
f(AB) &= \frac{2}{30} \approx 0,067 \text{ soit } 6,7\,\%
\end{align*}
On vérifie que la somme des fréquences vaut bien 100 \%.\par
2) Le \textbf{groupe sanguin} est un caractère à \textbf{variation discontinue} : il n'existe que quatre modalités nettes (A, B, AB, O), sans forme intermédiaire. La \textbf{taille} est un caractère à \textbf{variation continue} : entre 1,42 m et 1,78 m, toutes les valeurs intermédiaires sont possibles.\par
3) Non, cette affirmation n'est pas acceptable : elle repose sur une seule observation et sur une supposition. Pour déterminer l'origine d'un caractère, il faut comparer les caractères de \textbf{plusieurs} enfants avec ceux de leurs \textbf{deux} parents et recueillir des données nombreuses. Une conclusion scientifique exige des comparaisons et des preuves, pas une impression.
\end{exoresolu}

\begin{exercice}[À toi de jouer]
On mesure les tailles (en m) de 10 élèves : 1,50 ; 1,62 ; 1,55 ; 1,70 ; 1,48 ; 1,60 ; 1,58 ; 1,65 ; 1,52 ; 1,75. 1) Calculer la taille moyenne de ce groupe. 2) Déterminer l'étendue des valeurs. 3) Dire si ce caractère est à variation continue ou discontinue, et justifier. 4) Expliquer pourquoi on ne peut pas affirmer, à partir de ces seules mesures, que la taille est un caractère héréditaire.
\end{exercice}

En résumé, l'inventaire et la mesure des différences entre individus révèlent une \textbf{diversité considérable} au sein de l'espèce humaine : caractères à variation continue dessinant des courbes en cloche, caractères à variation discontinue répartis en catégories nettes, et caractères uniques comme les empreintes digitales. Cette diversité, correctement décrite et mesurée, pose désormais la question de son \textbf{origine} : c'est l'objet des sections suivantes, consacrées aux caractères héréditaires puis aux caractères modifiés par l'environnement.

\coursec{Les caractères héréditaires}

\courssub{De l'observation familiale à l'hypothèse de transmission}

Dans la section précédente, nous avons inventorié les différences entre individus (taille, couleur de peau, groupe sanguin, etc.) et distingué les caractères mesurables des caractères qualitatifs. Une question s'impose naturellement : \textbf{d'où proviennent ces différences ?} Pourquoi un enfant a-t-il les yeux marron comme sa mère et le groupe sanguin B comme son père, alors que son grand frère a les yeux bleus et le groupe A ?

Observons une situation familière. Dans de nombreuses familles camerounaises, comme ailleurs, on remarque des \textbf{ressemblances frappantes} entre parents et enfants : le nez du père, le sourire de la mère, la forme du lobe de l'oreille, ou encore le groupe sanguin. Ces ressemblances ne sont pas dues au hasard ni au mode de vie (manger les mêmes aliments, vivre sous le même toit). Elles suggèrent qu'une \textbf{information} a été transmise de la génération précédente à la suivante.

\begin{exemplebox}[Hypothèse de travail]
Certains caractères se transmettent des parents aux enfants par l'intermédiaire des gamètes lors de la fécondation. On dit que ces caractères sont \textbf{héréditaires}.
\end{exemplebox}

Pour valider cette hypothèse, les généticiens étudient des \textbf{arbres généalogiques} sur plusieurs générations en suivant un caractère précis, facile à identifier et ne changeant pas au cours de la vie : le \textbf{groupe sanguin} (système ABO).

\courssub{Étude documentaire : la transmission du groupe sanguin dans une famille}

Considérons une famille réelle observée sur deux générations. Le père est de groupe sanguin \textbf{A}, la mère de groupe \textbf{B}. Ils ont quatre enfants : deux sont de groupe \textbf{A}, un de groupe \textbf{B} et un de groupe \textbf{AB}. Aucun enfant n'est de groupe \textbf{O}.

\begin{experience}[Analyse d'un arbre généalogique pour le groupe sanguin]
\textbf{Document :} Arbre généalogique d'une famille (2 générations) pour le système ABO.
\textbf{Question :} La répartition des groupes sanguins chez les enfants est-elle compatible avec une transmission héréditaire depuis les parents ?

\end{experience}
\begin{popfigure}
\begin{tikzpicture}[
    scale=0.85,
    every node/.style={transform shape},
    male/.style={draw, thick, rectangle, rounded corners=4pt, minimum width=2.2cm, minimum height=1cm, fill=popblue!15, inner sep=4pt},
    female/.style={draw, thick, circle, minimum width=2.2cm, minimum height=1cm, fill=poppink!15, inner sep=4pt},
    child/.style={draw, thick, minimum width=1.8cm, minimum height=0.85cm, inner sep=3pt},
    line/.style={thick, popdark},
    labelstyle/.style={font=\small, text=popdark, anchor=north}
]
% Generation 1 (Parents)
\node[female, label=above:\textbf{Mère (Groupe B)}] (mere) at (0, 3) {};
\node[male, label=above:\textbf{Père (Groupe A)}] (pere) at (5, 3) {};
% Union line
\draw[line] (mere.east) -- (pere.west);
\draw[line] ($(mere.east)!0.5!(pere.west)$) -- ++(0, -0.8) coordinate (junction);
% Generation 2 (Enfants) - 4 enfants
% Enfant 1 (Fille, Groupe A)
\node[child, rectangle, rounded corners=3pt, fill=popblue!10, label=below:Fille 1 (A)] (e1) at (-3.5, 0.5) {};
% Enfant 2 (Garçon, Groupe B)
\node[child, circle, fill=poppink!10, label=below:Garçon 1 (B)] (e2) at (-0.5, 0.5) {};
% Enfant 3 (Fille, Groupe AB)
\node[child, rectangle, rounded corners=3pt, fill=poppurple!10, label=below:Fille 2 (AB)] (e3) at (2.5, 0.5) {};
% Enfant 4 (Garçon, Groupe A)
\node[child, circle, fill=popblue!10, label=below:Garçon 2 (A)] (e4) at (5.5, 0.5) {};
% Lines from junction to children
\foreach \child in {e1, e2, e3, e4} {
    \draw[line] (junction) |- (\child.north);
}
% Legend
\node[labelstyle, align=left] at (8.5, 3) {\textbf{Légende :}\\[2pt]
$\square$ = Homme / Garçon\\
$\bigcirc$ = Femme / Fille\\
\textcolor{popblue}{Bleu} = Groupe A\\
\textcolor{poppink}{Rose} = Groupe B\\
\textcolor{poppurple}{Violet} = Groupe AB\\
Trait horizontal = Union\\
Trait vertical = Descendance};
\end{tikzpicture}
\legende{Arbre généalogique simplifié sur deux générations montrant la répartition des groupes sanguins ABO. Les enfants présentent des groupes sanguins compatibles avec ceux des parents.}
\end{popfigure}

\textbf{Interprétation :}
\begin{itemize}
    \item Les enfants ne possèdent \textbf{que des groupes sanguins compatibles} avec ceux des parents (A, B, AB). Aucun groupe "inventé" (ex. : groupe C inexistant) n'apparaît.
    \item Le caractère "groupe sanguin" est présent \textbf{dès la naissance} et ne change pas au cours de la vie, quel que soit l'alimentation ou l'environnement.
    \item La régularité de cette transmission, observable sur de nombreuses familles, n'est pas due au hasard.
\end{itemize}

\begin{aretenir}
La régularité de la transmission du groupe sanguin, observable sur les arbres généalogiques, confirme l'hypothèse : \textbf{le groupe sanguin est un caractère héréditaire}. Il en va de même pour la couleur des yeux, la forme du lobe de l'oreille (attaché ou libre), et certaines maladies comme la drépanocytose.
\end{aretenir}

\courssub{Le support physique de l'hérédité : chromosomes et gènes}

Comment l'information "groupe A" ou "lobe attaché" voyage-t-elle du parent à l'enfant ? La réponse se trouve au cœur de nos cellules.

\begin{definition}[Support de l'hérédité]
L'information héréditaire est portée par les \textbf{gènes}, qui sont des portions d'\textbf{ADN} situées sur les \textbf{chromosomes} logés dans le \textbf{noyau} de chaque cellule.
\end{definition}

\begin{propriete}[Transmission chromosomique - Notion admise en 3ème]
\begin{itemize}
    \item La cellule humaine (sauf gamètes) possède \textbf{46 chromosomes} organisés en \textbf{23 paires} (caryotype humain). C'est le nombre \textbf{diploïde} ($2n = 46$).
    \item Les \textbf{gamètes} (spermatozoïde chez l'homme, ovule chez la femme) ne possèdent que \textbf{23 chromosomes} (un seul chromosome par paire). C'est le nombre \textbf{haploïde} ($n = 23$). Ils résultent de la méiose (étudiée en classe supérieure, admis ici).
    \item Lors de la \textbf{fécondation}, la fusion du spermatozoïde (23 chromosomes) et de l'ovule (23 chromosomes) reconstitue le nombre diploïde : \textbf{23 + 23 = 46 chromosomes} dans l'œuf (zygote).
    \item L'enfant reçoit donc \textbf{la moitié de ses chromosomes de son père} (via le spermatozoïde) et \textbf{l'autre moitié de sa mère} (via l'ovule).
\end{itemize}
\end{propriete}

\begin{popfigure}
\begin{tikzpicture}[scale=1.1, every node/.style={font=\small}]
% Styles
\tikzset{
    chrom/.style={draw=popdark, very thick, rounded corners=2pt, minimum height=1.2cm, inner sep=2pt},
    chromM/.style={chrom, fill=popblue!30, minimum width=0.6cm},
    chromF/.style={chrom, fill=poppink!30, minimum width=0.6cm},
    chromZ/.style={chrom, minimum width=0.6cm, fill=poppurple!20, pattern=north east lines, pattern color=poppurple!50},
    arrow/.style={->, very thick, popdark, shorten >=2pt, shorten <=2pt},
    labelstyle/.style={font=\small, text=popdark, anchor=north},
    labelbox/.style={rectangle, draw=popdark, fill=popcream, rounded corners=4pt, inner sep=4pt, font=\small\bfseries, align=center}
}
% PERE (Gamete male)
\node[labelbox, align=center] (pere_label) at (-5, 4){PÈRE\\ (Cellules somatiques : $2n=46$)};
\node[chromM, label={below:$n=23$}] (sperm) at (-5, 2.5) {SPERMATOZOÏDE};
\draw[arrow] (-5, 3.3) -- (-5, 2.7);
% MERE (Gamete female)
\node[labelbox, align=center] (mere_label) at (5, 4){MÈRE\\ (Cellules somatiques : $2n=46$)};
\node[chromF, label={below:$n=23$}] (ovule) at (5, 2.5) {OVULE};
\draw[arrow] (5, 3.3) -- (5, 2.7);
% FECONDATION
\node[labelbox, fill=popgold!50, align=center] (fecond) at (0, 1.2){FÉCONDATION\\ (Fusion des noyaux)};
\draw[arrow] (-4.4, 2.5) -- (-1.5, 1.5);
\draw[arrow] (4.4, 2.5) -- (1.5, 1.5);
% ZYGOTE (Enfant)
\node[labelbox, align=center] (zygote_label) at (0, -0.5){ŒUF (ZYGOTE) \\ puis ENFANT\\ ($2n=46$)};
\node[chromZ, label={below:$2n=46$}] (zygote) at (0, -1.8) {Noyau de l'œuf : 23 paires};
\draw[arrow] (0, 0.5) -- (0, -0.1);
% Chromosome pairs representation in zygote (simplified)
\foreach \i/\x in {1/-1.2, 2/-0.6, 3/0, 4/0.6, 5/1.2} {
    \draw[chromM, minimum width=0.3cm, minimum height=0.6cm] (\x-0.1, -2.2) -- ++(0, -0.6);
    \draw[chromF, minimum width=0.3cm, minimum height=0.6cm] (\x+0.1, -2.2) -- ++(0, -0.6);
    \draw[popdark, thin] (\x-0.1, -2.5) -- (\x+0.1, -2.5); % centromere link visual
}
% Legend
\node[labelstyle, align=left, anchor=north west] at (-6, -2.5) {
\textbf{Légende :}\\
\textcolor{popblue!70!black}{\rule{3mm}{3mm}} Chromosome d'origine paternelle\\
\textcolor{poppink!70!black}{\rule{3mm}{3mm}} Chromosome d'origine maternelle\\
\textcolor{poppurple!70!black}{\rule{3mm}{3mm}} Paire homologue (1 paternel + 1 maternel)\\
\textbf{Gènes :} portions d'ADN sur les chromosomes portant l'info du caractère.
};
\end{tikzpicture}
\legende{Schéma de la transmission des chromosomes lors de la fécondation humaine. Le zygote reçoit 23 chromosomes du père (bleu) et 23 de la mère (rose), restaurant le caryotype diploïde à 46 chromosomes (23 paires). Les gènes, portés par ces chromosomes, déterminent les caractères héréditaires.}
\end{popfigure}

\courssub{Cas particuliers : les vrais et les faux jumeaux}

L'étude des jumeaux apporte une preuve éclatante du rôle du patrimoine chromosomique (génotype) dans la détermination des caractères.

\begin{exemplebox}[Vrais jumeaux (Jumeaux monozygotes)]
Ils proviennent de \textbf{un seul ovule fécondé par un seul spermatozoïde} (un seul zygote). Ce zygote se divise en deux masses cellulaires distinctes très tôt (vers le 3\textsuperscript{e}-4\textsuperscript{e} jour).
\begin{itemize}
    \item Ils possèdent \textbf{exactement le même patrimoine héréditaire} (mêmes 46 chromosomes, mêmes gènes, même ADN).
    \item Ils sont \textbf{toujours du même sexe}.
    \item Ils se ressemblent \textbf{fortement} (même groupe sanguin, même couleur d'yeux, mêmes empreintes digitales, mêmes prédispositions génétiques aux maladies).
\end{itemize}
\textit{Différences minimes possibles :} de légères variations peuvent apparaître à cause de l'environnement (nutrition, maladies, mode de vie), mais le "plan de construction" génétique est identique.
\end{exemplebox}

\begin{exemplebox}[Faux jumeaux (Jumeaux dizygotes)]
Ils proviennent de \textbf{deux ovules différents fécondés par deux spermatozoïdes différents} (deux zygotes distincts, deux fécondations simultanées).
\begin{itemize}
    \item Ils partagent en moyenne \textbf{50\% de leur patrimoine héréditaire}, exactement comme des frères et sœurs nés à des dates différentes.
    \item Ils peuvent être \textbf{de sexes différents} (un garçon, une fille).
    \item Ils ne se ressemblent \textbf{pas plus} que des frères et sœurs ordinaires.
\end{itemize}
\end{exemplebox}

\begin{aretenir}[Preuve par les jumeaux]
La concordance quasi-totale des caractères héréditaires chez les vrais jumeaux (même patrimoine génétique) et la discordance fréquente chez les faux jumeaux (patrimoines génétiques différents) démontrent que \textbf{la ressemblance familiale repose sur la communauté de gènes}, et non sur le seul fait de grandir ensemble.
\end{aretenir}

\courssub{Exemples de caractères héréditaires et distinction cruciale}

Voici une liste non exhaustive de caractères dont l'hérédité est bien établie et observable dans nos populations :

\begin{center}
\begin{tabularx}{\linewidth}{|l|X|X|}\hline
\textbf{Caractère} & \textbf{Type} & \textbf{Observation courante} \\ \hline
Groupe sanguin (ABO, Rhésus) & Qualitatif & Déterminé à la naissance, invariant, crucial pour transfusions (hôpitaux de Yaoundé, Douala, etc.) \\ \hline
Couleur des yeux / cheveux & Qualitatif & Variété phénotypique visible (marron, noir, bleu, vert) \\ \hline
Forme du lobe de l'oreille & Qualitatif & Attaché (collé à la joue) vs Libre (pendant) — caractère classique d'initiation génétique \\ \hline
Drépanocytose (anémie falciforme) & Maladie génétique & \textbf{Très fréquente au Cameroun} (zone endémique palustre). Transmission héréditaire : parents porteurs sains peuvent avoir un enfant malade. \\ \hline
Daltonisme & Anomalie vision & Difficulté à distinguer rouge/vert. Lié aux chromosomes sexuels (plus fréquent chez les garçons). \\ \hline
\end{tabularx}
\end{center}

\begin{attention}[Piège majeur : Caractère héréditaire vs Caractère acquis]
\textbf{Erreur fréquente au BEPC :} Dire qu'un enfant « a hérité des gros muscles de son père parce qu'il fait de la musculation » ou « a hérité de la cicatrice de sa mère ».
\begin{itemize}
    \item \textbf{Caractère héréditaire :} Porté par les gènes (ADN), présent dans les gamètes, transmis à la fécondation. Ex : groupe sanguin, couleur des yeux, drépanocytose.
    \item \textbf{Caractère acquis (modifié par l'environnement) :} Apparaît au cours de la vie sous l'effet du milieu (nutrition, exercice, maladie, apprentissage, accident). \textbf{Il ne modifie pas l'ADN des gamètes et NE SE TRANSMET PAS à la descendance.}
\end{itemize}
\textbf{Astuce de rédaction :} Pour prouver qu'un caractère est héréditaire, on vérifie : (1) Il existe chez les parents/ancêtres. (2) Il apparaît dès la naissance (ou se manifeste selon un programme génétique précis, comme la calvitie). (3) Il ne dépend pas du mode de vie (sport, régime, éducation).
\end{attention}

\courssub{Méthode : Démontrer qu'un caractère est héréditaire}

\begin{methode}[Démarche d'investigation pour identifier un caractère héréditaire]
\textbf{Étape 1 : Observer} une famille sur au moins deux générations (parents $\to$ enfants). Noter la présence/absence du caractère.
\textbf{Étape 2 : Comparer} la fréquence du caractère chez les enfants. Si le caractère réapparaît systématiquement d'une génération à l'autre, c'est un indice fort.
\textbf{Étape 3 : Vérifier l'antériorité} : le caractère est-il visible dès la naissance (ou à un âge programmé) sans apprentissage ni condition de vie particulière ?
\textbf{Étape 4 : Exclure l'environnement} : des individus au patrimoine génétique identique (vrais jumeaux) élevés dans des milieux différents conservent-ils le caractère ? Si oui $\to$ héréditaire. Si le caractère change selon le milieu $\to$ acquis (ou interaction avec l'environnement).
\textbf{Étape 5 : Conclure} : « Le caractère X est héréditaire car il se transmet de génération en génération, est déterminé dès la fécondation et ne dépend pas du mode de vie. »
\end{methode}

\begin{exoresolu}[Analyse de la transmission de la drépanocytose]
\textbf{Énoncé :} Au Cameroun, un couple d'asymptomatiques (porteurs sains, phénotype normal) a quatre enfants. Deux enfants sont porteurs sains, un est drépanocytaire majeur et un est non porteur. Expliquer pourquoi la drépanocytose est un caractère héréditaire et non acquis.

\textbf{Solution détaillée :}
\begin{enumerate}
    \item \textbf{Observation du pedigree :} Les parents sont phénotypiquement sains (pas de crise drépanocytaire), mais ils ont un enfant malade. Cela prouve qu'ils portent l'information génétique de la maladie sans l'exprimer (ils sont porteurs sains). Cette information a été transmise par les gamètes lors de la fécondation.
    \item \textbf{Régularité familiale :} Sur 4 enfants, on observe les trois situations possibles (non porteur, porteur sain, malade). Cette répartition, retrouvée dans de nombreuses familles, exclut le hasard ou une contamination environnementale.
    \item \textbf{Antériorité et invariance :} L'information génétique responsable de la maladie est présente dans l'ADN de l'enfant dès la fécondation (23 chromosomes paternels + 23 maternels). La maladie se déclare vers 6 mois (quand l'hémoglobine fœtale diminue), indépendamment de l'alimentation ou de l'hygiène de vie.
    \item \textbf{Preuve par les jumeaux :} Si l'enfant malade avait un vrai jumeau, il serait obligatoirement malade aussi (même zygote, même patrimoine génétique). Si c'était un faux jumeau, il pourrait être sain, porteur ou malade.
    \item \textbf{Conclusion :} La drépanocytose est un \textbf{caractère héréditaire} (maladie génétique). L'information responsable est transmise par les gamètes. Elle n'est pas attrapée comme un paludisme (maladie infectieuse) ni causée par la nourriture.
\end{enumerate}
\end{exoresolu}

\courssub{Bilan de la section}

\begin{aretenir}[Synthèse : Les caractères héréditaires]
\begin{itemize}
    \item Un \textbf{caractère héréditaire} est un caractère transmis des parents aux enfants via les \textbf{gamètes} (spermatozoïde, ovule) lors de la \textbf{fécondation}.
    \item Son support physique est le \textbf{gène}, portion d'ADN portée par les \textbf{chromosomes} du noyau.
    \item L'enfant reçoit \textbf{23 chromosomes du père + 23 chromosomes de la mère = 46 chromosomes (23 paires)}.
    \item Exemples : Groupe sanguin, couleur des yeux, lobe d'oreille, drépanocytose, daltonisme.
    \item \textbf{Vrais jumeaux (monozygotes) :} Même patrimoine génétique $\to$ forte ressemblance.
    \item \textbf{Faux jumeaux (dizygotes) :} Patrimoines génétiques différents (50\% communs) $\to$ ressemblance fraternelle ordinaire.
    \item \textbf{Distinction fondamentale :} Héréditaire (génétique, transmissible) $\neq$ Acquis (environnemental, non transmissible à la descendance).
\end{itemize}
\end{aretenir}

\begin{savaistu}[La drépanocytose au Cameroun : un enjeu de santé publique]
La drépanocytose est la \textbf{maladie génétique la plus fréquente au Cameroun} (et en Afrique centrale/ouest). On estime qu'environ \textbf{2 à 3\% des nouveau-nés} sont drépanocytaires majeurs et \textbf{20 à 25\% de la population} sont porteurs sains.
\begin{itemize}
    \item \textbf{Contexte local :} Le Cameroun est une zone à forte transmission du paludisme. La présence du gène de la drépanocytose dans la population est liée à l'histoire de cette maladie (résistance relative au paludisme grave pour les porteurs sains), mais cet aspect évolutif sera approfondi au lycée.
    \item \textbf{Dépistage néonatal :} Le Cameroun déploie le dépistage systématique à la naissance (test sur sang de cordon ou talon) pour une prise en charge précoce (pénicilline prophylactique, vaccins, acide folique, éducation des parents dans les centres spécialisés comme le Centre de la Drépanocytose de l'Hôpital Général de Yaoundé).
    \item \textbf{Conseil génétique :} Un couple de porteurs sains a un risque à chaque grossesse d'avoir un enfant malade. L'information génétique permet aux familles de se préparer et de bénéficier d'un suivi médical adapté.
\end{itemize}
Cet exemple montre qu'un caractère \textbf{invisible à l'œil nu} (molécule d'hémoglobine modifiée) est un caractère héréditaire majeur, dont la compréhension sauve des vies.
\end{savaistu}

\begin{exercice}[Entraînement BEPC - Caractères héréditaires vs Acquis]
Classer les caractères suivants en deux colonnes : \textbf{Héréditaires (H)} et \textbf{Acquis (A)}. Justifier brièvement chaque classement.
\begin{enumerate}
    \item Groupe sanguin A Rhésus +.
    \item Cicatrice sur le front suite à une chute à vélo.
    \item Muscles hypertrophiés obtenus par 3 ans de musculation intensive.
    \item Couleur des yeux marron foncé.
    \item Drépanocytose (enfant malade né de parents sains).
    \item Parler couramment le français et l'ewondo.
    \item Forme du lobe de l'oreille libre.
    \item Bronzage intense après un séjour à Kribi.
\end{enumerate}
\end{exercice}

\begin{corrige}[Exercice : Héréditaire vs Acquis]
\begin{center}
\begin{tabularx}{\linewidth}{|l|l|X|}\hline
\textbf{Caractère} & \textbf{Classe} & \textbf{Justification} \\ \hline
Groupe sanguin A Rh+ & H & Déterminé par des gènes, présent dès la conception, invariant toute la vie. \\ \hline
Cicatrice & A & Résultat d'un traumatisme (environnement), ne modifie pas l'ADN des gamètes. \\ \hline
Muscles hypertrophiés & A & Adaptation physiologique à l'effort (entraînement), non codée dans le génome. \\ \hline
Yeux marron & H & Pigmentation de l'iris sous contrôle génétique, stable dès la naissance. \\ \hline
Drépanocytose & H & Maladie génétique transmise par les gamètes des parents porteurs sains. \\ \hline
Langues parlées & A & Apprentissage culturel, aucun gène ne code pour une langue précise. \\ \hline
Lobe libre & H & Caractère morphologique transmis par les gènes, visible dès la naissance. \\ \hline
Bronzage & A & Réaction physiologique (production de mélanine) aux UV, réversible, non transmis. \\ \hline
\end{tabularx}
\end{center}
\textbf{Règle mnémotechnique :} « Ce qui est dans le \textbf{noyau} (ADN) voyage dans les \textbf{gamètes} ; ce qui est dans le \textbf{corps} (muscle, cicatrice, mémoire) reste dans le corps. »
\end{corrige}

\coursec{Les caractères modifiés par l'environnement}

\courssub{Transition : de l'hérédité à l'influence du milieu}

Dans la section précédente, nous avons vu que les \textbf{caractères héréditaires} (groupe sanguin, couleur de base de l'œil, forme du lobe de l'oreille) sont portés par les gènes présents sur les chromosomes, et qu'ils se transmettent de génération en génération selon des lois précises. Mais si tu observes tes camarades de classe, tu constates que deux vrais jumeaux, qui possèdent \emph{exactement} le même patrimoine génétique, ne sont pas toujours identiques en tout point : l'un peut être plus massif, avoir la peau plus foncée à cause du soleil, porter une cicatrice ou parler une langue que l'autre ne connaît pas. D'où viennent ces différences ? C'est l'objet de cette section : comprendre comment l'\textbf{environnement} sculpte le phénotype, sans modifier le génotype.

\begin{definition}[Environnement d'un individu]
L'\cle{environnement} est l'ensemble des conditions de vie qui entourent un individu tout au long de son existence : \textbf{alimentation} (qualité, quantité), \textbf{climat} (température, ensoleillement), \textbf{activités physiques} (sport, travail manuel), \textbf{maladies} subies, \textbf{éducation} (scolarité, langue parlée, pratiques culturelles), \textbf{hygiène}, etc. C'est le « milieu extérieur » par opposition au « milieu intérieur » (génome).
\end{definition}

\begin{definition}[Caractère modifié par l'environnement (caractère acquis)]
Un \cle{caractère modifié par l'environnement} (ou \cle{caractère acquis}) est un caractère qui \textbf{apparaît au cours de la vie} sous l'influence des conditions de vie (environnement). Il \textbf{n'est pas présent dans le génome} et, par conséquent, \textbf{il ne se transmet pas aux descendants} (non héréditaire).
\end{definition}

\courssub{Étude de cas 1 — La taille : un caractère héréditaire sensible à la nutrition}

La taille est un exemple classique de caractère \textbf{à déterminisme mixte} : elle a une base génétique forte (les enfants de parents grands tendent à être grands), mais son expression finale dépend fortement de l'alimentation et de l'état de santé pendant la croissance.

\begin{experience}[Comparaison de la croissance staturo-pondérale selon l'état nutritionnel]
\textbf{Question :} Comment l'alimentation influence-t-elle l'expression du caractère « taille » chez des enfants de même potentiel génétique ?
\par\textbf{Document :} Suivi longitudinal de deux groupes d'enfants camerounais de 0 à 18 ans, issus de milieux socio-économiques comparables sur le plan génétique mais différents sur le plan nutritionnel.
\begin{itemize}
    \item \textbf{Groupe A (Bien nourris) :} Alimentation diversifiée, apport calorique et protidique suffisant (\textasciitilde 2\,200 kcal/j à 10 ans), suivi médical régulier.
    \item \textbf{Groupe B (Malnutris chroniques) :} Alimentation à base de féculents pauvres en protéines, apport calorique insuffisant (\textasciitilde 1\,300 kcal/j à 10 ans), infections répétées (paludisme, diarrhées).
\end{itemize}
\par\textbf{Observations (données moyennes) :}
\begin{center}
\begin{tabular}{|c|c|c|}\hline
\textbf{Âge (ans)} & \textbf{Taille moyenne Groupe A (cm)} & \textbf{Taille moyenne Groupe B (cm)} \\ \hline
0 (naissance) & 50 & 49 \\ \hline
5 & 110 & 100 \\ \hline
10 & 138 & 125 \\ \hline
15 & 165 & 148 \\ \hline
18 (adulte) & 172 & 158 \\ \hline
\end{tabular}
\end{center}
\par\textbf{Interprétation :}
\begin{enumerate}
    \item À la naissance, les tailles sont proches (différence 1 cm) : le potentiel génétique de départ est similaire.
    \item L'écart se creuse progressivement pendant la croissance (10 cm à 5 ans, 14 cm à 18 ans).
    \item Le Groupe B n'atteint pas sa \textbf{taille cible génétique} : l'environnement (malnutrition, maladies) a \textbf{limité l'expression} des gènes de la croissance.
\end{enumerate}
\par\textbf{Conclusion :} La taille est un caractère \textbf{héréditaire dans son potentiel}, mais \textbf{modifié par l'environnement} dans sa réalisation. Une amélioration de la nutrition (transition nutritionnelle) fait augmenter la taille moyenne d'une population en une génération (phénomène de \emph{sécularisation de la croissance}), sans modification du patrimoine génétique.
\end{experience}

\begin{popfigure}
\begin{tikzpicture}
\begin{axis}[
    width=\linewidth,
    height=0.6\linewidth,
    xlabel={Âge (ans)},
    ylabel={Taille moyenne (cm)},
    xmin=0, xmax=18,
    ymin=45, ymax=180,
    xtick={0,5,10,15,18},
    ytick={50,100,150,170},
    legend pos=north west,
    grid=major,
    popgrid,
    popcurve/.style={thick, mark=*},
]
\addplot[popblue, popcurve] coordinates {
    (0,50) (5,110) (10,138) (15,165) (18,172)
};
\addlegendentry{Groupe A : Bien nourris}
\addplot[poporange, popcurve] coordinates {
    (0,49) (5,100) (10,125) (15,148) (18,158)
};
\addlegendentry{Groupe B : Malnutris chroniques}
\end{axis}
\end{tikzpicture}
\legende{Évolution de la taille moyenne selon l'état nutritionnel. L'écart entre les deux courbes illustre l'action limitante de la malnutrition sur l'expression du potentiel génétique de croissance.}
\end{popfigure}

\begin{attention}[Piège fréquent : « La taille n'est pas héréditaire car elle change avec la nourriture »]
\textbf{FAUX.} La taille \textbf{est} héréditaire (les gènes fixent la \emph{plage} possible). L'environnement (nutrition) détermine \emph{où} dans cette plage l'individu se situera. Preuve : deux enfants bien nourris de parents petits resteront plus petits que deux enfants bien nourris de parents grands. Ne confonds pas \emph{non-hérédité} et \emph{sensibilité à l'environnement}.
\end{attention}

\courssub{Étude de cas 2 — Les vrais jumeaux élevés séparément : l'expérience naturelle décisive}

Les \textbf{vrais jumeaux (monozygotes)} proviennent de la division d'un même œuf fécondé : ils possèdent \textbf{exactement le même génotype} (même ADN nucléaire). S'ils sont élevés dans des milieux différents, toutes les différences observées entre eux sont \textbf{nécessairement dues à l'environnement}.

\begin{exemplebox}[Cas réel : Jumeaux séparés à la naissance (Étude type Minnesota / Bouchard)]
Deux vrais jumeaux, \textbf{Paul} et \textbf{Pierre}, séparés à la naissance et adoptés par des familles différentes.
\begin{itemize}
    \item \textbf{Paul} : élevé en zone urbaine (Yaoundé), alimentation riche et variée, scolarité continue, pratique le football 3 fois/semaine.
    \item \textbf{Pierre} : élevé en zone rurale enclavée, alimentation à base de manioc/maïs (pauvre en protéines), scolarité interrompue à 12 ans, travail agricole manuel intense, expositions solaires prolongées.
\end{itemize}
\par\textbf{Comparaison à 25 ans :}
\begin{center}
\begin{tabular}{|l|c|c|}\hline
\textbf{Caractère} & \textbf{Paul} & \textbf{Pierre} \\ \hline
Groupe sanguin (ABO) & A+ & A+ \\ \hline
Couleur des yeux & Marron & Marron \\ \hline
Forme du lobe d'oreille & Attaché & Attaché \\ \hline
\textbf{Masse corporelle} & \textbf{78 kg} & \textbf{62 kg} \\ \hline
\textbf{Teint de la peau (bras)} & \textbf{Clair (protégé)} & \textbf{Très foncé (bronzé)} \\ \hline
Langue parlée couramment & Français, Anglais & Langue locale, Pidgin \\ \hline
Cicatrices & Aucune visible & Cicatrice au genou (chute) \\ \hline
Force musculaire (dynamomètre) & 45 kg & 55 kg \\ \hline
\end{tabular}
\end{center}
\par\textbf{Analyse :}
\begin{itemize}
    \item \textbf{Identiques (Hérédité) :} Groupe sanguin, couleur des yeux, lobes d'oreilles $\rightarrow$ Caractères \textbf{héréditaires} (génotype identique).
    \item \textbf{Différents (Environnement) :} Masse (-16 kg), teint (bronzage), langue, cicatrice, force musculaire $\rightarrow$ Caractères \textbf{acquis/modifiés par l'environnement}.
\end{itemize}
\end{exemplebox}

\begin{exoresolu}[Analyse de l'origine des caractères chez des jumeaux]
\textbf{Énoncé :} Léa et Lou sont de vrais jumeaux. À 18 ans, Léa mesure 1,65 m pour 55 kg, elle a la peau mate, parle français et espagnol, et porte une cicatrice d'appendicectomie. Lou mesure 1,65 m pour 72 kg, a la peau très foncée, parle seulement français, et n'a pas de cicatrice abdominale. Ils ont tous deux les yeux noisette et le groupe sanguin O+.
\par\textbf{Question :} Classe chaque caractère dans le tableau ci-dessous en justifiant.
\begin{center}
\begin{tabular}{|l|c|c|}\hline
\textbf{Caractère} & \textbf{Héréditaire} & \textbf{Acquis (Environnement)} \\ \hline
Taille (1,65 m) & & \\ \hline
Masse (55 vs 72 kg) & & \\ \hline
Couleur des yeux (noisette) & & \\ \hline
Teint de la peau (mate vs très foncée) & & \\ \hline
Groupe sanguin (O+) & & \\ \hline
Langues parlées & & \\ \hline
Cicatrice & & \\ \hline
\end{tabular}
\end{center}
\tcblower
\textbf{Solution détaillée :}
\begin{center}
\begin{tabular}{|l|c|c|p{5cm}|}\hline
\textbf{Caractère} & \textbf{Héréditaire} & \textbf{Acquis} & \textbf{Justification} \\ \hline
Taille (1,65 m) & $\checkmark$ & & Identique chez les deux ; potentiel génétique réalisé (nutrition supposée correcte pour les deux). \\ \hline
Masse (55 vs 72 kg) & & $\checkmark$ & Différence de 17 kg ! Même génotype $\rightarrow$ différence due à l'alimentation et l'activité physique. \\ \hline
Couleur des yeux & $\checkmark$ & & Identique, ne change pas avec le milieu. \\ \hline
Teint de la peau & $\checkmark$ (base) & $\checkmark$ (bronzage) & La couleur \emph{de base} (mate) est héréditaire. Le teint \emph{très foncé} de Lou est un bronzage (acquis, réversible). \\ \hline
Groupe sanguin & $\checkmark$ & & Identique, déterminé génétiquement. \\ \hline
Langues parlées & & $\checkmark$ & Apprentissage culturel pur. \\ \hline
Cicatrice & & $\checkmark$ & Résultat d'une chirurgie (événement de vie), absent chez Lou. \\ \hline
\end{tabular}
\end{center}
\end{exoresolu}

\courssub{Étude de cas 3 — Le teint de la peau : couleur de base (héréditaire) vs bronzage (acquis)}

C'est l'exemple le plus fréquent de confusion au BEPC. Il faut distinguer deux composantes :

\begin{propriete}[Déterminisme de la pigmentation cutanée]
\begin{itemize}
    \item \textbf{La couleur de base (phototype) :} Déterminée par la quantité et le type de mélanine (eumélanine noire/brune, phéomélanine rouge/jaune) programmée génétiquement. C'est un \textbf{caractère héréditaire} (polygénique). Elle ne change pas au cours de la vie (sauf pathologie type vitiligo).
    \item \textbf{Le bronzage :} Augmentation transitoire de la production de mélanine par les mélanocytes \textbf{sous l'effet des rayons ultraviolets (UV) du soleil}. C'est une \textbf{réaction de défense} de la peau. C'est un \textbf{caractère acquis (modifié par l'environnement)}.
\end{itemize}
\end{propriete}

\begin{attention}[Erreur classique au BEPC : « La couleur de la peau est un caractère acquis »]
\textbf{FAUX.} Seule la \emph{variation} de couleur (bronzage) est acquise. La couleur \emph{constitutive} (celle de l'intérieur du bras, non exposé au soleil) est héréditaire. Un enfant de parents noirs né et élevé en Europe (peu de soleil) aura la peau noire (hérédité). Un enfant de parents blancs élevé au Cameroun (beaucoup de soleil) aura la peau blanche (hérédité) mais pourra bronzer (acquis). \textbf{Le bronzage ne se transmet pas aux enfants.}
\end{attention}

\courssub{Autres exemples de caractères acquis (modifiés par l'environnement)}

\begin{itemize}
    \item \textbf{Cicatrices, mutilations, piercings, scarifications traditionnelles :} Lésions tissulaires ou modifications corporelles intentionnelles. Disparaissent à la mort, absentes chez le nouveau-né.
    \item \textbf{Développement musculaire (hypertrophie) :} Résultat de l'entraînement sportif ou du travail physique. Un haltérophile a des muscles volumineux ; ses enfants naissent avec des muscles normaux.
    \item \textbf{Savoirs, langues, écriture, religion, coutumes :} \textbf{Capital culturel} acquis par l'éducation et l'imitation. Non génétique.
    \item \textbf{Maladies liées au mode de vie (non infectieuses) :} Obésité, diabète de type 2, hypertension artérielle, cirrhose alcoolique, certains cancers (tabac). Il existe souvent une \emph{prédisposition génétique} (hérédité), mais le \emph{déclenchement} nécessite l'environnement (alimentation, sédentarité, toxiques).
    \item \textbf{Allergies :} Prédisposition génétique (atopie) + exposition à l'allergène (environnement).
\end{itemize}

\begin{propriete}[Preuve que les caractères acquis ne sont pas héréditaires]
L'observation constante depuis l'Antiquité (et confirmée par la génétique moderne) montre que :
\begin{enumerate}
    \item Les enfants de parents \textbf{scarifiés} (scarifications rituelles au visage ou au corps, fréquentes dans certaines cultures camerounaises comme les Bamiléké ou les Haoussa) naissent \textbf{sans aucune scarification}.
    \item Les enfants de parents \textbf{musclés} (haltérophiles, manutentionnaires) naissent avec une musculature \textbf{normale}, pas hypertrophiée.
    \item Les enfants de parents \textbf{amputés} naissent avec \textbf{tous leurs membres}.
    \item Les enfants de parents \textbf{tatoués} ou \textbf{piercés} naissent sans tatouage ni piercing.
\end{enumerate}
\par\textbf{Conclusion rigoureuse :} Seules les informations inscrites dans l'ADN des \textbf{gamètes} (spermatozoïdes, ovocytes) sont transmises. Les modifications du corps (soma) survenues pendant la vie ne modifient pas l'ADN de la lignée germinale.
\end{propriete}

\courssub{Bilan synthétique : Un caractère = Hérédité + Environnement}

La très grande majorité des caractères observables (phénotype) résultent de l'\textbf{interaction} entre le programme génétique et les conditions de vie.

\begin{popfigure}
\begin{tikzpicture}[
    node distance=1.5cm and 2cm,
    box/.style={draw, thick, rounded corners, minimum width=3.5cm, minimum height=1.8cm, text centered, font=\small},
    arrow/.style={-Stealth, thick},
    plus/.style={font=\Large\bfseries, text=popdark}
]
\node[box, fill=popblueL, align=center] (G){\textbf{GÉNOTYPE}\\(Hérédité)\\\textit{Gènes, Allèles}\\\textit{ex: potentiel taille, groupe sanguin}};
\node[box, fill=poporangeL, right=of G, align=center] (E){\textbf{ENVIRONNEMENT}\\(Conditions de vie)\\\textit{Nutrition, Climat, Activités}\\\textit{Maladies, Éducation, Culture}};
\node[box, fill=popgreenL, below=of $(G)!0.5!(E)$, align=center] (P){\textbf{PHÉNOTYPE}\\(Caractère observable)\\\textit{Taille réelle, Masse, Teint,}\\\textit{Langue, Force, Santé}};
\draw[arrow, popblue] (G.south) -- ++(0,-0.5) -| (P.north west) node[plus, pos=0.25, left] {+};
\draw[arrow, poporange] (E.south) -- ++(0,-0.5) -| (P.north east) node[plus, pos=0.25, right] {+};
\node[align=center, font=\footnotesize, text=popmuted] at ($(P.south)+(0,-0.8)$) {Exemple : Taille adulte = Potentiel génétique (G) $\times$ Nutrition/Santé (E)};
\end{tikzpicture}
\legende{Schéma-bilan : Le phénotype (caractère observable) résulte de l'interaction entre le Génotype (information héréditaire) et l'Environnement (conditions de vie). La flèche « + » symbolise l'action conjointe des deux facteurs.}
\end{popfigure}

\begin{methode}[Comment déterminer l'origine d'un caractère ? (3 questions clés)]
Face à un caractère observé, pose-toi ces trois questions dans l'ordre :
\begin{enumerate}
    \item \textbf{Est-il présent dès la naissance (ou avant, in utero) ?}
    \begin{itemize}
        \item OUI $\rightarrow$ Fortement suspecté \textbf{héréditaire} (ou anomalie de développement).
        \item NON (apparaît plus tard) $\rightarrow$ Probablement \textbf{acquis} (ou héréditaire à expression tardive, ex: calvitie, maladie de Huntington).
    \end{itemize}
    \item \textbf{Se retrouve-t-il chez les parents / la fratrie ?}
    \begin{itemize}
        \item OUI (ressemblance familiale) $\rightarrow$ Argument pour \textbf{héréditaire}.
        \item NON (cas isolé) $\rightarrow$ Argument pour \textbf{acquis} (ou mutation \emph{de novo}, rare).
    \end{itemize}
    \item \textbf{Disparaît-il ou change-t-il si les conditions de vie changent ?}
    \begin{itemize}
        \item OUI (réversible : bronzage, masse musculaire, langue parlée) $\rightarrow$ \textbf{Acquis}.
        \item NON (irréversible : groupe sanguin, couleur des yeux, cicatrice ancienne) $\rightarrow$ \textbf{Héréditaire} ou \textbf{Acquis irréversible} (cicatrice, amputation).
    \end{itemize}
\end{enumerate}
\par\textbf{Règle d'or :} Seule la concordance des réponses (naissance + famille + stabilité) permet d'affirmer l'hérédité. La discordance signale l'environnement.
\end{methode}

\begin{savaistu}[La taille, indicateur de santé publique au Cameroun]
L'\textbf{Organisation Mondiale de la Santé (OMS)} et le \textbf{MINSANTÉ} utilisent la \textbf{taille pour l'âge} (et le poids pour la taille) comme indicateurs majeurs de l'état nutritionnel des populations.
\begin{itemize}
    \item Le \textbf{retard de croissance (stunting)} : taille trop basse pour l'âge (< -2 écarts-types par rapport à la médiane OMS). Il reflète une \textbf{malnutrition chronique} pendant les 1000 premiers jours (grossesse + 2 ans). Au Cameroun, environ \textbf{29\% des enfants de moins de 5 ans} sont en retard de croissance (EDS-MICS 2018).
    \item L'\textbf{émaciation (wasting)} : poids trop bas pour la taille. Reflet d'une malnutrition aiguë.
\end{itemize}
Ces enfants ont le \emph{même potentiel génétique} que les autres, mais l'environnement (pauvreté, insécurité alimentaire, infections répétées) les empêche de l'atteindre. Améliorer l'environnement (nutrition, santé, eau/assainissement) permet de \textbf{gagner des centimètres} en une génération : c'est la preuve vivante que la taille est un caractère à déterminisme mixte.
\end{savaistu}

\begin{aretenir}[Bilan de la section : Caractères héréditaires vs Caractères acquis]
\begin{itemize}
    \item \textbf{Caractère héréditaire :} Porté par les gènes (ADN des chromosomes), présent dans les gamètes, se transmet à la descendance, stable au cours de la vie (ex: groupe sanguin, couleur des yeux, empreintes digitales).
    \item \textbf{Caractère acquis (modifié par l'environnement) :} Apparaît au cours de la vie sous l'effet du milieu (nutrition, soleil, effort, apprentissage, traumatisme), \textbf{ne se transmet pas} aux enfants (ex: bronzage, masse musculaire, cicatrice, langues, obésité).
    \item \textbf{La plupart des caractères complexes (taille, intelligence, risque de maladies, personnalité) sont à déterminisme mixte :} $Phénotype = Génotype + Environnement + (G \times E)$.
    \item \textbf{Preuve de non-hérédité des acquis :} Les modifications corporelles des parents (scarifications, muscles, cicatrices) ne se retrouvent jamais chez le nouveau-né.
\end{itemize}
\end{aretenir}

\begin{exercice}[Entraînement BEPC — Origine des caractères]
\textbf{1.} Un couple de parents à peau mate (phénotype normal) a un enfant albinos (peau très claire, cheveux blancs, problèmes de vue). Quelle est l'origine du caractère « albinisme » chez l'enfant ? Justifie en précisant le génotype probable des parents.
\par\textbf{2.} Un lutteur traditionnel (mbapatte) très musclé a un fils qui, à 10 ans, a une morphologie fine et peu musclée. Le père dit : « Il n'a pas hérité de ma force. » Est-ce une affirmation scientifiquement correcte ? Explique.
\par\textbf{3.} Deux vrais jumeaux sont séparés à la naissance. L'un grandit au village (alimentation traditionnelle, travail aux champs, soleil), l'autre en ville (alimentation variée, école, bureau). À 30 ans, on compare leur groupe sanguin, leur taille, leur masse, leur teint, leur langue maternelle. Remplis le tableau :
\begin{center}
\begin{tabular}{|l|c|c|}\hline
\textbf{Caractère} & \textbf{Identique ou Différent ?} & \textbf{Origine (Héréditaire / Acquis / Mixte)} \\ \hline
Groupe sanguin & & \\ \hline
Taille & & \\ \hline
Masse corporelle & & \\ \hline
Teint de la peau (avant-bras) & & \\ \hline
Langue maternelle & & \\ \hline
\end{tabular}
\end{center}
\par\textbf{4.} Pourquoi les scarifications rituelles faites sur le visage des parents ne se transmettent-elles pas aux enfants ? Réponds en utilisant les termes « gamètes », « ADN », « soma ».
\end{exercice}

\begin{corrige}[Exercice n°1 — Entraînement BEPC]
\textbf{1.} L'albinisme est un caractère \textbf{héréditaire} à transmission \textbf{récessive}. Les parents, bien que phénotypiquement normaux (peau mate), sont tous deux \textbf{hétérozygotes} (génotype $A//a$). L'enfant a reçu l'allèle récessif $a$ de son père et l'allèle récessif $a$ de sa mère : son génotype est $a//a$ (homozygote récessif), d'où le phénotype albinos.
\par\textbf{2.} \textbf{Non, l'affirmation n'est pas correcte.} La force musculaire (hypertrophie) du père est un \textbf{caractère acquis} dû à l'entraînement intense (environnement). Elle n'est pas codée dans ses gènes (seul le \emph{potentiel} à développer du muscle l'est, c'est la part héréditaire). Le fils a hérité du \emph{potentiel génétique} (qui peut être grand), mais il ne l'a pas encore \emph{exprimé} car il n'a pas l'environnement (entraînement) pour le faire. Il n'a pas « hérité de la force acquise », il a hérité du \emph{génotype}.
\par\textbf{3.}
\begin{center}
\begin{tabular}{|l|c|c|}\hline
\textbf{Caractère} & \textbf{Identique / Différent} & \textbf{Origine} \\ \hline
Groupe sanguin & Identique & Héréditaire \\ \hline
Taille & Proche (Identique si nutrition similaire) & Mixte (Héréditaire + Nutrition) \\ \hline
Masse corporelle & Différent (Ville > Village souvent) & Acquis (Alimentation/Activité) \\ \hline
Teint (avant-bras) & Différent (Village plus foncé) & Acquis (Bronzage/Soleil) \\ \hline
Langue maternelle & Différent & Acquis (Culture/Éducation) \\ \hline
\end{tabular}
\end{center}
\par\textbf{4.} Les scarifications touchent le \textbf{soma} (corps, cellules somatiques). Elles ne modifient pas l'\textbf{ADN} présent dans les \textbf{gamètes} (spermatozoïdes/ovocytes). Seule l'information génétique contenue dans les gamètes est transmise à la descendance. L'enfant se développe à partir du zygote (fusion de deux gamètes intacts), sans trace des lésions cutanées parentales.
\end{corrige}

\coursec{Bilan : unité et diversité, applications à la vie courante}

Dans la section précédente, nous avons vu que l'environnement peut modifier certains de nos caractères : une cicatrice, une peau bronzée, une musculature développée par le sport. Mais ces modifications ne s'ajoutent pas au hasard : elles s'inscrivent sur un fond commun, celui de l'espèce humaine, et sur un héritage propre à chacun, reçu des parents. Il est temps de rassembler toutes les pièces du puzzle : comment un même être humain peut-il être à la fois semblable à tous les autres et absolument unique ? C'est ce bilan que nous dressons ici, avant de découvrir comment ces connaissances servent concrètement, chaque jour, dans les hôpitaux, les services d'identification et la vie citoyenne au Cameroun.

\courssub{Synthèse : l'unité et la diversité de l'espèce humaine}

\textbf{Une observation de départ}\par

Observe les élèves de ta classe. Tous ont deux yeux, un nez, une bouche, deux bras et deux jambes ; tous marchent debout, parlent un langage articulé et possèdent un cerveau développé. Pourtant, aucun n'est identique à un autre : les tailles, les teints, les groupes sanguins, les empreintes digitales diffèrent. Cette double constatation — \cle{ressemblance} générale et \cle{différences} individuelles — est le cœur de cette leçon.

\begin{propriete}[Bilan de la leçon : unité et diversité]
L'espèce humaine est à la fois \cle{une} et \cle{diverse}.
\begin{itemize}
\item \textbf{Une} : tous les humains partagent les mêmes caractères fondamentaux (bipédie, station debout, cerveau développé, langage articulé, main préhensile avec pouce opposable) et peuvent se croiser entre eux en donnant des descendants \textbf{féconds}. C'est le critère d'appartenance à une même espèce.
\item \textbf{Diverse} : au sein de cette unité, chaque individu se distingue par des \textbf{caractères héréditaires variables} (groupe sanguin, couleur des yeux, forme du visage) et par des \textbf{caractères modifiés par l'environnement} (cicatrices, bronzage, masse corporelle, savoirs acquis).
\end{itemize}
\end{propriete}

\begin{attention}[Ne pas opposer hérédité et environnement]
L'erreur la plus fréquente consiste à classer chaque caractère « soit héréditaire, soit acquis ». En réalité, de nombreux caractères dépendent \textbf{des deux à la fois} : la taille, la masse corporelle ou la santé résultent d'un héritage reçu des parents \emph{et} des conditions de vie (alimentation, soins, activité physique). Un enfant peut avoir reçu un potentiel de grande taille, mais rester petit s'il est mal nourri pendant sa croissance. Au BEPC, précise toujours cette nuance quand le sujet porte sur la taille, la masse ou la santé.
\end{attention}

\courssub{Classer les caractères : un tableau-bilan}

Pour vérifier que tu maîtrises la leçon, l'exercice classique consiste à trier une liste de caractères en trois colonnes. Voici le tableau de référence à connaître.

\begin{center}
\begin{tabularx}{\linewidth}{|X|X|X|}
\hline
\rowcolor{popblueL} \textbf{Caractères de l'espèce humaine} & \textbf{Caractères héréditaires variables} & \textbf{Caractères acquis (modifiés par l'environnement)} \\
\hline
Station debout, bipédie & Groupe sanguin (A, B, AB, O) & Cicatrices, tatouages \\
\hline
Cerveau développé & Couleur des yeux & Bronzage \\
\hline
Langage articulé & Forme du visage, teint & Masse corporelle (selon l'alimentation) \\
\hline
Main préhensile, pouce opposable & Empreintes digitales & Musculature développée par le sport \\
\hline
Croisements féconds entre tous les humains & Drépanocytose (caractère héréditaire) & Langue parlée, savoirs, habits \\
\hline
\end{tabularx}
\end{center}

\begin{methode}[Réussir le tri des caractères]
Pour classer un caractère sans erreur, pose-toi trois questions dans l'ordre :
\begin{enumerate}
\item \textbf{Tous les humains le possèdent-ils ?} Si oui, c'est un caractère de l'espèce (ex. : la bipédie).
\item \textbf{Est-il transmis par les parents et présent dès la naissance, mais variable d'un individu à l'autre ?} Si oui, c'est un caractère héréditaire variable (ex. : le groupe sanguin).
\item \textbf{Est-il apparu sous l'effet du milieu de vie, sans être transmis aux enfants ?} Si oui, c'est un caractère acquis (ex. : une cicatrice).
\end{enumerate}
\end{methode}

\courssub{L'individu : au croisement de l'hérédité et de l'environnement}

Chaque être humain est le résultat d'une rencontre : celle d'un \cle{héritage} reçu de ses parents et d'un \cle{environnement} dans lequel il grandit, le tout à l'intérieur d'une espèce unique qui lui donne ses caractères fondamentaux. Le schéma suivant résume cette idée essentielle.

\begin{popfigure}
\begin{center}
\begin{tikzpicture}[font=\small]
% Cadre espèce
\draw[rounded corners=8pt, popblue!40, fill=popblueL, thick] (-5.2,-3.2) rectangle (5.2,3.2);
\node[popdark, font=\small\bfseries, align=center] at (0,2.75) {L'espèce humaine : caractères communs à tous\\ (bipédie, langage, cerveau développé)};
% Boîte hérédité
\draw[rounded corners=5pt, fill=poppurpleL, draw=poppurple, thick] (-4.6,0.4) rectangle (-1.4,2.1);
\node[align=center, popdark] at (-3,1.25) {\textbf{Hérédité}\\[2pt] caractères reçus des parents\\ (groupe sanguin, couleur\\ des yeux, empreintes)};
% Boîte environnement
\draw[rounded corners=5pt, fill=popgreenL, draw=popgreen, thick] (1.4,0.4) rectangle (4.6,2.1);
\node[align=center, popdark] at (3,1.25) {\textbf{Environnement}\\[2pt] alimentation, climat,\\ mode de vie, éducation};
% Individu
\draw[fill=poporangeL, draw=poporange, very thick] (0,-1.6) circle (1.15);
\node[align=center, popdark, font=\bfseries] at (0,-1.6) {L'INDIVIDU\\ unique};
% Flèches
\draw[->, very thick, poppurple] (-3,0.35) -- (-0.9,-0.9);
\draw[->, very thick, popgreen!60!black] (3,0.35) -- (0.9,-0.9);
\draw[->, very thick, popblue] (0,-2.85) -- (0,-2.35);
\node[popblue, font=\footnotesize] at (2.1,-2.6) {cadre commun à tous};
\end{tikzpicture}
\end{center}
\legende{Schéma de synthèse : chaque individu est unique car il naît de la combinaison d'une hérédité propre et d'un environnement propre, au sein d'une seule espèce humaine qui fournit à tous les mêmes caractères fondamentaux.}
\end{popfigure}

\begin{savaistu}[Une diversité qui est une richesse]
La diversité humaine n'est pas un défaut : c'est une \textbf{richesse biologique}. Grâce à la variété de ses caractères, notre espèce a pu s'adapter à tous les milieux de la planète : les populations des hauts plateaux supportent mieux le manque d'oxygène, celles des régions très ensoleillées ont une peau plus riche en mélanine, qui protège mieux des rayons du soleil. Au Cameroun, de la forêt du Sud aux savanes de l'Extrême-Nord, cette diversité témoigne de la formidable capacité d'adaptation de l'espèce humaine.
\end{savaistu}

\courssub{Application 1 : les transfusions sanguines dans les hôpitaux}

Au service des urgences de l'Hôpital Central de Yaoundé, une accidentée de la route a perdu beaucoup de sang. Avant toute transfusion, le médecin détermine son \textbf{groupe sanguin}. Pourquoi ? Parce que le groupe sanguin est un \cle{caractère héréditaire} : il est fixé dès la naissance, ne change jamais au cours de la vie, et détermine la compatibilité entre le sang du donneur et celui du receveur. Une transfusion incompatible provoque une réaction grave (agglutination des globules rouges), parfois mortelle.

\begin{exemplebox}[Le donneur universel : le groupe O]
Les personnes de groupe \textbf{O} peuvent donner leur sang à \textbf{tous les receveurs} (groupes O, A, B et AB) : on les appelle \textbf{donneurs universels}. À l'inverse, les personnes de groupe \textbf{AB} peuvent recevoir le sang de tous les groupes : ce sont les \textbf{receveurs universels}. C'est pourquoi les banques de sang des hôpitaux camerounais recherchent particulièrement les donneurs de groupe O, très demandés lors des urgences. Connaître son groupe sanguin et donner son sang, c'est sauver des vies grâce à un caractère héréditaire.
\end{exemplebox}

\courssub{Application 2 : l'identification des personnes}

Comment prouver qu'une personne est bien celle qu'elle prétend être ? En utilisant ses caractères \textbf{uniques}. Les \cle{empreintes digitales}, dessins des crêtes de la peau au bout des doigts, sont des caractères dont la disposition fine est propre à chaque individu : même de vrais jumeaux n'ont pas les mêmes empreintes. C'est pourquoi elles sont relevées lors de l'établissement de la \textbf{Carte Nationale d'Identité} et des passeports au Cameroun. De même, la \textbf{reconnaissance faciale}, qui mesure les proportions du visage (caractères en grande partie héréditaires), permet de déverrouiller un téléphone ou de contrôler une identité à l'aéroport. La diversité biologique devient ainsi un outil de sécurité.

\courssub{Application 3 : la santé, entre hérédité et mode de vie}

La distinction entre caractères héréditaires et caractères acquis a des conséquences directes sur la prévention des maladies.

\begin{itemize}
\item \textbf{Maladies héréditaires} : la \cle{drépanocytose} (ou anémie falciforme), fréquente en Afrique centrale, est une maladie du sang transmise par les parents : les globules rouges des malades prennent une forme de faucille et transportent mal l'oxygène. On ne peut pas « l'attraper » ni la guérir par l'hygiène de vie ; la prévention passe par l'information des familles et le dépistage. Savoir qu'elle est héréditaire évite les fausses accusations et les croyances erronées.
\item \textbf{Maladies liées au mode de vie} : l'obésité, certaines maladies du cœur ou les caries dentaires dépendent largement de l'\textbf{hygiène alimentaire} et de l'activité physique. Celles-là peuvent être prévenues par nos comportements : manger équilibré, limiter les graisses et les boissons sucrées, faire du sport.
\end{itemize}

Comprendre l'origine d'un problème de santé permet donc de choisir la bonne attitude : dépistage et accompagnement pour l'héréditaire, changement d'habitudes pour l'acquis.

\courssub{Application 4 : une leçon de citoyenneté}

La science apporte ici un enseignement qui dépasse la biologie. Puisque \textbf{tous les humains appartiennent à une seule et même espèce} — mêmes caractères fondamentaux, croisements féconds entre tous les peuples —, il n'existe aucun fondement biologique aux hiérarchies entre groupes humains. Les différences visibles (teint, forme du visage) sont des variations superficielles à l'intérieur d'une unité profonde.

\begin{aretenir}[Unité biologique et égalité des droits]
L'unité de l'espèce humaine justifie l'\textbf{égalité de dignité et de droits} entre tous les êtres humains. Les préjugés ethniques, tribaux ou racistes sont contraires aux faits scientifiques : biologiquement, nous sommes une seule famille humaine, diverse dans ses apparences, une dans sa nature.
\end{aretenir}

\courssub{Méthode : rédiger une réponse justifiée au BEPC}

Au BEPC, une question de synthèse demande souvent de « dégager une conclusion justifiée » à partir de documents. Une bonne réponse se construit toujours en trois temps.

\begin{methode}[J'affirme, je justifie, je conclus]
\begin{enumerate}
\item \textbf{J'affirme} : j'énonce clairement l'idée (ex. : « L'espèce humaine est une et diverse »).
\item \textbf{Je justifie} : j'appuie l'affirmation sur un fait ou un résultat tiré des documents (ex. : « le document montre que tous les individus possèdent deux bras, deux jambes et marchent debout, mais que leurs groupes sanguins diffèrent »).
\item \textbf{Je conclus} : je relie le tout à la notion du cours (ex. : « les caractères communs traduisent l'unité de l'espèce ; les caractères variables, héréditaires ou acquis, traduisent sa diversité »).
\end{enumerate}
\end{methode}

\begin{exoresolu}[Classer des caractères et conclure]
\textbf{Énoncé.} Voici une liste de caractères observés chez des élèves : (a) la bipédie ; (b) le groupe sanguin A ; (c) une cicatrice au genou ; (d) le langage articulé ; (e) la couleur des yeux ; (f) le bronzage. 1) Classe ces caractères en trois catégories. 2) Rédige une conclusion justifiée sur l'unité et la diversité de l'espèce humaine.
\tcblower
\textbf{Solution détaillée.}
1) \textbf{Classement.}
\begin{itemize}
\item \emph{Caractères de l'espèce} (communs à tous les humains) : (a) la bipédie et (d) le langage articulé.
\item \emph{Caractères héréditaires variables} (transmis par les parents, différents selon les individus) : (b) le groupe sanguin A et (e) la couleur des yeux.
\item \emph{Caractères acquis} (dus à l'environnement, non transmis aux enfants) : (c) la cicatrice et (f) le bronzage.
\end{itemize}
2) \textbf{Conclusion rédigée en trois temps.} \emph{J'affirme} : l'espèce humaine est une et diverse. \emph{Je justifie} : tous les élèves étudiés marchent debout et parlent un langage articulé, ce qui montre qu'ils partagent les mêmes caractères fondamentaux ; cependant, leurs groupes sanguins et leurs couleurs des yeux diffèrent (caractères héréditaires variables), et certains portent des cicatrices ou un bronzage (caractères acquis sous l'effet du milieu). \emph{Je conclus} : les caractères communs expriment l'unité de l'espèce humaine, tandis que les caractères héréditaires variables et les caractères acquis expliquent que chaque individu soit unique.
\end{exoresolu}

\begin{exercice}[À toi de jouer]
Dans un centre de santé de Douala, on relève chez quatre patients : la taille, le groupe sanguin, la présence d'un tatouage et la station debout. 1) Indique, pour chaque caractère, sa catégorie. 2) Explique pourquoi la taille ne peut pas être classée simplement comme « héréditaire ». 3) Rédige une conclusion justifiée en trois temps sur l'unité et la diversité humaines.
\end{exercice}

\begin{corrige}[Exercice]
1) Station debout : caractère de l'espèce. Groupe sanguin : caractère héréditaire variable. Tatouage : caractère acquis. Taille : caractère mixte (hérédité et environnement).
2) La taille dépend d'un potentiel hérité des parents, mais aussi de l'alimentation et de la santé pendant la croissance : une malnutrition peut empêcher ce potentiel de se réaliser. Elle dépend donc des deux facteurs.
3) \emph{J'affirme} : l'espèce humaine est une et diverse. \emph{Je justifie} : les quatre patients se tiennent tous debout (caractère commun de l'espèce), mais leurs groupes sanguins diffèrent (hérédité variable) et l'un porte un tatouage (caractère acquis). \emph{Je conclus} : l'unité vient des caractères de l'espèce, la diversité vient des caractères héréditaires variables et des modifications dues à l'environnement.
\end{corrige}

\begin{aretenir}[L'essentiel de la leçon]
\begin{itemize}
\item Tous les humains partagent les \textbf{caractères de l'espèce} : bipédie, cerveau développé, langage, croisements féconds $\Rightarrow$ \textbf{unité}.
\item Chacun se distingue par des caractères \textbf{héréditaires variables} (groupe sanguin, empreintes) et \textbf{acquis} (cicatrices, bronzage) $\Rightarrow$ \textbf{diversité}.
\item Certains caractères (taille, masse, santé) dépendent \textbf{à la fois} de l'hérédité et de l'environnement.
\item Ces connaissances servent la médecine (transfusions, dépistage de la drépanocytose), l'identification (CNI, empreintes) et la citoyenneté (égalité des droits, rejet des préjugés).
\end{itemize}
\end{aretenir}

\newpage

\coursec{Activité d'intégration}

\courssub{Objectifs de l'activité}
\begin{itemize}
    \item Mobiliser les notions d'\cle{unité de l'espèce humaine}, de \cle{caractère héréditaire} et de \cle{caractère acquis (modifié par l'environnement)} pour analyser des données réelles issues d'enquêtes menées au Cameroun.
    \item Pratiquer une \cle{démarche d'investigation scientifique} : observer des documents, formuler des hypothèses, interpréter des données chiffrées, argumenter et rédiger une conclusion justifiée.
    \item Développer des \cle{savoir-faire} transversaux : travail collaboratif, lecture de tableaux et d'histogrammes, communication orale et écrite.
\end{itemize}

\courssub{Situation}
Dans le cadre d'un projet d'éducation à la santé et à la citoyenneté, votre classe de 3ème du \textbf{Lycée Bilingue de Bafoussam} participe à une enquête scientifique pilotée par le \textbf{Centre Pasteur du Cameroun} et l'\textbf{INS (Institut National de la Statistique)}. La question centrale est : \og \emph{Nos différences viennent-elles de nos parents ou de notre environnement ?} \fg\par
Trois documents, anonymisés, sont mis à votre disposition. Ils portent sur des populations camerounaises (région de l'Ouest, départements de la Menoua et du Koung-Khi).

\medskip
\noindent\textbf{Document 1 : Enquête familiale sur deux générations (10 familles, 40 individus).}
Le tableau ci-dessous récapitule, pour chaque famille, les \cle{tailles} (en cm), \cle{masses} (en kg) et \cle{groupes sanguins} (système ABO) des parents (P) et de leurs deux enfants (E1, E2) âgés de 18 à 22 ans. Tous sont nés et ont grandi à Bafoussam (milieu urbain).

\begin{center}
\begin{tabularx}{\linewidth}{|c|*{4}{>{\centering\arraybackslash}X|}}\hline
\rowcolor{popblueL}
\textbf{Famille} & \textbf{Père (P)} & \textbf{Mère (P)} & \textbf{Enfant 1 (E1)} & \textbf{Enfant 2 (E2)} \\ \hline
1 & T: 178, M: 78, GS: A & T: 162, M: 60, GS: O & T: 175, M: 70, GS: A & T: 165, M: 58, GS: O \\ \hline
2 & T: 168, M: 65, GS: B & T: 158, M: 55, GS: O & T: 166, M: 62, GS: B & T: 160, M: 54, GS: O \\ \hline
3 & T: 182, M: 85, GS: AB & T: 170, M: 68, GS: A & T: 179, M: 80, GS: A & T: 172, M: 72, GS: B \\ \hline
4 & T: 170, M: 72, GS: O & T: 160, M: 58, GS: O & T: 168, M: 68, GS: O & T: 162, M: 56, GS: O \\ \hline
5 & T: 175, M: 75, GS: A & T: 165, M: 62, GS: B & T: 173, M: 71, GS: A & T: 167, M: 60, GS: B \\ \hline
6 & T: 165, M: 63, GS: O & T: 155, M: 52, GS: A & T: 163, M: 60, GS: O & T: 158, M: 53, GS: A \\ \hline
7 & T: 180, M: 82, GS: B & T: 168, M: 65, GS: AB & T: 177, M: 78, GS: B & T: 170, M: 66, GS: A \\ \hline
8 & T: 172, M: 70, GS: A & T: 160, M: 59, GS: O & T: 169, M: 67, GS: A & T: 163, M: 57, GS: O \\ \hline
9 & T: 169, M: 66, GS: O & T: 159, M: 54, GS: B & T: 167, M: 64, GS: O & T: 161, M: 55, GS: B \\ \hline
10 & T: 177, M: 76, GS: A & T: 164, M: 61, GS: O & T: 174, M: 73, GS: A & T: 166, M: 59, GS: O \\ \hline
\end{tabularx}
\end{center}
\legende{Document 1 : Tailles (T, en cm), masses (M, en kg) et groupes sanguins (GS, système ABO) de 10 familles bafoussamaises (2 générations).}

\medskip
\noindent\textbf{Document 2 : Cas de vrais jumeaux monozygotes séparés à la naissance.}
\emph{Jean} et \emph{Paul} sont des jumeaux \cle{monozygotes} (issus d'un même ovule fécondé par un même spermatozoïde, qui s'est divisé en deux embryons : ils possèdent donc \emph{exactement la même information génétique}). Séparés à la naissance (adoption), ils ont grandi dans des milieux très différents :
\begin{itemize}
    \item \textbf{Jean} a été élevé à \textbf{Yaoundé} (quartier Bastos, milieu urbain favorisé) : alimentation variée et équilibrée (protéines animales quotidiennes, fruits, légumes), scolarité continue, accès aux soins de santé (vaccinations, suivi médical), pratique sportive régulière (football, natation).
    \item \textbf{Paul} a été élevé dans un \textbf{village de l'arrondissement de Santchou} (milieu rural défavorisé, zone de culture du riz et du maïs) : alimentation à base de glucides (riz, maïs, manioc) avec peu de protéines animales, scolarité interrompue à 12 ans pour les travaux champêtres, accès aux soins limité (centre de santé à 15 km), pas de pratique sportive organisée.
\end{itemize}
À l'âge de 25 ans, les caractères suivants sont mesurés :

\begin{center}
\begin{tabularx}{\linewidth}{|l|>{\centering\arraybackslash}X|>{\centering\arraybackslash}X|}\hline
\rowcolor{popblueL}
\textbf{Caractère} & \textbf{Jean (Yaoundé)} & \textbf{Paul (Santchou)} \\ \hline
Groupe sanguin (ABO) & A & A \\ \hline
Couleur des yeux & Marron foncé & Marron foncé \\ \hline
Type de lobe d'oreille & Libre (détaché) & Libre (détaché) \\ \hline
Taille & 1,78 m & 1,65 m \\ \hline
Masse & 72 kg & 54 kg \\ \hline
Indice de Masse Corporelle (IMC) & 22,7 kg/m$^2$ & 19,9 kg/m$^2$ \\ \hline
Tension artérielle (systolique/diastolique) & 118 / 72 mmHg & 132 / 84 mmHg \\ \hline
Niveau d'études & Licence (Bac+3) & Niveau 3ème (BEPC non obtenu) \\ \hline
Langue courante & Français, Anglais, Ewondo & Fe'efe'e, Pidgin, peu de Français \\ \hline
\end{tabularx}
\end{center}
\legende{Document 2 : Comparaison de deux vrais jumeaux monozygotes élevés dans des milieux contrastés.}

\medskip
\noindent\textbf{Document 3 : Histogramme des tailles d'élèves de 3ème dans deux établissements.}
Une campagne de mesure de la taille a été réalisée en mars 2026 auprès de tous les élèves de 3ème (âge moyen 14,5 ans) de deux collèges publics :
\begin{itemize}
    \item \textbf{Collège A :} \emph{Lycée de Nkolbisson} (Yaoundé, zone urbaine favorisée, cantine scolaire, infirmerie). Effectif : 320 élèves.
    \item \textbf{Collège B :} \emph{CES de Banka} (arrondissement de Banka, zone rurale enclavée, pas de cantine, poste de santé à 8 km). Effectif : 285 élèves.
\end{itemize}
Les résultats sont présentés sous forme d'histogrammes (classes de 3 cm de largeur).

\begin{center}
\begin{popfigure}
\begin{tikzpicture}[scale=0.7]
\begin{axis}[
    ybar,
    bar width=6pt,
    width=\linewidth,
    height=8cm,
    xlabel={Taille (cm)},
    ylabel={Effectif},
    xtick={141,144,147,150,153,156,159,162,165,168,171,174,177},
    xticklabel style={rotate=45, anchor=east, font=\tiny},
    ymin=0,
    ymax=50,
    legend style={at={(0.5,-0.32)}, anchor=north, legend columns=-1},
    title={Répartition des tailles : Collège A (Yaoundé) vs Collège B (Banka)},
    title style={font=\small},
]
\addplot[fill=popblue!70, draw=popblue] coordinates {
    (141,2) (144,5) (147,12) (150,25) (153,38) (156,42) (159,35) (162,28) (165,18) (168,10) (171,4) (174,1) (177,0)
};
\addplot[fill=poporange!70, draw=poporange] coordinates {
    (141,5) (144,15) (147,30) (150,40) (153,38) (156,32) (159,25) (162,18) (165,12) (168,6) (171,3) (174,1) (177,0)
};
\legend{Collège A (Yaoundé), Collège B (Banka)}
\end{axis}
\end{tikzpicture}
\legende{Document 3 : Histogrammes comparatifs des tailles d'élèves de 3ème (14,5 ans) en milieu urbain favorisé (bleu) et rural défavorisé (orange). Chaque classe a une largeur de 3 cm.}
\end{popfigure}
\end{center}

\courssub{Consignes}
Travaillez par groupes de 4 à 5 élèves. Chaque groupe rend un rapport écrit structuré répondant aux questions ci-dessous. Un porte-parole présentera la synthèse à l'oral (3 min max).

\begin{enumerate}
    \item \textbf{Unité de l'espèce :} À partir du Document 1, relevez \textbf{au moins trois caractères communs à tous les 40 individus} (parents et enfants). Que conclure sur l'unité de l'espèce humaine ?
    \item \textbf{Caractères héréditaires :} En vous appuyant \textbf{uniquement sur le Document 1}, identifiez le(s) caractère(s) qui semblent \textbf{transmis des parents aux enfants}. Justifiez votre réponse en citant des \textbf{données précises du tableau} (ex : \og Famille 3 : père AB, mère A $\to$ enfants A et B \fg). Pourquoi la taille et la masse ne peuvent-elles pas être considérées comme \emph{strictement} héréditaires d'après ce document seul ?
    \item \textbf{Caractères modifiés par l'environnement :} Comparez Jean et Paul (Document 2). Classifiez \textbf{chacun des 9 caractères listés} dans l'un des deux tableaux ci-dessous : \og Caractère héréditaire (identique chez les jumeaux) \fg ou \og Caractère modifié par l'environnement (différent chez les jumeaux) \fg. Justifiez \textbf{deux choix} par une phrase explicative.
    \item \textbf{Influence du milieu de vie sur un caractère quantitatif :} Analysez l'histogramme (Document 3).
    \begin{enumerate}
        \item Quelle est la \textbf{classe de taille la plus fréquentée (classe modale)} pour chaque collège ?
        \item Les élèves du Collège A sont-ils \emph{globalement} plus grands que ceux du Collège B ? Justifiez par une lecture de la distribution (position, étalement).
        \item Proposez \textbf{deux facteurs environnementaux} (liés au milieu de vie) susceptibles d'expliquer cet écart de taille à 14,5 ans.
    \end{enumerate}
    \item \textbf{Synthèse argumentée :} Rédigez, \textbf{en cinq lignes maximum}, une conclusion répondant à la question initiale : \og \emph{Nos différences viennent-elles de nos parents ou de notre environnement ?} \fg. Votre texte doit utiliser les termes \cle{caractères de l'espèce}, \cle{caractères héréditaires} et \cle{caractères acquis (ou modifiés par l'environnement)}.
\end{enumerate}

\courssub{Ressources à mobiliser}
\begin{aretenir}[Rappels de cours utiles pour l'activité]
\begin{itemize}
    \item \textbf{Caractère :} tout trait observable, mesurable ou descriptible chez un individu (ex : taille, groupe sanguin, langue parlée).
    \item \textbf{Caractères de l'espèce humaine :} caractères communs à \emph{tous} les êtres humains (ex : bipédie, main préhensile, cerveau développé, langage articulé, 46 chromosomes dans chaque cellule, appartenance au système ABO \emph{comme système} — mais pas le même groupe sanguin pour tous).
    \item \textbf{Caractère héréditaire :} caractère transmis des parents aux enfants par l'intermédiaire de l'information génétique (ADN) portée par les chromosomes. Il se retrouve chez les descendants et est \emph{identique} chez de vrais jumeaux monozygotes.
    \item \textbf{Caractère acquis (modifié par l'environnement) :} caractère qui dépend des conditions de vie (nutrition, santé, éducation, culture, activité physique...). Il \emph{diffère} chez de vrais jumeaux élevés dans des milieux différents.
    \item \textbf{Caractère à déterminisme mixte :} la plupart des caractères quantitatifs (taille, masse, IMC, tension artérielle) résultent de l'action combinée de \textbf{l'hérédité et de l'environnement}. L'hérédité fixe un \emph{potentiel}, l'environnement détermine la \emph{réalisation} de ce potentiel.
    \item \textbf{Lecture d'histogramme :} la classe modale est la classe dont l'effectif est le plus grand (barre la plus haute). La position globale de deux distributions se compare par leurs valeurs centrales ; l'étalement se compare par la largeur des distributions.
\end{itemize}
\end{aretenir}

\courssub{Démarche et corrigé commenté}
\begin{exoresolu}[Corrigé détaillé de l'activité d'intégration]
\textbf{Question 1 — Unité de l'espèce (Document 1).}
\par
\textbf{Observation :} dans le tableau, \emph{tous} les 40 individus présentent :
\begin{itemize}
    \item un \textbf{groupe sanguin} appartenant au système ABO (A, B, AB ou O) : le \emph{système} ABO est commun à tous les humains, même si le groupe précis varie d'un individu à l'autre ;
    \item une \textbf{taille} et une \textbf{masse} mesurables, caractères quantitatifs que l'on retrouve chez tout être humain ;
    \item (implicite dans le contexte) les caractères généraux de l'espèce : station debout (bipédie), main préhensile, langage articulé, cerveau développé, 46 chromosomes par cellule.
\end{itemize}
\textbf{Conclusion :} ces caractères partagés par l'ensemble des individus étudiés (et par tout être humain) sont les \cle{caractères de l'espèce humaine}. Ils attestent de l'\textbf{unité de l'espèce} \emph{Homo sapiens} : malgré leurs différences, tous les humains appartiennent à une seule et même espèce.

\medskip
\textbf{Question 2 — Caractères héréditaires (Document 1).}
\par
\textbf{Analyse du groupe sanguin :} le groupe sanguin ABO est déterminé par l'information génétique reçue des deux parents : chaque enfant reçoit une moitié de l'information génétique de son père et une moitié de sa mère.
\begin{itemize}
    \item \textbf{Famille 1 :} père A, mère O $\to$ enfants A et O : chaque enfant possède un groupe \emph{déjà présent chez l'un des parents}. \textbf{Compatible avec une transmission héréditaire.}
    \item \textbf{Famille 3 :} père AB, mère A $\to$ enfants A et B : l'enfant 2 est de groupe B, qu'il ne peut recevoir que de son père AB. \textbf{Compatible.}
    \item \textbf{Famille 7 :} père B, mère AB $\to$ enfants B et A : l'enfant 2 est de groupe A, reçu de sa mère AB. \textbf{Compatible.}
    \item \textbf{Dans les 10 familles :} le groupe sanguin de chaque enfant est \emph{toujours} explicable à partir des groupes des deux parents. \textbf{Aucune exception n'est observée.}
\end{itemize}
\textbf{Conclusion :} le \textbf{groupe sanguin ABO} est un \cle{caractère héréditaire} : il se transmet des parents aux enfants selon des règles précises, et il ne change jamais au cours de la vie, quelles que soient les conditions de vie.

\par
\textbf{Analyse de la taille et de la masse :}
\begin{itemize}
    \item \textbf{Famille 1 :} parents 178 cm et 162 cm $\to$ enfants 175 cm et 165 cm : les enfants sont \emph{intermédiaires} entre les parents, mais ne leur sont pas identiques.
    \item \textbf{Famille 3 :} parents 182 cm et 170 cm $\to$ enfants 179 cm et 172 cm.
    \item \textbf{Famille 4 :} parents 170 cm et 160 cm $\to$ enfants 168 cm et 162 cm : les enfants sont légèrement \emph{plus grands} que la moyenne des parents.
\end{itemize}
La taille et la masse \textbf{ne suivent pas un schéma de transmission simple} : elles varient de façon continue et les enfants sont souvent intermédiaires, parfois au-delà des valeurs parentales. \textbf{D'après ce document seul, on ne peut pas les considérer comme strictement héréditaires} : ce sont des caractères à \textbf{déterminisme mixte} (hérédité + environnement), ce que les Documents 2 et 3 vont confirmer.

\medskip
\textbf{Question 3 — Caractères héréditaires vs acquis (Document 2).}
\par
Jean et Paul sont \textbf{monozygotes} $\implies$ leur \textbf{information génétique est identique}. Toute différence observée entre eux est donc \emph{nécessairement} due à l'environnement (nutrition, santé, éducation, culture, activité physique...). C'est une \og expérience naturelle \fg{} très précieuse pour le biologiste.

\begin{center}
\begin{tabularx}{\linewidth}{|>{\centering\arraybackslash}X|>{\centering\arraybackslash}X|}\hline
\rowcolor{popblueL}
\textbf{Caractère héréditaire (identique)} & \textbf{Caractère modifié par l'environnement (différent)} \\ \hline
Groupe sanguin (A) & Taille (1,78 m contre 1,65 m) \\ \hline
Couleur des yeux (marron foncé) & Masse (72 kg contre 54 kg) \\ \hline
Type de lobe d'oreille (libre) & IMC (22,7 contre 19,9 kg/m$^2$) \\ \hline
& Tension artérielle (118/72 contre 132/84 mmHg) \\ \hline
& Niveau d'études (Licence contre niveau 3ème) \\ \hline
& Langue courante (Français/Anglais/Ewondo contre Fe'efe'e/Pidgin) \\ \hline
\end{tabularx}
\end{center}

\par
\textbf{Justifications (exemples attendus) :}
\begin{itemize}
    \item \textbf{Taille :} Jean (1,78 m) est plus grand que Paul (1,65 m) car il a bénéficié, pendant toute sa croissance, d'une \textbf{alimentation riche en protéines animales, en vitamines et en sels minéraux}, alors que Paul a eu une alimentation surtout glucidique (riz, maïs, manioc), pauvre en protéines. Leur potentiel génétique de taille était le même, mais seul Jean a pu le réaliser pleinement.
    \item \textbf{Langue courante :} Jean parle français et anglais (langues apprises à l'école à Yaoundé) ; Paul parle fe'efe'e et pidgin (langues de son entourage villageois). La langue n'est \textbf{pas inscrite dans les gènes} : elle s'acquiert par \textbf{imitation et apprentissage} dans le milieu culturel où l'on grandit.
\end{itemize}

\medskip
\textbf{Question 4 — Influence du milieu sur la taille (Document 3).}
\par
\textbf{a) Classes modales (barres les plus hautes) :}
\begin{itemize}
    \item \textbf{Collège A (Yaoundé) :} classe \textbf{[156 ; 159 cm[} (effectif 42 élèves).
    \item \textbf{Collège B (Banka) :} classe \textbf{[150 ; 153 cm[} (effectif 40 élèves).
\end{itemize}
\par
\textbf{b) Comparaison globale :} oui, les élèves du Collège A sont \textbf{globalement plus grands} que ceux du Collège B.
\begin{itemize}
    \item \textbf{Position :} toute la distribution du Collège A est \textbf{décalée vers les grandes tailles} : classe modale à [156 ; 159[ contre [150 ; 153[, et des effectifs encore importants jusqu'à 168--171 cm, alors que le Collège B a ses plus gros effectifs entre 144 et 156 cm.
    \item \textbf{Étalement :} le Collège B présente des effectifs non négligeables dans les plus petites classes (141--147 cm), presque vides au Collège A : les élèves les plus petits sont nettement plus nombreux à Banka.
\end{itemize}
\par
\textbf{c) Deux facteurs environnementaux explicatifs :}
\begin{enumerate}
    \item \textbf{Nutrition :} cantine scolaire à Nkolbisson (repas équilibrés quotidiens : protéines, lait, légumes) contre absence de cantine à Banka (repas familiaux souvent essentiellement glucidiques, avec des carences en protéines, fer, zinc et vitamines) : or la croissance en taille exige des protéines et des micronutriments en quantité suffisante.
    \item \textbf{Accès aux soins :} infirmerie sur place et hôpitaux proches à Yaoundé (dépistage et traitement rapides du paludisme et des vers intestinaux) contre un poste de santé à 8 km à Banka (retard de prise en charge ; les parasitoses chroniques provoquent anémie et malabsorption, donc retard de croissance).
\end{enumerate}
\emph{Autres facteurs acceptables :} qualité de l'eau et de l'assainissement (diarrhées répétées), travail physique précoce des enfants aux champs à Banka, niveau d'instruction des parents.

\medskip
\textbf{Question 5 — Synthèse argumentée (5 lignes max).}
\par
\og \textbf{Tous les humains partagent les mêmes \cle{caractères de l'espèce} (bipédie, main préhensile, système ABO...), ce qui affirme l'unité de l'espèce humaine. Mais chaque individu est unique : il reçoit de ses parents des \cle{caractères héréditaires} (groupe sanguin, couleur des yeux...) et acquiert au cours de sa vie des \cle{caractères modifiés par l'environnement} (taille, masse, langue, niveau d'études...). Nos différences viennent donc \emph{des deux} : l'hérédité fixe le potentiel, l'environnement détermine sa réalisation.} \fg
\end{exoresolu}

\begin{attention}[Erreurs fréquentes à éviter]
\begin{itemize}
    \item Dire \og le groupe sanguin est un caractère de l'espèce, donc tout le monde a le même groupe \fg : c'est le \emph{système} ABO qui est commun à tous, pas le groupe précis (A, B, AB ou O varient d'un individu à l'autre).
    \item Affirmer que la taille est \og purement héréditaire \fg : les Documents 2 et 3 montrent clairement l'effet de la nutrition et de la santé. La taille est un caractère à \textbf{déterminisme mixte}.
    \item Croire qu'un caractère acquis (musculature, langue parlée, niveau d'études) peut se transmettre aux enfants : \textbf{les caractères acquis ne sont pas héréditaires}.
    \item Comparer deux histogrammes en ne citant que les barres les plus hautes : il faut décrire la \textbf{position globale} et l'\textbf{étalement} des deux distributions.
\end{itemize}
\end{attention}

\courssub{Bilan de l'activité}
Cette activité d'intégration a permis de mettre en évidence, par l'analyse de données camerounaises, les trois piliers de la leçon :
\begin{enumerate}
    \item \textbf{L'unité de l'espèce humaine} est attestée par des caractères universels (système ABO, bipédie, main préhensile, organisation générale du corps) observés chez tous les individus du Document 1.
    \item \textbf{L'hérédité} se manifeste par la transmission fidèle de certains caractères (groupe sanguin, couleur des yeux, type de lobe d'oreille) des parents aux enfants, confirmée par l'identité de ces traits chez les vrais jumeaux (Document 2).
    \item \textbf{L'environnement} module l'expression de la majorité des caractères quantitatifs (taille, masse, IMC, tension artérielle) et détermine entièrement les caractères culturels (langue, niveau d'études), comme le prouvent les différences entre les jumeaux (Document 2) et le décalage des distributions de taille selon le milieu de vie (Document 3).
\end{enumerate}
\par
\textbf{Compétences évaluées (BEPC) :}
\begin{itemize}
    \item \textbf{Savoir :} définir et distinguer caractère de l'espèce, caractère héréditaire, caractère acquis.
    \item \textbf{Savoir-faire :} lire et interpréter un tableau de données familiales, un cas de jumeaux, un histogramme comparatif ; classer des caractères ; argumenter à partir de données chiffrées ; rédiger une conclusion synthétique et justifiée.
    \item \textbf{Savoir-être :} travail en groupe, écoute, esprit critique, rigueur dans l'argumentation, respect des délais.
\end{itemize}
\par
\emph{Application à la vie courante :} comprendre que la santé, la réussite scolaire et le développement physique ne sont pas \og écrits dans les gènes \fg{} seulement : l'amélioration des conditions de vie (nutrition, santé, éducation) permet à chaque enfant d'atteindre son plein potentiel. C'est le sens des politiques publiques de cantines scolaires, de vaccination, d'accès à l'eau potable et de scolarisation obligatoire au Cameroun.

\newpage

\coursec{Méthodes et exercices résolus}

\courssub{Méthodes et exercices résolus}

\begin{methode}[Identifier et classer les caractères : caractères d'espèce, héréditaires ou modifiés par l'environnement]
\textbf{Objectif :} À partir d'une liste de caractères observés chez l'homme, les ranger dans le bon tableau.
\textbf{Démarche en 4 étapes :}
\begin{enumerate}
    \item \textbf{Lister les caractères} proposés (ex. : nombre de membres, groupe sanguin, couleur de la peau, langue parlée, cicatrice, taille).
    \item \textbf{Tester l'universalité :} Ce caractère est-il présent chez \emph{tous} les êtres humains (sauf pathologie/accident) ? \textbf{OUI} $\rightarrow$ \cle{Caractère d'espèce} (ex. : 4 membres, 2 yeux, cerveau, respiration pulmonaire). \textbf{NON} $\rightarrow$ aller à l'étape 3.
    \item \textbf{Tester l'hérédité :} Ce caractère se transmet-il des parents aux enfants (ressemblance familiale) ? \textbf{OUI} $\rightarrow$ \cle{Caractère héréditaire} (ex. : groupe sanguin, couleur des yeux, type de lobe d'oreille, phénylcétonurie). \textbf{NON} $\rightarrow$ aller à l'étape 4.
    \item \textbf{Conclure :} Si le caractère n'est ni d'espèce, ni héréditaire, c'est un \cle{Caractère modifié par l'environnement} (acquis) (ex. : langue parlée, religion, cicatrice, musculature due au sport, bronzage, poids).
\end{enumerate}
\textbf{Réflexe :} Un caractère héréditaire est déterminé par des \cle{gènes} (ADN). Un caractère acquis n'est pas inscrit dans l'ADN des cellules germinales et ne se transmet pas à la descendance.
\end{methode}

\begin{exoresolu}[Classification de caractères observés dans une famille à Yaoundé]
On observe les caractères suivants chez les membres d'une famille : 
\begin{itemize}
    \item A : Possède un foie.
    \item B : Parle le français et l'ewondo.
    \item C : A le groupe sanguin A.
    \item D : A une cicatrice au genou gauche suite à une chute.
    \item E : A les yeux marron.
    \item F : Respire l'oxygène de l'air.
    \item G : Est droitier.
    \item H : A une forte musculature due à la pratique quotidienne du football.
\end{itemize}
\textbf{Classez chaque caractère dans l'un des trois tableaux ci-dessous et justifiez brièvement chaque classement.}
\tcblower
\textbf{Solution détaillée}

\textbf{Étape 1 : Analyse de chaque caractère.}
\begin{itemize}
    \item \textbf{A (Foie) :} Tous les humains en possèdent un (sauf anomalie majeure). $\rightarrow$ \textbf{Caractère d'espèce}.
    \item \textbf{B (Langues) :} S'apprend dans l'entourage familial/scolaire. Un bébé né de parents ewondo/francophones élevé au Japon parlera japonais. $\rightarrow$ \textbf{Caractère acquis (environnement)}.
    \item \textbf{C (Groupe sanguin A) :} Déterminé par des allèles sur le chromosome 9. Se transmet selon les lois de Mendel. $\rightarrow$ \textbf{Caractère héréditaire}.
    \item \textbf{D (Cicatrice) :} Résultat d'un traumatisme (chute). Ne modifie pas l'ADN des gamètes. $\rightarrow$ \textbf{Caractère acquis (environnement)}.
    \item \textbf{E (Yeux marron) :} Pigmentation de l'iris sous contrôle génétique (polygénique). Ressemblance parents/enfants fréquente. $\rightarrow$ \textbf{Caractère héréditaire}.
    \item \textbf{F (Respiration O$_2$) :} Métabolisme commun à tous les humains. $\rightarrow$ \textbf{Caractère d'espèce}.
    \item \textbf{G (Droitier) :} Latéralisation complexe (génétique + environnement/culture). Au programme de 3ème, on le classe souvent comme \textbf{caractère héréditaire} (prédominance génétique) mais attention : l'apprentissage forcé peut le modifier. \textit{Note : On l'accepte souvent en héréditaire au BEPC.}
    \item \textbf{H (Musculature football) :} Résultat de l'entraînement (hypertrophie musculaire). Disparaît à l'arrêt du sport. Non transmissible aux enfants. $\rightarrow$ \textbf{Caractère acquis (environnement)}.
\end{itemize}

\textbf{Étape 2 : Tableaux de classement.}

\begin{center}
\begin{tabularx}{\linewidth}{|l|X|}\hline
\rowcolor{popblueL} \textbf{Caractères d'espèce} & \textbf{Justification} \\ \hline
A : Possède un foie & Présent chez tout être humain viable (unité structurale). \\ 
F : Respire l'oxygène & Métabolisme énergétique universel chez l'homme. \\ \hline
\end{tabularx}
\end{center}

\begin{center}
\begin{tabularx}{\linewidth}{|l|X|}\hline
\rowcolor{popgreenL} \textbf{Caractères héréditaires} & \textbf{Justification} \\ \hline
C : Groupe sanguin A & Déterminé par des allèles (gènes), transmission parentale. \\ 
E : Yeux marron & Pigmentation génétique, ressemblance familiale. \\ 
G : Droitier & Prédisposition génétique forte (latéralisation). \\ \hline
\end{tabularx}
\end{center}

\begin{center}
\begin{tabularx}{\linewidth}{|l|X|}\hline
\rowcolor{poporangeL} \textbf{Caractères acquis (environnement)} & \textbf{Justification} \\ \hline
B : Langues parlées & Apprentissage culturel, non codé dans l'ADN. \\ 
D : Cicatrice & Accident de la vie, lésion tissulaire non génétique. \\ 
H : Musculature football & Adaptation phénotypique à l'effort (phénoplastie). \\ \hline
\end{tabularx}
\end{center}

\textbf{Conclusion :} L'espèce humaine se définit par des caractères constants (d'espèce). La diversité individuelle repose sur le mélange de caractères héréditaires (génotype) et de caractères acquis (influence du milieu).
\end{exoresolu}

\begin{methode}[Interpréter un document (tableau, graphique, photo) pour distinguer hérédité et environnement]
\textbf{Objectif :} Utiliser des données expérimentales ou observationnelles (jumeaux, adoption, population) pour trancher.
\textbf{Démarche en 4 étapes :}
\begin{enumerate}
    \item \textbf{Identifier le protocole :} Compare-t-on des \cle{jumeaux monozygotes (vrais jumeaux)} séparés à la naissance ? Des enfants adoptés vs parents biologiques/adoptifs ? Une population avant/après un changement de milieu ?
    \item \textbf{Repérer la variable étudiée :} Taille, QI, groupe sanguin, maladie, poids, etc.
    \item \textbf{Analyser les concordances/discordances :}
    \begin{itemize}
        \item \textbf{Forte concordance chez jumeaux monozygotes (même ADN) même élevés à part} $\rightarrow$ \textbf{Fort déterminisme génétique (héréditaire)}.
        \item \textbf{Concordance chez jumeaux dizygotes (frères/jumeaux faux) ou parents/enfants adoptifs partageant le milieu} $\rightarrow$ \textbf{Influence de l'environnement}.
        \item \textbf{Discordance chez jumeaux monozygotes élevés à part} $\rightarrow$ \textbf{Influence de l'environnement}.
    \end{itemize}
    \item \textbf{Rédiger la conclusion :} « Le caractère X est principalement héréditaire car... » OU « Le caractère Y est fortement influencé par l'environnement car... » OU « Le caractère Z résulte de l'interaction génotype $\times$ milieu ».
\end{enumerate}
\textbf{Formule clé :} \cle{Phénotype = Génotype + Environnement} (Interaction G $\times$ E).
\end{methode}

\begin{exoresolu}[Étude de la taille chez des jumeaux monozygotes séparés à la naissance]
On a mesuré la taille adulte de 10 paires de jumeaux monozygotes (vrais jumeaux) séparés à la naissance et élevés dans des familles différentes (milieux socio-économiques et nutritionnels variés). Résultats (en cm) :

\begin{center}
\begin{tabular}{|c|c|c|c|}\hline
\rowcolor{popblueL} \textbf{Paire} & \textbf{Jumeau 1 (Milieu favorisé)} & \textbf{Jumeau 2 (Milieu défavorisé)} & \textbf{Écart} \\ \hline
1 & 178 & 176 & 2 \\
2 & 165 & 163 & 2 \\
3 & 182 & 179 & 3 \\
4 & 170 & 168 & 2 \\
5 & 175 & 173 & 2 \\
6 & 160 & 158 & 2 \\
7 & 185 & 182 & 3 \\
8 & 172 & 170 & 2 \\
9 & 168 & 166 & 2 \\
10 & 177 & 175 & 2 \\ \hline
\end{tabular}
\end{center}

\textbf{1. Calculez la moyenne des tailles pour chaque groupe et l'écart moyen.}
\textbf{2. Interprétez ces résultats : la taille est-elle un caractère héréditaire, acquis, ou les deux ? Justifiez en utilisant la notion d'interaction G $\times$ E.}
\tcblower
\textbf{Solution détaillée}

\textbf{1. Calculs :}
\begin{itemize}
    \item Moyenne Jumeaux 1 (Favorisé) : $\frac{178+165+182+170+175+160+185+172+168+177}{10} = \frac{1732}{10} = \mathbf{\num{173,2}\,\text{cm}}$.
    \item Moyenne Jumeaux 2 (Défavorisé) : $\frac{176+163+179+168+173+158+182+170+166+175}{10} = \frac{1710}{10} = \mathbf{\num{171,0}\,\text{cm}}$.
    \item Écart moyen : $\frac{2+2+3+2+2+2+3+2+2+2}{10} = \frac{22}{10} = \mathbf{\num{2,2}\,\text{cm}}$.
\end{itemize}

\textbf{2. Interprétation et Conclusion :}
\begin{itemize}
    \item \textbf{Observation majeure :} Malgré des milieux de vie très différemment favorisés (nutrition, santé), la taille des jumeaux reste \textbf{très proche} (écart moyen seulement \num{2,2}\,\text{cm}, soit $\approx \num{1,3}\,\%$).
    \item \textbf{Argument génétique :} Les jumeaux monozygotes partagent \textbf{100\,\% de leur ADN} (même génotype). La forte concordance (ressemblance) de leur taille, \emph{même élevés à part}, prouve que la taille est \textbf{fortement déterminée par l'hérédité} (polygénique).
    \item \textbf{Argument environnemental :} On note cependant une \textbf{légère différence systématique} : les jumeaux du milieu favorisé sont en moyenne \num{2,2}\,\text{cm} plus grands. La nutrition, l'accès aux soins, l'absence de maladies chroniques pendant la croissance (milieu) ont permis une \textbf{meilleure expression du potentiel génétique}.
    \item \textbf{Synthèse :} La taille est un caractère à \textbf{déterminisme mixte (interaction Génotype $\times$ Environnement)}. Le génotype fixe une \cle{valeur cible} (norme de réaction), l'environnement (nutrition, santé) détermine si cette cible est atteinte.
\end{itemize}
\textbf{Phrase de conclusion type BEPC :} \og La taille chez l'homme est un caractère héréditaire à expression polygénique, dont la réalisation phénotypique finale dépend significativement des conditions environnementales (nutrition, santé), illustrant l'interaction Génotype $\times$ Milieu. \fg
\end{exoresolu}

\begin{methode}[Rédiger une argumentation scientifique structurée (CER : Claim - Evidence - Reasoning)]
\textbf{Objectif :} Répondre à une question de type « Montrez que... », « Justifiez que... », « Expliquez pourquoi... » avec rigueur.
\textbf{Structure en 3 parties (CER) :}
\begin{enumerate}
    \item \textbf{Claim (Affirmation/Thèse) :} Réponse directe à la question en une phrase. \textit{Ex : « Le groupe sanguin est un caractère strictement héréditaire. »}
    \item \textbf{Evidence (Preuves/Données) :} Citez \textbf{précisément} les faits du document (chiffres, observations, résultats d'expérience) ou les connaissances du cours. \textit{Ex : « Le document 1 montre que l'enfant a le groupe A alors que la mère est O et le père B... » ou « Le groupe sanguin est déterminé par le gène ABO sur le chromosome 9. »}
    \item \textbf{Reasoning (Raisonnement/Explication) :} Faites le lien logique entre la preuve et l'affirmation en mobilisant le mécanisme biologique. \textit{Ex : « Comme les allèles se séparent à la méiose et se recombinent à la fécondation, l'enfant reçoit un allèle de chaque parent. Le groupe sanguin ne change pas au cours de la vie et n'est pas influencé par l'alimentation. »}
\end{enumerate}
\textbf{Connecteurs logiques à utiliser :} \emph{Or, car, donc, par conséquent, en effet, cela s'explique par, ce qui prouve que, en revanche, cependant.}
\textbf{Interdiction :} « C'est évident », « Tout le monde sait que », « Normalement », « Probablement ».
\end{methode}

\begin{exoresolu}[Argumentation : Le lobe d'oreille attaché est-il un caractère héréditaire ?]
\textbf{Document :} Arbre généalogique simplifié d'une famille sur 3 générations.
\begin{itemize}
    \item Génération I : Père (lobes libres), Mère (lobes attachés).
    \item Génération II : 3 enfants : Fille 1 (lobes attachés), Fils 1 (lobes libres), Fille 2 (lobes attachés).
    \item Génération III : Fille 1 (lobes attachés) $\times$ Homme (lobes libres) $\rightarrow$ 2 enfants : lobes attachés et lobes libres.
\end{itemize}
\textbf{Question :} À l'aide de l'arbre généalogique, montrez que le caractère « lobe d'oreille attaché » est héréditaire et déduisez le mode de transmission (dominant/récessif).
\tcblower
\textbf{Solution détaillée (Modèle de rédaction CER)}

\textbf{Affirmation (Claim) :} Le caractère « lobe d'oreille attaché » est un \textbf{caractère héréditaire} à transmission \textbf{autosomique récessive}.

\textbf{Preuves (Evidence) :}
\begin{enumerate}
    \item \textbf{Transmission parentale :} Dans la génération I, le père a les lobes libres et la mère les lobes attachés. Ils ont 3 enfants dont 2 ont les lobes attachés (Fille 1, Fille 2) et 1 a les lobes libres (Fils 1). L'apparition du phénotype « attaché » chez les enfants alors que le père est « libre » prouve la transmission.
    \item \textbf{Saut de génération / Réapparition :} Le père (Gén I) est phénotypiquement « libre » mais a des enfants « attachés ». Il est donc \textbf{porteur sain} (hétérozygote).
    \item \textbf{Croisement test (Gén III) :} La Fille 1 (lobes attachés) mariée à un homme (lobes libres) donne 2 enfants : un « attaché » et un « libre ». Si « attaché » était dominant, une personne « attachée » (homozygote dominant ou hétérozygote) croisée avec un « libre » (récessif) ne pourrait pas donner de « libre » si elle était homozygote, ou donnerait 50\,\% « libre » si hétérozygote. Mais ici, le parent « attaché » donne un enfant « libre » $\rightarrow$ le parent « attaché » est homozygote pour l'allèle récessif. Le parent « libre » est donc hétérozygote.
\end{enumerate}

\textbf{Raisonnement (Reasoning) :}
\begin{itemize}
    \item Notons $L$ l'allèle « lobes libres » et $l$ l'allèle « lobes attachés ».
    \item Le père (Gén I) est phénotype « libre » mais a des enfants « attachés » ($ll$) $\rightarrow$ Son génotype est \textbf{$Ll$ (hétérozygote)}.
    \item La mère (Gén I) est phénotype « attaché » $\rightarrow$ Son génotype est \textbf{$ll$ (homozygote récessif)}.
    \item Croisement $Ll \times ll$ $\rightarrow$ 50\,\% $Ll$ (libres), 50\,\% $ll$ (attachés). \textbf{Correspond aux effectifs observés (1 libre, 2 attachés $\approx$ ratio 1:1)}.
    \item La Fille 1 est $ll$. Son mari est phénotype « libre » mais a un enfant « attaché » ($ll$) $\rightarrow$ Le mari est \textbf{$Ll$}.
    \item Croisement $ll \times Ll$ $\rightarrow$ 50\,\% $Ll$ (libres), 50\,\% $ll$ (attachés). \textbf{Correspond aux effectifs observés (1 libre, 1 attaché)}.
\end{itemize}

\textbf{Conclusion :} L'analyse des croisements dans la fratrie et la descendance confirme que l'allèle « lobe attaché » ($l$) est \textbf{récessif} par rapport à l'allèle « lobe libre » ($L$). Le caractère suit une transmission \textbf{autosomique récessive} (non liée au sexe, touche les deux sexes équitablement).
\end{exoresolu}

\begin{methode}[Utiliser les groupes sanguins (système ABO) pour exclure une paternité (Preuve par l'impossible)]
\textbf{Objectif :} Déterminer si un homme peut être le père biologique d'un enfant connaissant les groupes sanguins de la mère, de l'enfant et de l'homme mis en cause.
\textbf{Rappel génétique (Programme 3ème) :} 3 allèles : $I^A$, $I^B$, $i$ ($I^A$ et $I^B$ codominants, dominants sur $i$).
\begin{center}
\begin{tabular}{|c|c|}\hline
\rowcolor{popblueL} \textbf{Phénotype (Groupe)} & \textbf{Génotypes possibles} \\ \hline
A & $I^A I^A$ ou $I^A i$ \\ \hline
B & $I^B I^B$ ou $I^B i$ \\ \hline
AB & $I^A I^B$ \\ \hline
O & $ii$ \\ \hline
\end{tabular}
\end{center}
\textbf{Démarche en 3 étapes (Méthode de l'exclusion) :}
\begin{enumerate}
    \item \textbf{Déduire les allèles obligatoires de l'enfant :} Regarder le groupe de la mère. L'enfant a reçu \textbf{un allèle de la mère}. L'autre allèle \textbf{vient obligatoirement du père biologique}.
    \item \textbf{Identifier l'allèle paternel obligatoire :} Enlever la contribution maternelle du génotype de l'enfant. L'allèle restant est \textbf{l'allèle que le père \emph{doit} posséder}.
    \item \textbf{Tester l'homme mis en cause :} 
    \begin{itemize}
        \item S'il \textbf{ne possède pas} cet allèle $\rightarrow$ \textbf{EXCLUSION CERTAINE} (Il n'est pas le père).
        \item S'il \textbf{possède} cet allèle $\rightarrow$ \textbf{NON-EXCLUSION} (Il \emph{peut} être le père, mais ce n'est pas une preuve formelle ; seul l'ADN le prouve).
    \end{itemize}
\end{enumerate}
\textbf{Attention :} On ne prouve JAMAIS la paternité par les groupes sanguins, on ne fait que \textbf{ne pas l'exclure}.
\end{methode}

\begin{exoresolu}[Exclusion de paternité par les groupes sanguins ABO]
Une mère de groupe sanguin \textbf{A} a un enfant de groupe sanguin \textbf{AB}. Deux hommes prétendent être le père :
\begin{itemize}
    \item Homme 1 : Groupe \textbf{B}.
    \item Homme 2 : Groupe \textbf{O}.
\end{itemize}
\textbf{Lequel des deux peut être exclu ? Justifiez rigoureusement.}
\tcblower
\textbf{Solution détaillée}

\textbf{Étape 1 : Génotypes possibles de la mère (Groupe A).}
$I^A I^A$ ou $I^A i$.

\textbf{Étape 2 : Génotype de l'enfant (Groupe AB).}
Un seul génotype possible : \textbf{$I^A I^B$} (codominance).
L'enfant possède \textbf{obligatoirement} un allèle $I^A$ et un allèle $I^B$.

\textbf{Étape 3 : Déterminer l'allèle maternel transmis.}
La mère (Groupe A) ne possède \textbf{pas} l'allèle $I^B$ (ses génotypes possibles sont $I^A I^A$ ou $I^A i$).
$\rightarrow$ L'allèle $I^A$ de l'enfant \emph{peut} venir de la mère.
$\rightarrow$ L'allèle $I^B$ de l'enfant \textbf{ne peut PAS venir de la mère}.
$\rightarrow$ \textbf{L'allèle $I^B$ a OBLIGATOIREMENT été transmis par le père biologique.}

\textbf{Étape 4 : Test des hommes mis en cause.}
\begin{itemize}
    \item \textbf{Homme 1 (Groupe B) :} Génotypes possibles : $I^B I^B$ ou $I^B i$. \textbf{Il possède l'allèle $I^B$.} $\rightarrow$ \textbf{NON-EXCLUSION.} Il peut être le père.
    \item \textbf{Homme 2 (Groupe O) :} Génotype unique : \textbf{$ii$}. Il ne possède \textbf{PAS} l'allèle $I^B$ (il n'a que des allèles $i$). $\rightarrow$ \textbf{EXCLUSION CERTAINE.} Il ne peut pas transmettre $I^B$. Il n'est pas le père biologique.
\end{itemize}

\textbf{Conclusion :} \textbf{L'Homme 2 (Groupe O) est exclu de la paternité.} L'Homme 1 (Groupe B) n'est pas exclu (compatibilité génétique), mais seul un test ADN (empreinte génétique) pourrait confirmer la paternité.
\end{exoresolu}

\begin{methode}[Appliquer les connaissances à la vie courante : Conseil génétique et dépistage néonatal (Contexte Cameroun)]
\textbf{Objectif :} Expliquer l'utilité du dépistage néonatal (ex. drépanocytose, phénylcétonurie) et du conseil génétique avant mariage (carte de génotype au Cameroun).
\textbf{Démarche en 3 étapes :}
\begin{enumerate}
    \item \textbf{Identifier la maladie génétique concernée :} Transmission autosomique récessive (allèle muté $m$, allèle sain $M$). Malades = $mm$. Porteurs sains (asymptomatiques) = $Mm$.
    \item \textbf{Calculer le risque pour la descendance :} Croisement de deux porteurs sains ($Mm \times Mm$).
    \begin{center}
    \begin{tabular}{|c|c|c|}\hline
    \rowcolor{popblueL} & $M$ & $m$ \\ \hline
    $M$ & $MM$ (Sain) & $Mm$ (Porteur) \\ \hline
    $m$ & $Mm$ (Porteur) & $\mathbf{mm}$ (\textbf{Malade}) \\ \hline
    \end{tabular}
    \end{center}
    $\rightarrow$ Risque \textbf{1/4 (25\,\%)} d'enfant malade à chaque grossesse. Risque 1/2 (50\,\%) porteur sain. Risque 1/4 (25\,\%) sain non porteur.
    \item \textbf{Formuler la recommandation (Savoir-être) :}
    \begin{itemize}
        \item \textbf{Avant mariage :} Connaître son \cle{statut génotypique} (carte AA, AS, SS, AC, SC, CC pour drépanocytose). Éviter le mariage à risque (ex. AS $\times$ AS, AS $\times$ SS, SS $\times$ SS) ou se préparer (diagnostic prénatal, PMA).
        \item \textbf{À la naissance :} Dépistage néonatal (test de Guthrie au talon) $\rightarrow$ Prise en charge précoce (pénicilline, vaccins, nutrition) $\rightarrow$ Réduction mortalité/invalidité.
    \end{itemize}
\end{enumerate}
\textbf{Vocabulaire clé :} \cle{Porteur sain (hétérozygote)}, \cle{Homozygote malade}, \cle{Conseil génétique}, \cle{Dépistage néonatal}.
\end{methode}

\begin{exoresolu}[Conseil génétique pour la drépanocytose au Cameroun]
Au Cameroun, la drépanocytose (anémie falciforme) touche environ \num{2}\,\% des naissances (homozygotes $SS$). Le trait drépanocytaire (hétérozygotes $AS$) concerne environ \num{25}\,\% de la population. Un couple consulte avant mariage. L'homme est $AS$, la femme est $AS$.
\textbf{1. Représentez le croisement génétique et calculez les probabilités de génotypes et phénotypes pour leurs futurs enfants.}
\textbf{2. En tant que professionnel de santé, rédigez le conseil que vous donnez à ce couple (ton bienveillant, vocabulaire précis, respect du choix).}
\tcblower
\textbf{Solution détaillée}

\textbf{1. Croisement et Probabilités :}
\textbf{Gamètes :} Homme $AS \rightarrow$ Gamètes $A$ ou $S$. Femme $AS \rightarrow$ Gamètes $A$ ou $S$.

\begin{center}
\begin{tabular}{|c|c|c|}\hline
\rowcolor{popblueL} \textbf{Fécondation} & \textbf{$A$ (Mère)} & \textbf{$S$ (Mère)} \\ \hline
\textbf{$A$ (Père)} & $AA$ & $AS$ \\ \hline
\textbf{$S$ (Père)} & $AS$ & $\mathbf{SS}$ \\ \hline
\end{tabular}
\end{center}

\textbf{Résultats (à chaque grossesse, indépendamment des précédentes) :}
\begin{itemize}
    \item \textbf{$AA$ (Homozygote normal) :} $\mathbf{\num{25}\,\%}$ ($1/4$). Enfant sain, non porteur.
    \item \textbf{$AS$ (Hétérozygote / Trait drépanocytaire) :} $\mathbf{\num{50}\,\%}$ ($2/4$). Enfant \textbf{porteur sain}, asymptomatique en général, transmet l'allèle $S$.
    \item \textbf{$SS$ (Homozygote malade / Drépanocytose majeure) :} $\mathbf{\num{25}\,\%}$ ($1/4$). Enfant \textbf{malade} : anémie hémolytique chronique, crises vaso-occlusives douloureuses, infections répétées, risque vital, espérance de vie réduite sans prise en charge lourde.
\end{itemize}

\textbf{2. Rédaction du conseil génétique (Modèle) :}

\og \textbf{Monsieur, Madame,}

\textbf{Vos résultats :} Vous êtes tous les deux porteurs sains du trait drépanocytaire (génotype $AS$). Vous êtes en bonne santé.

\textbf{Le risque pour vos enfants :} À chaque grossesse, il y a \textbf{1 chance sur 4 (\num{25}\,\%)} d'avoir un enfant atteint de drépanocytose majeure ($SS$), \textbf{2 chances sur 4 (\num{50}\,\%)} d'avoir un enfant porteur sain comme vous ($AS$), et \textbf{1 chance sur 4 (\num{25}\,\%)} d'avoir un enfant non porteur ($AA$). Ce risque est le \textbf{même pour chaque grossesse}, quel que soit le phénotype des enfants précédents.

\textbf{Options et accompagnement :}
\begin{itemize}
    \item \textbf{Diagnostic Prénatal (DPN) :} Possible dès le 1er trimestre (biopsie trophoblastique) ou 2ème trimestre (amniocentèse) pour connaître le génotype du fœtus. En cas de fœtus $SS$, l'interruption médicale de grossesse (IMG) est encadrée par la loi camerounaise (Loi n°90/035) si la pathologie est grave.
    \item \textbf{Dépistage Néonatal Systématique :} Si vous choisissez de ne pas faire de DPN, le test au talon (Guthrie) à la naissance (gratuit dans les hôpitaux publics partenaires du Programme National de Lutte contre la Drépanocytose) permet une prise en charge \textbf{immédiate} (prophylaxie pénicilline, vaccins pneumocoque/Haemophilus, acide folique, éducation thérapeutique) qui \textbf{transforme le pronostic vital}.
    \item \textbf{Fécondation In Vitro (FIV) avec DPI :} Diagnostic Préimplantatoire (sélection d'embryons $AA$ ou $AS$) possible dans des centres spécialisés (coût élevé).
\end{itemize}

\textbf{Notre rôle :} Vous informer sans vous influencer. Votre choix vous appartient, en toute liberté de conscience. Nous restons à votre disposition pour tout complément d'information ou soutien psychologique. \fg

\textbf{Conclusion :} La connaissance du statut génotypique ($AS/AS$) permet une \textbf{parentalité responsable} et une \textbf{prise en charge précoce} qui sauve des vies.
\end{exoresolu}

\begin{methode}[Analyser l'influence de l'environnement sur un caractère à déterminisme génétique (Norme de réaction)]
\textbf{Objectif :} Expliquer pourquoi des individus de même génotype (clones, jumeaux monozygotes, lignées pures plantes) peuvent avoir des phénotypes différents selon le milieu.
\textbf{Démarche en 3 étapes :}
\begin{enumerate}
    \item \textbf{Définir la norme de réaction :} Intervalle de valeurs phénotypiques qu'un \textbf{même génotype} peut prendre selon les conditions du milieu.
    \item \textbf{Analyser le document (courbe/histogramme) :} Tracer ou lire la courbe \textbf{Phénotype = f(Milieu)} pour un génotype donné.
    \begin{itemize}
        \item Courbe plate $\rightarrow$ Caractère \textbf{canalisé} (peu sensible au milieu, ex. groupe sanguin, nombre de doigts).
        \item Courbe pentue/large $\rightarrow$ Caractère \textbf{plastique} (très sensible au milieu, ex. taille, poids, QI, rendement agricole, couleur hortensia).
    \end{itemize}
    \item \textbf{Conclure sur l'interaction G $\times$ E :} « Ce caractère a une base génétique forte (héritabilité élevée) MAIS son expression finale dépend de [facteur environnemental précis : nutrition, lumière, température, éducation, pathogènes]. »
\end{enumerate}
\textbf{Exemple type BEPC :} Hauteur de plantes de même lignée pure cultivées avec/sans engrais / à l'ombre/au soleil. Taille de jumeaux monozygotes élevés dans des milieux différents.
\end{methode}

\begin{exoresolu}[Norme de réaction : Rendement de deux variétés de maïs selon l'apport en azote]
On cultive deux variétés de maïs (Variété A : locale ; Variété B : améliorée) dans des parcelles recevant des doses croissantes d'engrais azoté (0, 50, 100, 150 kg N/ha). Le rendement (en t/ha) est mesuré.

\begin{center}
\begin{tabular}{|c|c|c|}\hline
\rowcolor{popblueL} \textbf{Azote (kg/ha)} & \textbf{Variété A (Locale)} & \textbf{Variété B (Améliorée)} \\ \hline
0 & 1,2 & 1,0 \\ \hline
50 & 2,0 & 3,5 \\ \hline
100 & 2,5 & 6,0 \\ \hline
150 & 2,6 & 7,2 \\ \hline
\end{tabular}
\end{center}

\textbf{1. Représentez graphiquement (schéma descriptif) les normes de réaction des deux variétés.}
\textbf{2. Comparez la plasticité phénotypique des deux variétés.}
\textbf{3. Quelle variété conseiller à un paysan n'ayant pas accès à l'engrais ? Et à un paysan pouvant fertiliser à 150 kg N/ha ? Justifiez.}
\tcblower
\textbf{Solution détaillée}

\textbf{1. Représentation graphique (Description pour le BEPC - tracer sur papier millimétré) :}
\begin{itemize}
    \item Axes : Abscisse = Dose Azote (0 à 150 kg/ha). Ordonnée = Rendement (0 à 8 t/ha).
    \item \textbf{Courbe Variété A (Locale) :} Parte de (0 ; 1,2). Monte vite vers (50 ; 2,0). \textbf{Plateau rapide} vers (100 ; 2,5) et (150 ; 2,6). Forme \textbf{asymptotique / saturante}.
    \item \textbf{Courbe Variété B (Améliorée) :} Parte de (0 ; 1,0) (légèrement plus bas que A sans engrais). Monte \textbf{régulièrement et fortement} : (50 ; 3,5), (100 ; 6,0), (150 ; 7,2). Forme \textbf{quasi-linéaire croissante} sur l'intervalle.
\end{itemize}

\textbf{2. Comparaison de la plasticité :}
\begin{itemize}
    \item \textbf{Variété A : Faible plasticité (Norme de réaction étroite).} Son rendement varie peu (1,2 $\rightarrow$ 2,6 t/ha, amplitude 1,4 t/ha) malgré un triplement de l'azote. Elle est \textbf{peu sensible} à la richesse du milieu. C'est une variété \textbf{rustique}, adaptée aux milieux pauvres.
    \item \textbf{Variété B : Forte plasticité (Norme de réaction large).} Son rendement varie énormément (1,0 $\rightarrow$ 7,2 t/ha, amplitude 6,2 t/ha). Elle est \textbf{très sensible} à la fertilisation. C'est une variété \textbf{à haut potentiel} mais \textbf{exigeante} en intrants.
\end{itemize}

\textbf{3. Conseil agronomique (Application vie courante - Cameroun) :}
\begin{itemize}
    \item \textbf{Paysan SANS engrais (0 kg N/ha) :} Conseiller \textbf{Variété A (Locale)}. Rendement : 1,2 t/ha contre 1,0 t/ha pour B. La variété locale valorise mieux la fertilité naturelle du sol (racines profondes, association mycorhizienne, cycle court). \textbf{Argument :} « Sans intrant, la variété améliorée déprime (besoin azoté non satisfait), la variété locale assure la sécurité alimentaire de base. »
    \item \textbf{Paysan AVEC engrais (150 kg N/ha) :} Conseiller \textbf{Variété B (Améliorée)}. Rendement : 7,2 t/ha contre 2,6 t/ha pour A. La variété améliorée convertit efficacement l'azote en biomasse/grain (indice de récolte élevé, port dressé). \textbf{Argument :} « L'investissement en engrais (coût $\approx$ 150 kg $\times$ 300 FCFA/kg = 45 000 FCFA/ha) n'est rentable (revenu maïs $\approx$ 7,2 t $\times$ 150 FCFA/kg = 1 080 000 FCFA) qu'avec la variété B. La variété A sature et ne valorise pas l'engrais. »
\end{itemize}
\textbf{Conclusion :} Le choix de la variété (Génotype) doit s'adapter au milieu (Environnement = niveau de fertilisation) pour optimiser le Phénotype (Rendement). C'est l'application concrète de l'interaction \cle{G $\times$ E} en agriculture.
\end{exoresolu}

\begin{attention}[Les 5 erreurs qui coûtent des points au BEPC]
\begin{enumerate}
    \item \textbf{Confondre « Caractère d'espèce » et « Caractère héréditaire ».} \textit{Erreur :} « Avoir 2 yeux est un caractère héréditaire. » \textbf{Correction :} C'est un caractère d'espèce (universel). Héréditaire = variable entre individus (ex. couleur des yeux) et transmissible.
    \item \textbf{Dire qu'un caractère acquis se transmet aux enfants (Lamarckisme).} \textit{Erreur :} « Le forgeron a des bras musclés, ses enfants naîtront musclés. » / « La cicatrice du père se transmet au fils. » \textbf{Correction :} Seuls les caractères codés dans l'ADN des \textbf{gamètes} (cellules reproductrices) se transmettent. Les modifications du soma (corps) ne modifient pas le génotype germinal.
    \item \textbf{Affirmer que les groupes sanguins \emph{prouvent} la paternité.} \textit{Erreur :} « L'homme est groupe A, l'enfant est groupe A, donc c'est le père. » \textbf{Correction :} Les groupes sanguins ne permettent que l'\textbf{EXCLUSION} (non-exclusion $\neq$ preuve). Seule l'empreinte génétique (ADN) prouve la paternité (probabilité $> 99,99\,\%$).
    \item \textbf{Oublier l'interaction G $\times$ E pour les caractères quantitatifs (taille, poids, QI, rendement).} \textit{Erreur :} « La taille est 100\,\% génétique » ou « La taille est 100\,\% alimentation ». \textbf{Correction :} Dire « La taille est un caractère à déterminisme \textbf{polygénique} dont l'expression phénotypique dépend de la \textbf{nutrition et de la santé} pendant la croissance (norme de réaction). »
    \item \textbf{Mauvaise lecture d'arbre généalogique : Ne pas tester l'hypothèse récessive/dominante sur TOUS les croisements.} \textit{Erreur :} Conclure « dominant » parce qu'un parent malade a un enfant malade, sans vérifier si deux parents sains peuvent avoir un enfant malade (ce qui prouverait le récessif). \textbf{Réflexe :} Tester systématiquement les deux hypothèses (Dominante / Récessive) sur \textbf{tous} les couples parents/enfants de l'arbre. Une seule exception invalide l'hypothèse.
\end{enumerate}
\end{attention}

\newpage

\coursec{Exercices}

\courssub{Je vérifie mes connaissances}

\begin{exercice}[QCM : Caractères de l'espèce]
Parmi les propositions suivantes, laquelle correspond à un \textbf{caractère de l'espèce humaine} (commun à tous les individus sains) ?
\begin{enumerate}
    \item Avoir les yeux marron.
    \item Posséder un squelette interne avec une colonne vertébrale.
    \item Parler la langue française.
    \item Avoir une cicatrice au genou.
\end{enumerate}
\end{exercice}

\begin{exercice}[Vrai ou Faux justifié]
Indiquez si les affirmations suivantes sont vraies (V) ou fausses (F) en justifiant brièvement votre réponse.
\begin{enumerate}
    \item Deux individus de l'espèce humaine peuvent avoir des groupes sanguins différents.
    \item La taille d'un adulte dépend uniquement de son patrimoine génétique.
    \item Un caractère acquis se transmet obligatoirement à la descendance.
    \item Le critère de fécondité entre individus permet de définir l'appartenance à une même espèce.
\end{enumerate}
\end{exercice}

\begin{exercice}[Définitions et vocabulaire]
Donnez la définition précise des termes suivants en utilisant le vocabulaire scientifique de la leçon :
\begin{enumerate}
    \item Caractère héréditaire.
    \item Caractère acquis (ou modifié par l'environnement).
    \item Espèce biologique.
    \item Variabilité intra-spécifique.
\end{enumerate}
\end{exercice}

\begin{exercice}[Calcul direct : Fréquence phénotypique]
Dans une classe de 40 élèves de 3ème, on relève les phénotypes suivants pour le caractère \og lobe de l'oreille attaché / libre \fg : 28 élèves ont les lobes libres (allèle dominant) et 12 élèves ont les lobes attachés (allèle récessif).
Calculez la fréquence (en \%) de chaque phénotype dans cet échantillon.
\end{exercice}

\courssub{Je m'entraîne}

\begin{exercice}[Classement de caractères dans une famille]
Le tableau ci-dessous présente plusieurs caractères observés chez trois membres d'une même famille (Père, Mère, Enfant).

\begin{center}
\begin{tabularx}{\linewidth}{|l|X|X|X|}\hline
\textbf{Caractère} & \textbf{Père} & \textbf{Mère} & \textbf{Enfant} \\ \hline
Groupe sanguin (ABO) & A & B & O \\ \hline
Taille (cm) & 1,78 & 1,65 & 1,70 \\ \hline
Lobe de l'oreille & Libre & Attaché & Libre \\ \hline
Cicatrice au bras & Oui (accident vélo) & Non & Non \\ \hline
Couleur des yeux & Marron & Bleu & Marron \\ \hline
Langue parlée & Français, Bassa & Français, Anglais & Français \\ \hline
\end{tabularx}
\end{center}

\begin{enumerate}
    \item Pour chaque caractère du tableau, précisez s'il est \textbf{héréditaire} ou \textbf{acquis}.
    \item Justifiez votre classement pour le groupe sanguin, la taille, la cicatrice et la langue parlée.
\end{enumerate}
\end{exercice}

\begin{exercice}[Exploitation d'histogramme : Variabilité de la taille]
L'histogramme ci-dessous représente la répartition des tailles (en cm) de 40 élèves d'une classe de 3ème au Cameroun.

\begin{popfigure}
\begin{tikzpicture}
\begin{axis}[
    ybar,
    width=\linewidth,
    height=8cm,
    xlabel={Taille (cm)},
    ylabel={Effectif},
    xtick={145,150,155,160,165,170,175,180,185},
    ytick={0,2,4,6,8,10},
    ymin=0, ymax=11,
    bar width=12pt,
    tick label style={font=\small},
    label style={font=\small},
]
\addplot[fill=popblue!70, draw=popblue] coordinates {
    (145,1) (150,2) (155,5) (160,8) (165,10) (170,8) (175,4) (180,2)
};
\end{axis}
\end{tikzpicture}
\legende{Répartition des tailles de 40 élèves de 3ème.}
\end{popfigure}

\begin{enumerate}
    \item Déterminez la classe de taille la plus fréquente (classe modale).
    \item Calculez la taille moyenne approximative de la classe (utilisez les valeurs centrales des classes : 147,5 ; 152,5 ; 157,5 ; 162,5 ; 167,5 ; 172,5 ; 177,5 ; 182,5).
    \item Quelles sont les valeurs extrêmes (min et max) observées ?
    \item Proposez \textbf{deux explications scientifiques} (l'une liée à l'hérédité, l'autre à l'environnement) pour justifier cette variabilité de la taille au sein de la classe.
\end{enumerate}
\end{exercice}

\begin{exercice}[Étude de cas : Vrais jumeaux élevés séparément]
Deux vrais jumeaux monozygotes (issus du même ovule fécondé) ont été adoptés à la naissance par deux familles différentes vivant dans des environnements distincts (l'un en zone urbaine à Douala, l'autre en zone rurale dans le Nord). À 20 ans, on observe :
\begin{itemize}
    \item Même groupe sanguin (A Rh+), même couleur des yeux (marron), même empreintes digitales.
    \item Masse corporelle : Jumeau 1 = 72 kg ; Jumeau 2 = 58 kg.
    \item Métier : Jumeau 1 = Étudiant en informatique ; Jumeau 2 = Agriculteur.
    \item Langue courante : Jumeau 1 = Français, Anglais ; Jumeau 2 = Français, Fulfulde.
\end{itemize}

\begin{enumerate}
    \item Parmi les caractères listés, classez-les en \og Héréditaires \fg et \og Acquis / Influencés par l'environnement \fg.
    \item Expliquez pourquoi les vrais jumeaux ont \textbf{obligatoirement} le même groupe sanguin et la même couleur des yeux.
    \item Pourquoi leurs masses corporelles et leurs métiers sont-ils différents malgré un patrimoine génétique identique ?
    \item Rédigez une conclusion synthétique (3 à 4 lignes) sur l'origine des différences entre individus de la même espèce.
\end{enumerate}
\end{exercice}

\begin{exercice}[Critère de l'espèce et fécondité]
Un élève affirme : \og Un Pygmée Baka de la forêt du Sud-Cameroun et un Peul du Nord-Cameroun n'appartiennent pas à la même espèce car ils sont très différents physiquement et culturellement. \fg

\begin{enumerate}
    \item En vous appuyant sur le \textbf{critère de fécondité} (critère biologique de l'espèce), expliquez pourquoi cet élève a tort.
    \item Citez \textbf{trois caractères de l'espèce humaine} qui prouvent leur unité biologique fondamentale.
    \item Donnez un exemple de caractère qui varie entre ces deux populations (variabilité intra-spécifique) et précisez s'il est principalement héréditaire ou acquis.
\end{enumerate}
\end{exercice}

\courssub{Je me prépare au BEPC}

\begin{exercice}[Sujet type BEPC : Hérédité du groupe sanguin et diversité]
\textbf{Contexte :} Un couple consulte un centre de santé pour un conseil prénatal. Le père est de groupe sanguin A, la mère de groupe B. Leur premier enfant est de groupe O.

\textbf{Données :} Le système ABO est déterminé par trois allèles : $I^A$, $I^B$ (codominants) et $i$ (récessif).
\begin{enumerate}
    \item Déduisez les \textbf{génotypes} probables du père et de la mère pour qu'ils puissent avoir un enfant de groupe O. Justifiez.
    \item En vous basant sur cet exemple, expliquez pourquoi le groupe sanguin est considéré comme un \textbf{caractère héréditaire}.
    \item Citez \textbf{deux autres exemples de caractères héréditaires} chez l'homme (autres que le groupe sanguin) et \textbf{deux exemples de caractères acquis}.
    \item L'enfant grandit dans un environnement où l'alimentation est déséquilibrée (carences en protéines et vitamines). Précisez quels caractères de l'enfant (parmi : groupe sanguin, taille adulte, couleur de peau, langue parlée, intelligence scolaire) risquent d'être \textbf{modifiés par l'environnement} et lesquels resteront \textbf{inchangés}.
\end{enumerate}
\end{exercice}

\begin{exercice}[Sujet type BEPC : Unité biologique et lutte contre les discriminations]
\textbf{Situation :} Dans un village du Cameroun, une rumeur se propage : \og Les personnes albinos ne sont pas tout à fait humaines, elles appartiennent à une autre espèce, il faut les éviter. \fg Cette rumeur entraîne l'exclusion sociale et scolaire d'enfants albinos.

En tant qu'élève de 3ème instruit en SVTEEHB, vous êtes chargé(e) de rédiger un \textbf{texte argumenté} (entre 8 et 10 lignes) à destination du chef de village et des parents d'élèves pour combattre cette discrimination.

Votre texte doit obligatoirement s'appuyer sur les connaissances scientifiques suivantes :
\begin{itemize}
    \item La définition de l'espèce humaine (critère de fécondité).
    \item La distinction entre caractères de l'espèce (unitaires) et caractères variables (diversité).
    \item L'origine génétique de l'albinisme (caractère héréditaire récessif) et son absence d'impact sur la fécondité.
    \item La conclusion sur l'unité biologique de l'espèce humaine.
\end{itemize}
\end{exercice}

\begin{exercice}[Sujet type BEPC : Synthèse - Hérédité, Environnement, Unité]
Le document ci-dessous résume les résultats d'une étude comparative menée sur 100 paires de vrais jumeaux (monozygotes) et 100 paires de faux jumeaux (dizygotes) élevés ensemble dans les mêmes familles à Yaoundé. On a mesuré la \textbf{concordance} (pourcentage de paires où les deux jumeaux présentent le même phénotype) pour trois caractères.

\begin{center}
\begin{tabularx}{\linewidth}{|l|X|X|}\hline
\textbf{Caractère étudié} & \textbf{Concordance Vrais Jumeaux (Monozygotes)} & \textbf{Concordance Faux Jumeaux (Dizygotes)} \\ \hline
Groupe sanguin ABO & 100 \% & 50 \% \\ \hline
Myopie (port de lunettes) & 85 \% & 40 \% \\ \hline
Langue maternelle parlée & 100 \% & 100 \% \\ \hline
\end{tabularx}
\end{center}

\begin{enumerate}
    \item \textbf{Analyse des résultats :}
    \begin{enumerate}
        \item Comparez les concordances pour le groupe sanguin entre vrais et faux jumeaux. Quelle conclusion tirez-vous sur l'origine de ce caractère ?
        \item Pour la myopie, la concordance n'est pas de 100 \% chez les vrais jumeaux. Que cela suggère-t-il sur les facteurs intervenant dans ce caractère ?
        \item Pourquoi la concordance est-elle de 100 \% pour la langue maternelle chez les deux types de jumeaux ? Quelle est la nature de ce caractère ?
    \end{enumerate}
    \item \textbf{Modélisation :}
    \item À partir de l'ensemble des résultats, classez les trois caractères dans le tableau suivant en cochant la case correspondante :
    \begin{center}
    \begin{tabular}{|l|c|c|}\hline
    \textbf{Caractère} & \textbf{Principalement Héréditaire} & \textbf{Principalement Acquis / Environnement} \\ \hline
    Groupe sanguin & & \\ \hline
    Myopie & & \\ \hline
    Langue maternelle & & \\ \hline
    \end{tabular}
    \end{center}
    \item \textbf{Argumentation :} Un journaliste écrit : \og La myopie est une maladie génétique, on l'hérite de ses parents, point final. \fg En vous appuyant \textbf{uniquement} sur les données du tableau, critiquez cette affirmation trop simpliste.
    \item \textbf{Bilan :} Rappellez la définition d'un \textbf{caractère de l'espèce humaine} et donnez un exemple qui ne figure pas dans le tableau.
\end{enumerate}
\end{exercice}

\newpage

\coursec{Corrigés des exercices}

\courssub{Je vérifie mes connaissances}

\begin{corrige}[Exercice 1 — QCM : Caractères de l'espèce]
\textbf{Bonne réponse : proposition 2.}

\begin{itemize}
    \item \textbf{Proposition 1 (yeux marron) : FAUX.} La couleur des yeux varie d'un individu à l'autre (marron, bleu, vert...) : c'est un \cle{caractère variable} (héréditaire), pas un caractère de l'espèce.
    \item \textbf{Proposition 2 (squelette interne avec colonne vertébrale) : VRAI.} Tous les individus sains de l'espèce humaine possèdent un squelette interne et une colonne vertébrale : c'est un \cle{caractère de l'espèce humaine}, commun à tous.
    \item \textbf{Proposition 3 (parler français) : FAUX.} La langue parlée est un \cle{caractère acquis} (appris dans l'environnement) ; tous les humains ne parlent pas français.
    \item \textbf{Proposition 4 (cicatrice au genou) : FAUX.} Une cicatrice est un caractère acquis, propre à un individu, non commun à tous.
\end{itemize}

\textbf{Point de vigilance :} un caractère de l'espèce est présent chez \emph{tous} les individus sains ; un caractère variable distingue les individus entre eux.
\end{corrige}

\begin{corrige}[Exercice 2 — Vrai ou Faux justifié]
\begin{enumerate}
    \item \textbf{VRAI.} Le groupe sanguin est un caractère héréditaire \emph{variable} : deux individus peuvent être de groupes différents (A, B, AB ou O) tout en appartenant à la même espèce.
    \item \textbf{FAUX.} La taille dépend du patrimoine génétique \emph{et} de l'environnement (alimentation, santé pendant la croissance). C'est un caractère mixte (héréditaire et modifié par l'environnement).
    \item \textbf{FAUX.} Un caractère acquis (cicatrice, musculation, langue parlée) ne modifie pas l'information génétique : il \emph{ne se transmet pas} à la descendance.
    \item \textbf{VRAI.} Deux individus appartiennent à la même espèce s'ils peuvent se reproduire entre eux et donner des descendants \emph{féconds} : c'est le \cle{critère de fécondité}, critère biologique de l'espèce.
\end{enumerate}
\end{corrige}

\begin{corrige}[Exercice 3 — Définitions et vocabulaire]
\begin{enumerate}
    \item \textbf{Caractère héréditaire :} caractère inscrit dans l'information génétique (les gènes) des parents et transmis à la descendance lors de la reproduction. Exemples : groupe sanguin, couleur des yeux.
    \item \textbf{Caractère modifié par l'environnement (acquis) :} caractère apparu au cours de la vie de l'individu sous l'effet de son milieu de vie (alimentation, apprentissage, accidents, mode de vie). Il n'est pas inscrit dans les gènes et ne se transmet pas. Exemples : cicatrice, langue parlée.
    \item \textbf{Espèce biologique :} ensemble d'individus qui se ressemblent, peuvent se reproduire entre eux et donner des descendants \emph{féconds}.
    \item \textbf{Variabilité intra-spécifique :} ensemble des différences observées entre les individus d'une \emph{même} espèce, liées aux caractères héréditaires variables et aux caractères modifiés par l'environnement.
\end{enumerate}
\end{corrige}

\begin{corrige}[Exercice 4 — Calcul direct : Fréquence phénotypique]
Effectif total : $N = 40$ élèves.

\textbf{Lobes libres :}
\[ f_1 = \frac{28}{40}\times 100 = 70\ \% \]

\textbf{Lobes attachés :}
\[ f_2 = \frac{12}{40}\times 100 = 30\ \% \]

\textbf{Vérification :} $70\ \% + 30\ \% = 100\ \%$ et $28 + 12 = 40$ : le calcul est cohérent.

\textbf{Conclusion :} dans cet échantillon, le phénotype \og lobes libres \fg{} (allèle dominant) est majoritaire avec $70\ \%$, contre $30\ \%$ pour le phénotype \og lobes attachés \fg{} (allèle récessif).

\textbf{Point de vigilance :} ne pas confondre fréquence (en \%) et effectif (nombre d'élèves) ; toujours vérifier que la somme des fréquences vaut $100\ \%$.
\end{corrige}

\courssub{Je m'entraîne}

\begin{corrige}[Exercice 5 — Classement de caractères dans une famille]
\textbf{1. Classement :}

\begin{center}
\begin{tabularx}{\linewidth}{|l|X|}\hline
\textbf{Caractère} & \textbf{Nature} \\ \hline
Groupe sanguin (ABO) & Héréditaire \\ \hline
Taille & Héréditaire \emph{et} influencée par l'environnement (mixte) \\ \hline
Lobe de l'oreille & Héréditaire \\ \hline
Cicatrice au bras & Modifié par l'environnement (acquis) \\ \hline
Couleur des yeux & Héréditaire \\ \hline
Langue parlée & Modifié par l'environnement (acquis) \\ \hline
\end{tabularx}
\end{center}

\textbf{2. Justifications :}
\begin{itemize}
    \item \textbf{Groupe sanguin :} il est déterminé par les gènes transmis par les parents et ne change jamais au cours de la vie ; il est donc héréditaire. (Un enfant O peut naître de parents A et B si ceux-ci sont hétérozygotes $I^A/i$ et $I^B/i$.)
    \item \textbf{Taille :} elle dépend des gènes des parents (l'enfant a une taille proche de celle de ses parents) mais aussi de l'alimentation et de la santé pendant la croissance : c'est un caractère mixte.
    \item \textbf{Cicatrice :} elle résulte d'un accident (chute de vélo) survenu au cours de la vie ; elle n'est inscrite dans aucun gène : c'est un caractère modifié par l'environnement, non transmissible.
    \item \textbf{Langue parlée :} elle s'apprend dans le milieu de vie (famille, école) ; le père parle Bassa, la mère Anglais, l'enfant Français : preuve qu'elle n'est pas héréditaire mais acquise.
\end{itemize}
\end{corrige}

\begin{corrige}[Exercice 6 — Exploitation d'histogramme : Variabilité de la taille]
\textbf{1. Classe modale :} l'effectif le plus élevé est 10 élèves, correspondant à la classe centrée sur \SI{165}{cm} (classe 162,5–167,5 cm).

\textbf{2. Taille moyenne approximative :} on utilise les centres de classes :
\begin{align*}
\bar{x} &= \frac{147{,}5\times 1 + 152{,}5\times 2 + 157{,}5\times 5 + 162{,}5\times 8 + 167{,}5\times 10 + 172{,}5\times 8 + 177{,}5\times 4 + 182{,}5\times 2}{40}\\
&= \frac{147{,}5 + 305 + 787{,}5 + 1300 + 1675 + 1380 + 710 + 365}{40}\\
&= \frac{6670}{40} = 166{,}75\ \text{cm}
\end{align*}
Vérification des effectifs : $1+2+5+8+10+8+4+2 = 40$. La taille moyenne est d'environ \SI{166,8}{cm}.

\textbf{3. Valeurs extrêmes :} minimum $\approx$ \SI{145}{cm} ; maximum $\approx$ \SI{182,5}{cm} (classe 180–185 cm). L'étendue est d'environ $37$ à \SI{40}{cm}.

\textbf{4. Deux explications :}
\begin{itemize}
    \item \textbf{Origine héréditaire :} les élèves ont des parents différents, donc des patrimoines génétiques différents qui fixent des tailles potentielles différentes.
    \item \textbf{Origine environnementale :} l'alimentation (protéines, calcium, vitamines), la santé pendant la croissance et les conditions de vie varient d'une famille à l'autre et modifient la taille réellement atteinte.
\end{itemize}
\end{corrige}

\begin{corrige}[Exercice 7 — Étude de cas : Vrais jumeaux élevés séparément]
\textbf{1. Classement :}
\begin{itemize}
    \item \textbf{Héréditaires :} groupe sanguin (A Rh+), couleur des yeux (marron), empreintes digitales.
    \item \textbf{Modifiés par l'environnement :} masse corporelle, métier, langue courante.
\end{itemize}

\textbf{2. Même groupe sanguin et même couleur des yeux :} les vrais jumeaux (monozygotes) proviennent du \emph{même} ovule fécondé : ils possèdent donc un \cle{patrimoine génétique identique}. Or le groupe sanguin et la couleur des yeux sont des caractères purement héréditaires, déterminés uniquement par les gènes : ils sont donc obligatoirement identiques chez les deux jumeaux.

\textbf{3. Masses et métiers différents :} la masse corporelle dépend de l'alimentation et du mode de vie ; le métier et la langue dépendent du milieu social et culturel. Ayant grandi dans des environnements différents (ville de Douala / zone rurale du Nord), les jumeaux ont acquis des caractères différents malgré des gènes identiques.

\textbf{4. Conclusion (exemple de rédaction) :} Les différences entre individus de la même espèce ont une double origine. Une partie est héréditaire, inscrite dans les gènes transmis par les parents (groupe sanguin, couleur des yeux). L'autre partie est acquise sous l'effet de l'environnement (alimentation, apprentissage, mode de vie). L'étude des vrais jumeaux élevés séparément le prouve : à gènes identiques, seuls les caractères influencés par le milieu diffèrent.
\end{corrige}

\begin{corrige}[Exercice 8 — Critère de l'espèce et fécondité]
\textbf{1. L'élève a tort :} selon le \cle{critère de fécondité}, deux individus appartiennent à la même espèce s'ils peuvent se reproduire entre eux et donner des descendants \emph{féconds}. Or un Baka et une Peule (ou inversement) peuvent avoir ensemble des enfants eux-mêmes capables de se reproduire. Ils appartiennent donc bien à la \emph{même espèce} : \emph{Homo sapiens}. Les différences physiques et culturelles ne remettent pas en cause cette unité biologique.

\textbf{2. Trois caractères de l'espèce humaine (communs à tous) :}
\begin{itemize}
    \item Un squelette interne avec une colonne vertébrale ;
    \item La station debout et la marche bipède ;
    \item Un cerveau développé logé dans une boîte crânienne (et un langage articulé) ;
    \item Deux bras et deux jambes, des mains à pouce opposable.
\end{itemize}

\textbf{3. Exemple de caractère variable :} la taille moyenne ou la couleur de la peau varient entre ces populations : ce sont des caractères principalement \textbf{héréditaires} (avec une influence de l'environnement pour la taille). La langue parlée (langue baka / fulfulde) est, elle, un caractère \textbf{modifié par l'environnement}.

\textbf{Point de vigilance :} ne jamais confondre \og race \fg{} ou ethnie avec \og espèce \fg : il n'existe qu'\emph{une seule} espèce humaine.
\end{corrige}

\courssub{Je me prépare au BEPC}

\begin{corrige}[Exercice 9 — Sujet type BEPC : Hérédité du groupe sanguin et diversité]
\textbf{1. Génotypes des parents :} l'enfant est de groupe O, donc de génotype $i/i$ : il a reçu un allèle $i$ de \emph{chaque} parent. Le père est de groupe A : il possède au moins un allèle $I^A$ ; puisqu'il a transmis $i$, son génotype est $I^A/i$. De même, la mère de groupe B est de génotype $I^B/i$.

Échiquier de croisement $I^A/i \times I^B/i$ :
\begin{center}
\begin{tabular}{|c|c|c|}\hline
 & $I^B$ & $i$ \\ \hline
$I^A$ & $I^A/I^B$ (AB) & $I^A/i$ (A) \\ \hline
$i$ & $I^B/i$ (B) & $i/i$ (O) \\ \hline
\end{tabular}
\end{center}
Chaque enfant a donc 1 chance sur 4 d'être de groupe O : le résultat observé est possible.

\textbf{2. Caractère héréditaire :} le groupe sanguin est déterminé par les allèles ($I^A$, $I^B$, $i$) transmis par les parents lors de la reproduction ; il est fixé dès la naissance et ne change jamais au cours de la vie, quel que soit l'environnement. C'est donc un caractère héréditaire.

\textbf{3. Autres exemples :}
\begin{itemize}
    \item \textbf{Héréditaires :} la couleur des yeux ; la forme du lobe de l'oreille (libre/attaché) ; la couleur de la peau.
    \item \textbf{Modifiés par l'environnement :} une cicatrice ; la langue parlée ; la pratique d'un sport ou d'un métier.
\end{itemize}

\textbf{4. Effet d'une alimentation carencée :}
\begin{itemize}
    \item \textbf{Modifiés par l'environnement :} la \emph{taille adulte} (retard de croissance par manque de protéines et vitamines) et le \emph{développement cognitif} (les carences peuvent perturber le développement et les apprentissages).
    \item \textbf{Inchangés :} le \emph{groupe sanguin} et la \emph{couleur de peau} (caractères purement héréditaires) ; la \emph{langue parlée} dépend du milieu linguistique, pas de l'alimentation.
\end{itemize}

\textbf{Barème indicatif (sur 10) :} génotypes justifiés 3 pts ; caractère héréditaire expliqué 2 pts ; exemples 2 pts ; classement final justifié 3 pts.

\textbf{Points de vigilance :} justifier les génotypes par la transmission obligatoire d'un allèle $i$ par chaque parent ; ne pas écrire que le groupe O est \og dominant \fg{} (l'allèle $i$ est récessif) ; distinguer clairement héréditaire et modifié par l'environnement dans la dernière question.
\end{corrige}

\begin{corrige}[Exercice 10 — Sujet type BEPC : Unité biologique et lutte contre les discriminations]
\textbf{Exemple de texte argumenté :}

\og Monsieur le Chef, chers parents, la rumeur qui exclut les enfants albinos est contraire à la science. Tous les êtres humains appartiennent à une seule espèce, \emph{Homo sapiens}, car un homme et une femme de n'importe quelle communauté peuvent avoir ensemble des enfants féconds : c'est le critère biologique de l'espèce. Les personnes albinos possèdent tous les caractères de l'espèce humaine : squelette interne, station debout, cerveau développé, langage. L'albinisme n'est qu'un caractère héréditaire récessif : un gène responsable de la fabrication de la mélanine, le pigment qui colore la peau, est modifié. Comme la couleur des yeux ou le groupe sanguin, c'est une simple variante de la diversité humaine ; il n'affecte ni l'intelligence, ni la santé globale, ni la capacité à avoir des enfants. Rejeter un enfant albinos, c'est donc rejeter un être humain à part entière, notre semblable. Accueillons tous les enfants à l'école : l'unité de l'espèce humaine est un fait scientifique, et sa diversité est une richesse. \fg

\textbf{Barème indicatif (sur 10) :} critère de fécondité correctement utilisé 2,5 pts ; distinction caractères de l'espèce / caractères variables 2,5 pts ; origine génétique récessive de l'albinisme et absence d'effet sur la fécondité 3 pts ; conclusion sur l'unité biologique et qualité argumentative 2 pts.

\textbf{Points de vigilance :} employer le vocabulaire scientifique exact (espèce, caractère héréditaire récessif, fécondité) ; ne jamais écrire qu'il existe plusieurs espèces humaines ; le texte doit rester respectueux et argumenté, sans simple opinion.
\end{corrige}

\begin{corrige}[Exercice 11 — Sujet type BEPC : Synthèse — Hérédité, Environnement, Unité]
\textbf{1. Analyse des résultats :}
\begin{enumerate}
    \item \textbf{Groupe sanguin :} concordance de $100\ \%$ chez les vrais jumeaux (gènes identiques) contre $50\ \%$ chez les faux jumeaux (gènes différents, comme des frères et sœurs ordinaires). Conclusion : le groupe sanguin est un caractère \textbf{entièrement héréditaire}, déterminé uniquement par les gènes.
    \item \textbf{Myopie :} la concordance est élevée chez les vrais jumeaux ($85\ \%$) mais pas totale, alors que leurs gènes sont identiques. Cela suggère que la myopie dépend des gènes \emph{et} de facteurs environnementaux (lecture prolongée, écrans, éclairage...) : c'est un caractère \textbf{mixte} (héréditaire et modifié par l'environnement).
    \item \textbf{Langue maternelle :} concordance de $100\ \%$ dans les deux cas car tous les jumeaux sont élevés dans les \emph{mêmes familles}, donc le même milieu linguistique. La langue est un caractère \textbf{modifié par l'environnement} (appris), totalement indépendant des gènes.
\end{enumerate}

\textbf{2. Modélisation (tableau complété) :}
\begin{center}
\begin{tabular}{|l|c|c|}\hline
\textbf{Caractère} & \textbf{Principalement Héréditaire} & \textbf{Principalement Modifié par l'environnement} \\ \hline
Groupe sanguin & X & \\ \hline
Myopie & X (avec influence de l'environnement) & \\ \hline
Langue maternelle & & X \\ \hline
\end{tabular}
\end{center}

\textbf{3. Argumentation :} l'affirmation du journaliste est trop simpliste. Si la myopie était \emph{purement} génétique, deux vrais jumeaux, qui ont exactement les mêmes gènes, seraient \emph{toujours} concordants ($100\ \%$). Or le tableau montre une concordance de seulement $85\ \%$ : dans $15\ \%$ des paires, un seul jumeau est myope. Des facteurs d'environnement (travail de près, écrans, lumière) interviennent donc aussi. La myopie a une composante héréditaire forte mais n'est pas \og uniquement \fg{} génétique.

\textbf{4. Bilan :} un \textbf{caractère de l'espèce humaine} est un caractère commun à tous les individus sains de l'espèce, qui traduit son unité biologique. Exemples (hors tableau) : le squelette interne avec colonne vertébrale, la station debout bipède, le cerveau développé, deux bras et deux jambes.

\textbf{Barème indicatif (sur 10) :} analyse des trois caractères 4 pts ; tableau correctement rempli 2 pts ; critique argumentée à partir des données 2,5 pts ; définition + exemple 1,5 pt.

\textbf{Points de vigilance :} fonder chaque conclusion \emph{uniquement} sur les données du tableau (comparer $100\ \%$ / $50\ \%$ / $85\ \%$) ; rappeler que vrais jumeaux = gènes identiques, faux jumeaux = gènes différents ; pour la myopie, citer explicitement l'écart à $100\ \%$ comme preuve du rôle de l'environnement.
\end{corrige}

\newpage

\coursec{Fiche bilan}

\courssub{Carte mentale de la leçon}

\begin{popfigure}
\begin{tikzpicture}[mindmap, grow cyclic, every node/.style={concept, text=white, font=\small},
  root concept/.append style={concept color=popblue, font=\small\bfseries},
  level 1/.append style={level distance=3.4cm, sibling angle=60},
  level 2/.append style={level distance=2.1cm, sibling angle=45, font=\scriptsize}]
\node[root concept]{Ressemblances et différences dans l'espèce humaine}
  child[concept color=popgreen]{ node {Notion de \cle{caractère}}
    child { node {caractère observable ou mesurable} }
  }
  child[concept color=poporange]{ node {Ressemblances}
    child { node {caractères de l'espèce : bipédie, station debout, langage, cerveau développé} }
  }
  child[concept color=poppurple]{ node {Différences entre individus}
    child { node {groupe sanguin, couleur des yeux, taille, empreintes digitales} }
  }
  child[concept color=popink]{ node {Caractères \cle{héréditaires}}
    child { node {transmis par les parents, présents à la naissance} }
  }
  child[concept color=popgold]{ node {Caractères modifiés par l'\cle{environnement}}
    child { node {taille, masse, teint de peau, cicatrices} }
  }
  child[concept color=poppink]{ node {Applications}
    child { node {identification, transfusion, respect des différences} }
  };
\end{tikzpicture}
\legende{Carte mentale de la leçon : l'espèce humaine est une (caractères communs) mais chaque individu est unique (caractères héréditaires et caractères modifiés par l'environnement).}
\end{popfigure}

\courssub{L'essentiel à retenir}

\begin{bilan}
\textbf{Définitions clés}
\begin{itemize}
  \item \cle{Espèce} : ensemble d'individus qui se ressemblent, peuvent se reproduire entre eux et donnent des descendants fertiles.
  \item \cle{Caractère} : tout aspect observable ou mesurable d'un individu (forme, couleur, taille, groupe sanguin...).
  \item \cle{Caractère héréditaire} : caractère transmis des parents aux enfants ; il est présent dès la naissance (groupe sanguin, couleur des yeux, forme du lobe de l'oreille).
  \item \cle{Caractère modifié par l'environnement} : caractère influencé par les conditions de vie (alimentation, soleil, accidents, exercice physique) : taille, masse, teint de la peau, cicatrices.
\end{itemize}

\textbf{Les caractères de l'espèce humaine (communs à tous les hommes)}
\begin{itemize}
  \item Station debout et marche bipède ; membres supérieurs libres avec pouce opposable.
  \item Cerveau volumineux et développé ; langage articulé ; pensée et fabrication d'outils.
  \item Pilosité réduite ; reproduction sexuée ; longue période de croissance.
\end{itemize}

\textbf{Tableau récapitulatif à connaître par c\oe ur}
\begin{center}
\begin{tabularx}{\linewidth}{|l|X|X|}
\hline
\rowcolor{popblueL} \textbf{Critère} & \textbf{Caractère héréditaire} & \textbf{Modifié par l'environnement} \\
\hline
Origine & transmis par les parents & conditions de vie \\
\hline
Présent à la naissance & oui & non (apparaît ou varie après) \\
\hline
Exemples & groupe sanguin, couleur des yeux, albinisme & taille, masse, teint de peau, cicatrices \\
\hline
Peut varier au cours de la vie & non & oui \\
\hline
\end{tabularx}
\end{center}

\textbf{Méthodes express}
\begin{itemize}
  \item Pour classer un caractère : se poser deux questions : \og est-il présent à la naissance ? \fg{} et \og peut-il changer selon le mode de vie ? \fg{}
  \item Pour comparer des individus : dresser l'inventaire des caractères communs (ressemblances) puis des caractères variables (différences).
  \item Pour rédiger une conclusion : citer les observations, dégager le caractère étudié, conclure sur son origine (héréditaire ou environnementale).
\end{itemize}
\end{bilan}

\begin{aretenir}[Checklist avant le BEPC]
\begin{itemize}
  \item[$\square$] Je sais définir une espèce et un caractère.
  \item[$\square$] Je cite au moins 4 caractères communs à tous les êtres humains.
  \item[$\square$] J'explique pourquoi tous les hommes appartiennent à la même espèce.
  \item[$\square$] Je cite au moins 4 caractères qui diffèrent d'un individu à l'autre.
  \item[$\square$] Je sais définir un caractère héréditaire et donner 3 exemples.
  \item[$\square$] Je sais définir un caractère modifié par l'environnement et donner 3 exemples.
  \item[$\square$] Je sais classer un caractère donné (héréditaire ou environnemental) en le justifiant.
  \item[$\square$] Je sais que certains caractères (taille, masse) dépendent à la fois de l'hérédité et de l'environnement.
  \item[$\square$] Je sais que les empreintes digitales sont uniques et servent à identifier les individus.
  \item[$\square$] Je rédige une conclusion en m'appuyant sur mes observations et en respectant le vocabulaire scientifique.
\end{itemize}
\end{aretenir}

\findecours

\end{document}
